% Limitation des dysfonctionnements des organes ou structures intervenant dans les mouvements réflexes — SVT Tle D — Cours signé OnBuch+
% Programme officiel MINESEC. Compiler avec : tectonic N07-reflexes-medullaires-myotatiques.tex
\newcommand{\DOCMATIERE}{SVT}
\newcommand{\DOCNIVEAU}{Tle D}
\newcommand{\DOCMODULE}{Module I : Le Monde Vivant — Activités réflexes}
\newcommand{\DOCLECON}{N07}
\newcommand{\DOCTITRE}{Limitation des dysfonctionnements des organes ou structures intervenant dans les mouvements réflexes}
% OnBuch+ — préambule « cours signés » ENRICHI (compilé avec tectonic / XeLaTeX)
% Système visuel « Pop » d'OnBuch+ : fond crème (jamais blanc), bordures encre
% épaisses, ombres « dures » décalées, titres Archivo Black en capitales.
% Version MATHÉMATIQUES : graphes (pgfplots/tikz), boîtes pédagogiques
% (définition, propriété, méthode, attention, le savais-tu, exercice résolu,
% exercice, TP, bilan), page de garde et signature OnBuch+ ancrée en bas.
%
% Macros attendues dans main.tex AVANT \input{preamble} :
%   \DOCMATIERE  \DOCNIVEAU  \DOCMODULE  \DOCLECON (numéro)  \DOCTITRE
\documentclass[11pt]{article}

\usepackage{fontspec}
\usepackage[french]{babel}
\usepackage[a4paper,margin=2cm,top=3.4cm,bottom=2.8cm,headheight=1.4cm,footskip=1.3cm]{geometry}
\usepackage{amsmath,amssymb,amsfonts}
\usepackage{mathtools} % \xmapsto, \coloneqq... souvent utilisés par les modèles
\usepackage{cancel} % simplifications barrées (bilans de forces)
% \abs{z} (module, valeur absolue) et \norm{u} (norme), utilisés par les modèles
\providecommand{\abs}{}\let\abs\relax\DeclarePairedDelimiter{\abs}{\lvert}{\rvert}
\providecommand{\norm}{}\let\norm\relax\DeclarePairedDelimiter{\norm}{\lVert}{\rVert}
\usepackage[table]{xcolor} % [table] : fournit aussi \rowcolors (alternance de lignes)
\usepackage{pagecolor}
\usepackage{tikz}
\usetikzlibrary{arrows.meta,positioning,calc,shapes.geometric,shapes.misc,shapes.arrows,decorations.pathmorphing,decorations.pathreplacing,decorations.markings,patterns,shadows,mindmap,trees,fit,backgrounds,matrix,intersections,angles,quotes}
\usepackage{pgfplots}
\usepackage[version=4]{mhchem} % équations nucléaires \ce{^{235}_{92}U}
\usepackage[european,siunitx]{circuitikz} % schémas électriques
\pgfplotsset{compat=1.18}
\usepgfplotslibrary{fillbetween,groupplots} % aires entre courbes (calcul intégral)
\usepackage{enumitem}
\usepackage{fancyhdr}
% [most] charge toutes les bibliothèques tcolorbox (listings, minted, raster,
% documentation...) et gonfle inutilement la mémoire TeX sur un document long
% (~300 boîtes) : on ne charge que ce qui est réellement utilisé.
\usepackage[skins,breakable]{tcolorbox}
\usepackage{graphicx}
\usepackage{colortbl}
\usepackage{booktabs}
\usepackage{tabularx}
\usepackage{diagbox} % case d'en-tête diagonale des tableaux à double entrée
% Colonnes extensibles centrée / gauche / droite (C, L, R), souvent utilisées par les modèles
\newcolumntype{C}{>{\centering\arraybackslash}X}
\newcolumntype{L}{>{\raggedright\arraybackslash}X}
\newcolumntype{R}{>{\raggedleft\arraybackslash}X}
\usepackage{array}
\usepackage{multicol}
\usepackage{siunitx}
% parse-numbers=false : accepte aussi \SI{2,5 \times 10^{-3}}{...}
\sisetup{locale=FR,output-decimal-marker={,},group-minimum-digits=4,parse-numbers=false}
\DeclareSIUnit\ohm{\ensuremath{\symup{\Omega}}} % U+2126 absent de la police
\DeclareSIUnit\division{div} % réglages d'oscilloscope (ms/div, V/div)
\DeclareSIUnit\tr{tr} % tours (vitesse de rotation, ex. \SI{1500}{\tr\per\minute})
\DeclareSIUnit\molar{M} % concentration molaire (ex. \SI{3}{\molar}, \SI{1}{\micro\molar})
\DeclareSIUnit\an{an} % année (le modèle confond parfois avec \year, un compteur TeX) — ex. \SI{5}{\kilo\meter\per\an}
\DeclareSIUnit\FCFA{FCFA} % franc CFA (ex. \SI{500}{\FCFA\per\kilo\gram})
\DeclareSIUnit\degreeBrix{\textdegree Brix} % teneur en sucre d'un sirop/jus (réfractométrie)
\DeclareSIUnit\month{mois} % le modèle utilise parfois \month, un compteur TeX, comme unité de durée
\usepackage{qrcode}
% \degree : commande fréquemment utilisée par les modèles pour un angle ou une
% température en dehors de \SI{}{} (non fournie par les paquets chargés).
\providecommand{\degree}{^{\circ}}
% \qed (fin de démonstration), utilisé par les modèles sans amsthm
\providecommand{\qed}{\hfill$\square$}
% \barre : double barre des valeurs interdites dans un tableau de variations (tkz-tab)
\providecommand{\barre}{\ensuremath{\|}}
% environnement proof (amsthm non chargé) utilisé par les modèles
\ifdefined\proof\else\newenvironment{proof}[1][Démonstration]{\par\noindent\textit{#1.}\ }{\qed\par}\fi
\usepackage{lastpage}
\usepackage{hyperref}
\hypersetup{hidelinks}

% ============================================================
% Palette « Pop » OnBuch+ — mêmes hex que la classe P (plus_ui.dart)
% ============================================================
\definecolor{poporange}{HTML}{F59321}
\definecolor{popdark}{HTML}{E07A0C}
\definecolor{popcream}{HTML}{FFF6E8}
\definecolor{popink}{HTML}{1C1714}
\definecolor{poppurple}{HTML}{7A5AE0}
\definecolor{popgreen}{HTML}{2E9E6A}
\definecolor{poppink}{HTML}{E0457B}
\definecolor{popblue}{HTML}{2F7DE1}
\definecolor{popgold}{HTML}{C98A12}
\definecolor{popmuted}{HTML}{8A7F76}
\definecolor{popline}{HTML}{EADFD2}
\definecolor{popgray}{HTML}{8A7F76}
\colorlet{popgrey}{popgray}
% Alias tolérants (le modèle nomme parfois les couleurs différemment)
\colorlet{popwhite}{white}
\colorlet{popblack}{popink}
\colorlet{popred}{poppink}
\colorlet{popyellow}{popgold}
% Teintes claires (fonds de tableaux/encadrés) — le modèle nomme parfois
% les couleurs avec un suffixe L (ex. popblueL) sans les avoir définies.
\colorlet{poporangeL}{poporange!20}
\colorlet{popdarkL}{popdark!20}
\colorlet{poppurpleL}{poppurple!20}
\colorlet{popgreenL}{popgreen!20}
\colorlet{poppinkL}{poppink!20}
\colorlet{popblueL}{popblue!20}
\colorlet{popgoldL}{popgold!20}
\colorlet{popredL}{popred!20}
\colorlet{popgrayL}{popgray!20}
% Teintes claires pour fonds de tableaux / zones
\colorlet{popblueL}{popblue!10!white}
\colorlet{poporangeL}{poporange!14!white}
\colorlet{popgreenL}{popgreen!12!white}
\colorlet{poppurpleL}{poppurple!12!white}
\colorlet{poppinkL}{poppink!12!white}
\colorlet{popgoldL}{popgold!14!white}

\pagecolor{popcream}

% ============================================================
% Typographie
% ============================================================
\defaultfontfeatures{Ligatures=TeX}
\setmainfont{PlusJakartaSans-Medium.ttf}[Path=../../../fonts/,BoldFont=PlusJakartaSans-Bold.ttf]
% Maths en sans-serif (Fira Math) avec chiffres et lettres droites de Plus
% Jakarta, pour que les formules se fondent dans le texte.
\usepackage[math-style=ISO,bold-style=ISO]{unicode-math}
\setmathfont{FiraMath-Regular.otf}[Path=../../../fonts/]
\setmathfont{PlusJakartaSans-Medium.ttf}[Path=../../../fonts/,range={up/num,up/{Latin,latin},\mathup}]
\usepackage{icomma} % virgule décimale sans espace en mode math
% Caractères mathématiques Unicode absents des polices de texte (Plus Jakarta,
% Archivo Black) : rendus par leur équivalent mathématique, en texte comme en
% formule — sinon ils s'affichent en carré vide (ex. « Arithmétique dans ℤ »).
\usepackage{newunicodechar}
\newunicodechar{ℝ}{\ensuremath{\mathbb{R}}}
\newunicodechar{ℤ}{\ensuremath{\mathbb{Z}}}
\newunicodechar{ℂ}{\ensuremath{\mathbb{C}}}
\newunicodechar{ℕ}{\ensuremath{\mathbb{N}}}
\newunicodechar{ℚ}{\ensuremath{\mathbb{Q}}}
\newunicodechar{𝔻}{\ensuremath{\mathbb{D}}}
\newunicodechar{α}{\ensuremath{\alpha}}
\newunicodechar{β}{\ensuremath{\beta}}
\newunicodechar{γ}{\ensuremath{\gamma}}
\newunicodechar{δ}{\ensuremath{\delta}}
\newunicodechar{ε}{\ensuremath{\varepsilon}}
\newunicodechar{θ}{\ensuremath{\theta}}
\newunicodechar{λ}{\ensuremath{\lambda}}
\newunicodechar{μ}{\ensuremath{\mu}}
\newunicodechar{σ}{\ensuremath{\sigma}}
\newunicodechar{τ}{\ensuremath{\tau}}
\newunicodechar{φ}{\ensuremath{\varphi}}
\newunicodechar{ψ}{\ensuremath{\psi}}
\newunicodechar{ω}{\ensuremath{\omega}}
\newunicodechar{Σ}{\ensuremath{\Sigma}}
% majuscules produites par \MakeUppercase dans les titres (α-aminés -> Α)
\newunicodechar{Α}{\ensuremath{\alpha}}
\newunicodechar{Β}{\ensuremath{\beta}}
\newunicodechar{Γ}{\ensuremath{\Gamma}}
\newunicodechar{Δ}{\ensuremath{\Delta}}
\newunicodechar{Ε}{\ensuremath{\varepsilon}}
\newunicodechar{Θ}{\ensuremath{\Theta}}
\newunicodechar{Λ}{\ensuremath{\Lambda}}
\newunicodechar{Μ}{\ensuremath{\mu}}
\newunicodechar{Τ}{\ensuremath{\tau}}
\newunicodechar{Φ}{\ensuremath{\Phi}}
\newunicodechar{Ψ}{\ensuremath{\Psi}}
\newunicodechar{Ω}{\ensuremath{\Omega}}
\newunicodechar{↦}{\ensuremath{\mapsto}}
\newunicodechar{∈}{\ensuremath{\in}}
\newunicodechar{∉}{\ensuremath{\notin}}
\newunicodechar{∧}{\ensuremath{\wedge}}
\newunicodechar{⇔}{\ensuremath{\Leftrightarrow}}
\newunicodechar{⟺}{\ensuremath{\Longleftrightarrow}}
\newunicodechar{⇒}{\ensuremath{\Rightarrow}}
\newunicodechar{⟹}{\ensuremath{\Longrightarrow}}
\newunicodechar{∘}{\ensuremath{\circ}}
\newunicodechar{∪}{\ensuremath{\cup}}
\newunicodechar{∩}{\ensuremath{\cap}}
\newunicodechar{≡}{\ensuremath{\equiv}}
\newunicodechar{⊂}{\ensuremath{\subset}}
\newunicodechar{⊥}{\ensuremath{\perp}}
\newunicodechar{∃}{\ensuremath{\exists}}
\newunicodechar{∀}{\ensuremath{\forall}}
\newunicodechar{∛}{\ensuremath{\sqrt[3]{\,}}}
\newunicodechar{∜}{\ensuremath{\sqrt[4]{\,}}}
\newunicodechar{⃗}{\ensuremath{{}^{\to}}}
\newunicodechar{ⁿ}{\ifmmode^{n}\else\textsuperscript{\ensuremath{n}}\fi}
\newunicodechar{ˣ}{\ifmmode^{x}\else\textsuperscript{\ensuremath{x}}\fi}
\newunicodechar{ᵇ}{\ifmmode^{b}\else\textsuperscript{\ensuremath{b}}\fi}
\newunicodechar{ʳ}{\ifmmode^{r}\else\textsuperscript{\ensuremath{r}}\fi}
\newunicodechar{ᵘ}{\ifmmode^{u}\else\textsuperscript{\ensuremath{u}}\fi}
\newunicodechar{ʸ}{\ifmmode^{y}\else\textsuperscript{\ensuremath{y}}\fi}
\newunicodechar{ᵃ}{\ifmmode^{a}\else\textsuperscript{\ensuremath{a}}\fi}
\newunicodechar{ᵉ}{\ifmmode^{e}\else\textsuperscript{\ensuremath{e}}\fi}
\newunicodechar{ᵏ}{\ifmmode^{k}\else\textsuperscript{\ensuremath{k}}\fi}
\newunicodechar{ᵖ}{\ifmmode^{p}\else\textsuperscript{\ensuremath{p}}\fi}
\newunicodechar{ᴺ}{\ifmmode^{N}\else\textsuperscript{\ensuremath{N}}\fi}
\newunicodechar{ᶜ}{\ifmmode^{c}\else\textsuperscript{\ensuremath{c}}\fi}
\newunicodechar{ʰ}{\ifmmode^{h}\else\textsuperscript{\ensuremath{h}}\fi}
\newunicodechar{ᵠ}{\ifmmode^{q}\else\textsuperscript{\ensuremath{q}}\fi}
\newunicodechar{ₙ}{\ifmmode_{n}\else\textsubscript{\ensuremath{n}}\fi}
\newunicodechar{ₐ}{\ifmmode_{a}\else\textsubscript{\ensuremath{a}}\fi}
\newunicodechar{ᵢ}{\ifmmode_{i}\else\textsubscript{\ensuremath{i}}\fi}
\newunicodechar{ᵣ}{\ifmmode_{r}\else\textsubscript{\ensuremath{r}}\fi}
\newunicodechar{ₖ}{\ifmmode_{k}\else\textsubscript{\ensuremath{k}}\fi}
\newunicodechar{ᵦ}{\ifmmode_{\beta}\else\textsubscript{\ensuremath{\beta}}\fi}
\newunicodechar{ₓ}{\ifmmode_{x}\else\textsubscript{\ensuremath{x}}\fi}
\newfontfamily\popdisplay{ArchivoBlack-Regular.ttf}[Path=../../../fonts/]
\newfontfamily\poplogo{SpaceGrotesk-Bold.ttf}[Path=../../../fonts/]
\newfontfamily\popsemi{PlusJakartaSans-SemiBold.ttf}[Path=../../../fonts/]
\newfontfamily\popxbold{PlusJakartaSans-ExtraBold.ttf}[Path=../../../fonts/]
\newfontfamily\popxboldls{PlusJakartaSans-ExtraBold.ttf}[Path=../../../fonts/,LetterSpace=3.0]

\setlist[enumerate]{leftmargin=1.7em,itemsep=5pt,topsep=4pt}
\setlist[itemize]{leftmargin=1.7em,itemsep=4pt,topsep=4pt,label={\color{poporange}\textbullet}}
\setlength{\parindent}{0pt}
\setlength{\parskip}{5pt}
\renewcommand{\arraystretch}{1.25}

% pgfplots : style Pop par défaut
\pgfplotsset{
  every axis/.append style={axis line style={popink,line width=1pt},tick style={popink},
    grid style={popline,line width=0.5pt},label style={font=\small},tick label style={font=\footnotesize}},
  popcurve/.style={poporange,line width=1.8pt,no markers,smooth},
}
% Même style disponible dans un \draw TikZ simple (hors pgfplots)
\tikzset{popcurve/.style={poporange,line width=1.8pt}}
\tikzset{popgrid/.style={popline,very thin}}
\colorlet{popcurve}{poporange} % le nom du style est parfois utilisé comme couleur

% ============================================================
% Pill / squiggle
% ============================================================
\newcommand{\poppill}[3]{%
  \tikz[baseline=-0.55ex]\node[draw=popink,line width=1.2pt,rounded corners=2.6pt,%
    fill=#2,inner xsep=6.5pt,inner ysep=3pt]{%
    \popxboldls\fontsize{7}{7}\selectfont\color{#3}\MakeUppercase{#1}};%
}
\newcommand{\popsquiggle}[1]{%
  \tikz[baseline=-0.4ex]{
    \draw[#1,line width=2.6pt,line cap=round]
      plot [smooth] coordinates {(0,0.09) (0.34,0.01) (0.68,0.21) (1.02,0.01) (1.36,0.10)};
  }%
}
% Mot-clé surligné (marqueur orange sous le texte)
\newcommand{\cle}[1]{{\popxbold\color{popdark}#1}}

% ============================================================
% En-tête / pied
% ============================================================
\pagestyle{fancy}
\fancyhf{}
\renewcommand{\headrulewidth}{0pt}
\renewcommand{\footrulewidth}{0pt}
\lhead{\raisebox{-0.15ex}{\poplogo\fontsize{13}{13}\selectfont\color{popink}OnBuch\color{poporange}+}}
\rhead{\poppill{\DOCMATIERE\ \textbullet\ \DOCNIVEAU\ \textbullet\ Leçon \DOCLECON}{white}{popink}}
\lfoot{\popxbold\fontsize{7.6}{7.6}\selectfont\color{popmuted}\MakeUppercase{Cours signé OnBuch+}\ \textbullet\ {\popsemi\DOCTITRE}}
\rfoot{\tikz[baseline=-0.6ex]\node[rounded corners=8.5pt,fill=popink,minimum height=17pt,minimum width=17pt,inner xsep=5pt,inner sep=0]{\popxbold\fontsize{8}{8}\selectfont\color{popcream}\thepage\,/\,\pageref*{LastPage}};}

% ============================================================
% Titres
% ============================================================
\newcommand{\chaptitre}[2]{%
  \begin{tcolorbox}[enhanced,colback=poporange,colframe=popink,boxrule=2.6pt,arc=11pt,
      drop shadow=popink,left=15pt,right=15pt,top=12pt,bottom=11pt,before skip=3pt,after skip=18pt]
    \poppill{#2}{popink}{popcream}\par\vspace{9pt}
    {\raggedright\hyphenpenalty=10000\exhyphenpenalty=10000\popdisplay\fontsize{22}{24}\selectfont\color{popcream}\MakeUppercase{#1}\par}
  \end{tcolorbox}
}
% Section numérotée : pastille orange avec le numéro + titre Archivo
\newcounter{popsec}
\newcommand{\coursec}[1]{%
  \stepcounter{popsec}\setcounter{popsub}{0}%
  \par\vspace{16pt}\noindent
  \begin{minipage}{\linewidth}
  \tikz[baseline=-0.7ex]\node[draw=popink,line width=1.6pt,fill=poporange,rounded corners=4pt,minimum size=22pt,inner sep=0]{\popdisplay\fontsize{12}{12}\selectfont\color{popcream}\thepopsec};%
  \hspace{8pt}{\popdisplay\fontsize{15}{17}\selectfont\color{popink}\MakeUppercase{#1}}\par
  \vspace{4pt}\noindent{\color{poporange}\rule{36pt}{3.6pt}}
  \end{minipage}\par\nopagebreak\vspace{8pt}
}
\newcounter{popsub}
\newcommand{\courssub}[1]{%
  \stepcounter{popsub}%
  \par\vspace{9pt}\noindent{\popxbold\fontsize{11.5}{13}\selectfont\color{popink}{\color{poporange}\thepopsec.\thepopsub}\hspace{6pt}#1}\par\nopagebreak\vspace{4pt}
}

% ============================================================
% Boîtes pédagogiques — même recette : fond blanc, bordure épaisse colorée,
% ombre dure encre, badge plein. Titre optionnel [#1].
% ============================================================
\tcbset{popbase/.style={breakable,enhanced,colback=white,boxrule=2.3pt,arc=9pt,
  shadow={3pt}{-3pt}{0pt}{popink},left=14pt,right=14pt,top=11pt,bottom=11pt,before skip=11pt,after skip=15pt,
  fontupper=\popsemi\fontsize{10.6}{14.5}\selectfont\color{popink},
  fontlower=\popsemi\fontsize{10.6}{14.5}\selectfont\color{popink}}}
% Coche dessinée (le glyphe ✓ n'existe pas dans les polices du document)
\newcommand{\popcheck}[1]{\tikz[baseline=-0.5ex]\draw[#1,line width=1.6pt,line cap=round,line join=round](-0.13,0.0)--(-0.04,-0.09)--(0.13,0.1);}
\providecommand{\coche}{\popcheck{popgreen}}
\providecommand{\resultat}[1]{\par\smallskip\noindent\fcolorbox{popgreen}{popgreenL}{\parbox{\dimexpr\linewidth-2\fboxsep-2\fboxrule\relax}{\textbf{Résultat.} #1}}\par\smallskip}
\newcommand{\popboxhead}[3]{\poppill{#1}{#2}{white}\ifx\relax#3\relax\else\hspace{6pt}{\popxbold\fontsize{10.6}{12}\selectfont\color{popink}#3}\fi\par\vspace{7pt}}

\newtcolorbox{aretenir}[1][]{popbase,colframe=popgreen,before upper={\popboxhead{À retenir}{popgreen}{#1}}}
\newtcolorbox{exemplebox}[1][]{popbase,colframe=poporange,before upper={\popboxhead{Exemple}{poporange}{#1}}}
\newtcolorbox{definition}[1][]{popbase,colframe=popblue,before upper={\popboxhead{Définition}{popblue}{#1}}}
\newtcolorbox{propriete}[1][]{popbase,colframe=poppurple,before upper={\popboxhead{Propriété}{poppurple}{#1}}}
\newtcolorbox{methode}[1][]{popbase,colframe=popgold,before upper={\popboxhead{Méthode}{popgold}{#1}}}
\newtcolorbox{attention}[1][]{popbase,colframe=poppink,before upper={\popboxhead{Attention}{poppink}{#1}}}
\newtcolorbox{experience}[1][]{popbase,colframe=popblue,colback=popblueL,before upper={\popboxhead{Expérience}{popblue}{#1}}}
% « Le savais-tu ? » : carte violette pleine (contexte, culture, Cameroun)
\newtcolorbox{savaistu}[1][]{popbase,colback=poppurple,colframe=popink,
  fontupper=\popsemi\fontsize{10.4}{14.2}\selectfont\color{white},
  before upper={\poppill{Le savais-tu\,?}{white}{poppurple}\ifx\relax#1\relax\else\hspace{6pt}{\popxbold\color{white}#1}\fi\par\vspace{7pt}}}
% Exercice résolu : énoncé en haut, \tcblower puis la solution sur fond crème
\newcounter{popexo}\newcounter{popexr}
\newtcolorbox{exoresolu}[1][]{popbase,colframe=poporange,colbacklower=popcream,
  segmentation style={popink,line width=1.2pt,dashed},
  before upper={\stepcounter{popexr}\popboxhead{Exercice résolu \thepopexr}{poporange}{#1}},
  before lower={\poppill{Solution}{popink}{popcream}\par\vspace{6pt}}}
% Exercice (non résolu en place) — la correction est dans la partie Corrigés
\newtcolorbox{exercice}[1][]{popbase,colframe=popink,
  before upper={\stepcounter{popexo}\popboxhead{Exercice \thepopexo}{popink}{#1}}}
% Corrigé
\newtcolorbox{corrige}[1][]{popbase,colframe=popgreen,colback=popgreenL,
  before upper={\popboxhead{Corrigé}{popgreen}{#1}}}
% Objectifs / prérequis / bilan
\newtcolorbox{objectifs}{popbase,colframe=popink,colback=poporangeL,
  before upper={\popboxhead{Objectifs de la leçon}{popink}{}}}
\newtcolorbox{prerequis}{popbase,colframe=popink,colback=popblueL,
  before upper={\popboxhead{Prérequis}{popblue}{}}}
\newtcolorbox{bilan}{popbase,colframe=popink,colback=white,boxrule=2.8pt,
  before upper={\popboxhead{Fiche bilan}{poporange}{L'essentiel à connaître pour l'examen}}}
% Figure encadrée avec légende
\newtcolorbox{popfigure}[1][]{enhanced,breakable=false,colback=white,colframe=popline,boxrule=1.4pt,arc=8pt,
  left=10pt,right=10pt,top=10pt,bottom=8pt,before skip=10pt,after skip=12pt,halign=center,
  lower separated=false,halign lower=center,
  fontlower=\popsemi\fontsize{9}{11.5}\selectfont\color{popmuted},
  #1}
\newcommand{\legende}[1]{\par\vspace{6pt}{\popsemi\fontsize{9}{11.5}\selectfont\color{popmuted}#1}}

% ============================================================
% Signature OnBuch+
% ============================================================
\newcommand{\poplogoinline}[1]{{\poplogo\fontsize{#1}{#1}\selectfont\color{popink}OnBuch\color{poporange}+}}
\newcommand{\signatureonbuch}{%
  \begin{tcolorbox}[enhanced,colback=popink,colframe=popink,boxrule=2.2pt,arc=10pt,
      drop shadow=poporange,left=14pt,right=14pt,top=10pt,bottom=10pt,before skip=0pt,after skip=0pt]
    \tikz[baseline=-0.7ex]\node[circle,fill=popgreen,draw=popcream,line width=1.3pt,minimum size=22pt,inner sep=0]{\popcheck{white}};%
    \hspace{9pt}\begin{minipage}[c]{0.62\linewidth}
      {\popxbold\fontsize{12}{13}\selectfont\color{popcream}Signé\ \textbullet\ {\poplogo OnBuch\color{poporange}+}}\par\vspace{2pt}
      {\raggedright\popsemi\fontsize{8}{10.5}\selectfont\color{popline}\MakeUppercase{Équipe pédagogique OnBuch+}\\\MakeUppercase{Conforme au programme officiel MINESEC}\par}
    \end{minipage}\hfill
    {\poplogo\fontsize{22}{22}\selectfont\color{popcream}OnBuch\color{poporange}+}
  \end{tcolorbox}%
}

% ============================================================
% Page de garde
% #1 titre · #2 matière · #3 niveau · #4 module · #5 n° leçon · #6 durée ·
% #7 description / objectifs (texte ou itemize)
% ============================================================
\newcommand{\pagegarde}[7]{%
  \thispagestyle{empty}
  \newgeometry{a4paper,margin=1.7cm,top=1.6cm,bottom=1.4cm}%
  \begin{tikzpicture}[remember picture,overlay]
    \fill[poporange!18] ([xshift=-3.2cm,yshift=-3.6cm]current page.north east) circle (4.6cm);
    \fill[poppurple!14] ([xshift=2.4cm,yshift=7.6cm]current page.south west) circle (3.8cm);
  \end{tikzpicture}%
  {\poplogo\fontsize{30}{30}\selectfont\color{popink}OnBuch\color{poporange}+}\hfill
  \raisebox{0.6ex}{\poppill{Cours signé \textbullet\ Premium}{popink}{popcream}}\par
  \vspace{1pt}\popsquiggle{poppurple}\par
  \vspace{8pt}
  \begin{tcolorbox}[enhanced,colback=poporange,colframe=popink,boxrule=2.8pt,arc=13pt,
      drop shadow=popink,left=18pt,right=18pt,top=12pt,bottom=13pt,after skip=10pt]
    \poppill{#2 \textbullet\ #3}{popink}{popcream}\hspace{5pt}\poppill{#4}{white}{popink}\par\vspace{8pt}
    {\popdisplay\fontsize{14}{14}\selectfont\color{popink}LEÇON #5}\par\vspace{4pt}
    {\raggedright\hyphenpenalty=10000\exhyphenpenalty=10000\popdisplay\fontsize{25}{27}\selectfont\color{popcream}\MakeUppercase{#1}\par}
    \vspace{8pt}
    {\popxbold\fontsize{9.5}{9.5}\selectfont\color{popink}\MakeUppercase{Durée conseillée : #6}}
  \end{tcolorbox}
  \begin{tcolorbox}[enhanced,colback=white,colframe=popink,boxrule=2.2pt,arc=10pt,
      drop shadow=popink,left=16pt,right=16pt,top=10pt,bottom=10pt,after skip=8pt]
    \poppill{Dans cette leçon}{poppurple}{white}\par\vspace{5pt}
    \popsemi\fontsize{9.3}{12.6}\selectfont\color{popink}#7
  \end{tcolorbox}
  \noindent\begin{minipage}[t]{0.487\linewidth}
    \begin{tcolorbox}[enhanced,colback=popgreen,colframe=popink,boxrule=2.2pt,arc=10pt,
        drop shadow=popink,left=11pt,right=11pt,top=10pt,bottom=10pt,height=3.1cm,valign=center]
      \tcbox[colback=white,colframe=popink,boxrule=1.3pt,arc=5pt,boxsep=0pt,nobeforeafter,
        left=3pt,right=3pt,top=3pt,bottom=3pt,valign=center]{\qrcode[height=1.9cm]{https://play.google.com/store/apps/details?id=com.luvvix.onbuch&hl=fr}}%
      \hspace{8pt}\begin{minipage}[c]{0.5\linewidth}
        {\popdisplay\fontsize{11}{12.5}\selectfont\color{white}\MakeUppercase{Télécharge OnBuch}}\par\vspace{4pt}
        {\raggedright\popsemi\fontsize{8.4}{11}\selectfont\color{white}Gratuit \textbullet\ Google Play\\Tuteur IA, annales, concours, résultats\par}
      \end{minipage}
    \end{tcolorbox}
  \end{minipage}\hfill
  \begin{minipage}[t]{0.487\linewidth}
    \begin{tcolorbox}[enhanced,colback=poppurple,colframe=popink,boxrule=2.2pt,arc=10pt,
        drop shadow=popink,left=11pt,right=11pt,top=10pt,bottom=10pt,height=3.1cm,valign=center]
      \tikz[baseline=-0.7ex]\node[circle,fill=white,draw=popink,line width=1.3pt,minimum size=1.6cm,inner sep=0]{\popdisplay\fontsize{24}{24}\selectfont\color{poppurple}+};%
      \hspace{8pt}\begin{minipage}[c]{0.6\linewidth}
        {\popdisplay\fontsize{11}{12.5}\selectfont\color{white}\MakeUppercase{Passe à OnBuch+}}\par\vspace{4pt}
        {\raggedright\popsemi\fontsize{8.4}{11}\selectfont\color{white}Tous les cours signés, Tuteur IA illimité, fiches PDF — onglet OnBuch+ de l'app\par}
      \end{minipage}
    \end{tcolorbox}
  \end{minipage}
  \vspace{8pt}
  \popcontenu
  \vfill
  \signatureonbuch
  \restoregeometry
  \newpage
}

% Rangée « Ce cours contient » (page de garde)
\newcommand{\poptile}[3]{% #1 couleur #2 gros texte #3 légende
  \begin{tcolorbox}[enhanced,colback=#1,colframe=popink,boxrule=2pt,arc=9pt,shadow={2.5pt}{-2.5pt}{0pt}{popink},
      left=6pt,right=6pt,top=9pt,bottom=9pt,halign=center,valign=center,height=1.75cm,width=\linewidth,nobeforeafter]
    {\popdisplay\fontsize{12.5}{13}\selectfont\color{white}\MakeUppercase{#2}}\par\vspace{3pt}
    {\centering\popsemi\fontsize{7.8}{9.5}\selectfont\color{white}#3\par}
  \end{tcolorbox}}
\newcommand{\popcontenu}{%
  \par\noindent{\popxboldls\fontsize{7.5}{7.5}\selectfont\color{popmuted}CE COURS SIGNÉ CONTIENT}\par\vspace{7pt}
  \noindent\begin{minipage}[t]{0.235\linewidth}\poptile{popblue}{Cours}{complet, illustré, pas à pas}\end{minipage}\hfill
  \begin{minipage}[t]{0.235\linewidth}\poptile{poporange}{Méthodes}{et exercices résolus}\end{minipage}\hfill
  \begin{minipage}[t]{0.235\linewidth}\poptile{poppink}{Exercices}{type examen + corrigés}\end{minipage}\hfill
  \begin{minipage}[t]{0.235\linewidth}\poptile{popgreen}{Fiche bilan}{l'essentiel à retenir}\end{minipage}\par}

% Sommaire
\newtcolorbox{sommaire}{enhanced,colback=white,colframe=popink,boxrule=2.2pt,arc=10pt,
  drop shadow=popink,left=18pt,right=18pt,top=13pt,bottom=13pt,before skip=6pt,after skip=20pt,
  before upper={\poppill{Sommaire}{poppurple}{white}\par\vspace{9pt}\popsemi\fontsize{10.6}{14.5}\selectfont\color{popink}}}

% Signature de fin de cours — poussée tout en bas de la dernière page
\newcommand{\findecours}{%
  \par\vfill\nopagebreak
  \begin{center}\popsquiggle{poporange}\end{center}\vspace{4pt}
  \signatureonbuch
}
\AtBeginDocument{\def\llbracket{\mathopen{[\mkern-3.5mu[}}\def\rrbracket{\mathclose{]\mkern-3.5mu]}}} % intervalles d'entiers (glyphe ⟦ absent de la police)
% Glyphes absents de Fira Math : redéfinis à partir de glyphes disponibles
\AtBeginDocument{%
  \def\setminus{\mathbin{\backslash}}\let\smallsetminus\setminus
  \def\vdots{\mathord{\vbox{\baselineskip3.2pt\lineskiplimit0pt\kern4pt\hbox{.}\hbox{.}\hbox{.}}}}%
  \def\ddots{\mathinner{\mkern1mu\raise6.5pt\hbox{.}\mkern2mu\raise3.6pt\hbox{.}\mkern2mu\raise0.7pt\hbox{.}\mkern1mu}}%
  \def\triangle{\mathord{\scalebox{0.9}{$\symup{\Delta}$}}}%
  \def\nabla{\mathord{\raisebox{\depth}{\scalebox{1}[-1]{$\symup{\Delta}$}}}}%
  \def\checkmark{\popcheck{popdark}}%
  \def\star{\mathbin{\ast}}\def\bigstar{\mathord{\ast}}%
  \def\diamond{\mathbin{\scalebox{0.6}{$\Diamond$}}}%
  \def\bigcup{\mathop{\mathchoice{\raisebox{-0.3em}{\scalebox{1.7}{$\cup$}}}{\scalebox{1.25}{$\cup$}}{\cup}{\cup}}\displaylimits}%
  \def\bigcap{\mathop{\mathchoice{\raisebox{-0.3em}{\scalebox{1.7}{$\cap$}}}{\scalebox{1.25}{$\cap$}}{\cap}{\cap}}\displaylimits}%
  \def\lesseqqgtr{\mathrel{\lessgtr}}\def\bigoplus{\mathop{\oplus}\displaylimits}%
}
\providecommand{\ssi}{si et seulement si} % abréviation fréquente
\AtBeginDocument{\providecommand{\wideparen}{\overparen}} % arc de cercle

%% FIGFIX-BEGIN v5 — correctifs globaux des figures (docs/onbuch-generation/tools/figfix)
\usepackage{adjustbox}\usepackage{etoolbox}\tcbuselibrary{raster}
% toute figure TikZ/pgfplots plus large que la ligne est réduite à la largeur disponible
\BeforeBeginEnvironment{tikzpicture}{\begin{adjustbox}{max width=\linewidth}}
\AfterEndEnvironment{tikzpicture}{\end{adjustbox}}
% légendes pgfplots lisibles par-dessus les courbes
\pgfplotsset{every axis/.append style={legend style={fill=white,fill opacity=0.92,text opacity=1,draw=black!15,inner sep=3pt}}}
% étiquettes d'arêtes (syntaxe edge["..."]) avec fond blanc
\tikzset{every edge quotes/.append style={fill=white,inner sep=1.2pt}}
%% FIGFIX-END
\begin{document}

\pagegarde{Limitation des dysfonctionnements des organes ou structures intervenant dans les mouvements réflexes}{SVT}{Tle D}{Module I : Le Monde Vivant — Activités réflexes}{N07}{8 h}{Cette leçon étudie la structure des centres nerveux et les mécanismes des mouvements réflexes (médullaires, myotatiques, conditionnels), établis à partir des expériences historiques de Bell et Magendie, de Pflüger et de la dégénérescence wallérienne. Elle permet d'expliquer les paralysies et les absences de réflexes observées en clinique.
\vspace{4pt}
\begin{multicols}{2}\small
\begin{itemize}
\item À la fin de cette leçon, je sais décrire les grandes lignes de la structure de l'encéphale et de la moelle épinière et identifier un neurone et ses différentes formes.
\item À la fin de cette leçon, je sais réaliser la dissection de l'encéphale d'un Mammifère et déméduller ou décérébrer un Batracien.
\item À la fin de cette leçon, je sais interpréter les résultats des expériences de Bell et Magendie et des expériences de dégénérescence wallérienne pour dégager le rôle des racines rachidiennes.
\item À la fin de cette leçon, je sais mettre en évidence un acte réflexe médullaire chez la grenouille et interpréter les expériences de Pflüger pour dégager les éléments qui y interviennent.
\item À la fin de cette leçon, je sais réaliser le schéma annoté d'un arc réflexe médullaire et d'un arc réflexe myotatique.
\item À la fin de cette leçon, je sais mettre en évidence les éléments intervenant dans le réflexe rotulien ou achilléen et expliquer son importance.
\end{itemize}\end{multicols}}

\begin{sommaire}
\begin{multicols}{2}
\begin{enumerate}
\item Structure des centres nerveux et du neurone
\item Rôle des racines rachidiennes : expériences de Bell et Magendie
\item Les réflexes médullaires : expériences de Pflüger et arc réflexe
\item Le réflexe myotatique : réflexe rotulien et achilléen
\item Les réflexes acquis ou conditionnels
\item Application : limitation et explication des dysfonctionnements
\item Activité d'intégration
\item Méthodes et exercices résolus
\item Exercices
\item Corrigés des exercices
\item Fiche bilan
\end{enumerate}
\end{multicols}
\end{sommaire}

\begin{objectifs}
\begin{itemize}
\item À la fin de cette leçon, je sais décrire les grandes lignes de la structure de l'encéphale et de la moelle épinière et identifier un neurone et ses différentes formes.
\item À la fin de cette leçon, je sais réaliser la dissection de l'encéphale d'un Mammifère et déméduller ou décérébrer un Batracien.
\item À la fin de cette leçon, je sais interpréter les résultats des expériences de Bell et Magendie et des expériences de dégénérescence wallérienne pour dégager le rôle des racines rachidiennes.
\item À la fin de cette leçon, je sais mettre en évidence un acte réflexe médullaire chez la grenouille et interpréter les expériences de Pflüger pour dégager les éléments qui y interviennent.
\item À la fin de cette leçon, je sais réaliser le schéma annoté d'un arc réflexe médullaire et d'un arc réflexe myotatique.
\item À la fin de cette leçon, je sais mettre en évidence les éléments intervenant dans le réflexe rotulien ou achilléen et expliquer son importance.
\item À la fin de cette leçon, je sais expliquer brièvement les mécanismes et l'importance des réflexes acquis ou conditionnels.
\item À la fin de cette leçon, je sais expliquer certains cas de paralysie ou d'absence de réflexe chez des malades (poliomyélite, section médullaire, tétanos).
\end{itemize}
\end{objectifs}

\begin{prerequis}
\begin{itemize}
\item Notion de message nerveux et de nerf (classe de Première : relation du système nerveux avec les mouvements)
\item Structure générale du muscle squelettique et notion de contraction musculaire
\item Organisation générale du système nerveux (nerfs, moelle, cerveau) vue en classe de Seconde
\item Notion de stimulus et de réponse (irritabilité du vivant)
\item Techniques de base de dissection et d'observation au microscope
\end{itemize}
\end{prerequis}

\newpage

\coursec{Structure des centres nerveux et du neurone}

Au stade omnisports de Bafoussam, un joueur reçoit un coup sur le genou lors d'un match de championnat. Sur la touche, le médecin du club sort un petit marteau à réflexes et tape doucement sous sa rotule : aussitôt, la jambe du joueur se détend vers l'avant, sans qu'il l'ait voulu. Le médecin sourit : « Ton système nerveux fonctionne bien. » Quelques mois plus tard, à l'hôpital central de Yaoundé, un moto-taximan victime d'un accident de la route ne sent plus ses jambes et ne peut plus les bouger, alors que ses muscles, à l'examen, semblent parfaitement intacts. Comment expliquer que le médecin puisse tester le système nerveux en tapant simplement sur un genou ? Pourquoi une lésion du dos peut-elle priver quelqu'un de tout mouvement et de toute sensibilité ?

Pour répondre à ces questions, il faut d'abord connaître le \emph{matériel} dont dispose le système nerveux : les \cle{centres nerveux} (encéphale et moelle épinière), les \cle{nerfs} qui les relient aux organes, et les \cle{neurones}, cellules qui fabriquent et conduisent le message nerveux. C'est l'objet de cette première partie : observer, disséquer, décrire, afin de poser les bases anatomiques indispensables à la compréhension des réflexes.

\courssub{Observation macroscopique des centres nerveux}

\textbf{Situation-problème.} Sur une table de dissection du laboratoire, on dépose une tête de mouton achetée à l'abattoir de Mvog-Mbi. À l'œil nu, on ne voit rien du système nerveux : il est enfermé dans les os. Comment atteindre et décrire les centres nerveux ?

\begin{methode}[Dissection de l'encéphale d'un Mammifère (mouton ou chèvre)]
\begin{enumerate}
\item \textbf{Dégager la boîte crânienne} : à l'aide d'une scie ou d'un ciseau fort, pratiquer deux découpes latérales du crâne, de part et d'autre de la ligne médiane, en prenant soin de ne pas enfoncer l'outil trop profondément pour ne pas léser l'encéphale.
\item \textbf{Retirer la calotte crânienne} et dégager délicatement les méninges (enveloppes protectrices) qui adhèrent à l'os.
\item \textbf{Sectionner le bulbe rachidien} au niveau du trou occipital, le plus en arrière possible, puis \textbf{extraire l'encéphale} en le soulevant doucement et en sectionnant les nerfs crâniens qui le retiennent.
\item \textbf{Identifier les parties} : sur la face dorsale, observer le cerveau et le cervelet ; retourner l'encéphale pour observer, sur la face ventrale, le bulbe rachidien prolongé par la moelle épinière.
\end{enumerate}
\end{methode}

\begin{experience}[Observation de l'encéphale de mouton]
\textbf{Matériel} : tête de mouton fraîche, scie, ciseaux, pinces, gants, plateau de dissection.\par
\textbf{Observations} :
\begin{itemize}
\item Le \cle{cerveau} est la partie la plus volumineuse ; sa surface présente de nombreux plis, les \emph{circonvolutions}, séparés par des sillons. Il est divisé en deux hémisphères cérébraux.
\item Le \cle{cervelet}, situé en arrière et en dessous du cerveau, est plus petit et présente des plis fins et parallèles (aspect feuilleté).
\item Le \cle{bulbe rachidien} est un renflement situé sous le cervelet ; il se prolonge vers le bas par un long cordon cylindrique : la \cle{moelle épinière}, logée dans le canal rachidien de la colonne vertébrale.
\end{itemize}
\textbf{Interprétation} : l'encéphale (cerveau, cervelet, bulbe rachidien) et la moelle épinière forment le \emph{système nerveux central}, ou \cle{centres nerveux}. Ils sont protégés par les os (crâne, vertèbres), ce qui traduit leur importance vitale et leur fragilité.
\end{experience}

\begin{aretenir}[Grandes lignes de la structure des centres nerveux]
Les centres nerveux comprennent :
\begin{itemize}
\item l'\cle{encéphale}, logé dans la boîte crânienne, qui comprend le \textbf{cerveau} (siège des activités volontaires et de l'intelligence), le \textbf{cervelet} (coordination des mouvements, équilibre) et le \textbf{bulbe rachidien} (contrôle des fonctions vitales : rythme cardiaque, respiration) ;
\item la \cle{moelle épinière}, logée dans le canal rachidien, centre des réflexes et voie de passage des messages entre l'encéphale et la périphérie.
\end{itemize}
\end{aretenir}

\begin{savaistu}[La moelle épinière humaine en chiffres]
Chez l'Homme, la moelle épinière mesure environ \SI{45}{cm} de long pour un diamètre d'environ \SI{1}{cm} ; elle ne descend que jusqu'à la deuxième vertèbre lombaire (L2). Il s'en détache \textbf{31 paires de nerfs rachidiens} : 8 paires cervicales, 12 paires thoraciques, 5 paires lombaires, 5 paires sacrées et 1 paire coccygienne. C'est la rupture de ce fragile cordon, lors d'un accident de moto ou d'une chute, qui explique les paralysies des deux jambes (paraplégie) fréquentes dans nos services de traumatologie.
\end{savaistu}

\courssub{Coupe transversale de la moelle épinière et racines rachidiennes}

\textbf{Observation.} On réalise une coupe fine et transversale de la moelle épinière de mouton et on l'examine à la loupe binoculaire. Deux régions de couleurs différentes apparaissent nettement.

\begin{itemize}
\item Au centre, une zone grisâtre, la \cle{substance grise}, dessine une forme de \textbf{H} (ou de papillon aux ailes déployées). Elle présente quatre saillies : deux \emph{cornes dorsales} (postérieures), fines, et deux \emph{cornes ventrales} (antérieures), plus larges. Au milieu du H se trouve un petit orifice, le \cle{canal épendymaire}, qui contient le liquide céphalo-rachidien.
\item En périphérie, une zone blanchâtre, la \cle{substance blanche}, entoure complètement la substance grise.
\end{itemize}

De part et d'autre de la moelle se rattachent les \cle{racines rachidiennes} :
\begin{itemize}
\item la \cle{racine dorsale} (postérieure) pénètre dans la corne dorsale ; elle porte un renflement, le \cle{ganglion spinal} ;
\item la \cle{racine ventrale} (antérieure) sort de la corne ventrale ; elle \textbf{ne porte pas de ganglion} ;
\item les deux racines \textbf{convergent} peu après leur sortie de la moelle et fusionnent pour former le \cle{nerf rachidien}, qui se dirige vers les organes (peau, muscles).
\end{itemize}

\begin{popfigure}
\begin{center}
\begin{tikzpicture}[scale=1.0, every node/.style={font=\small}]
% Substance blanche (moelle)
\fill[popcream, draw=popdark, thick] (0,0) ellipse (2.6 and 2.2);
% Substance grise en H (papillon)
\fill[popblue!45, draw=popdark] 
  (-0.35,1.5) .. controls (-1.0,1.5) and (-1.1,0.9) .. (-0.85,0.35)
  -- (-0.35,0.15) -- (-0.35,-0.15) -- (-1.15,-0.35)
  .. controls (-1.5,-0.6) and (-1.3,-1.5) .. (-0.75,-1.5)
  .. controls (-0.35,-1.5) and (-0.25,-0.7) .. (-0.15,-0.35)
  -- (0.15,-0.35)
  .. controls (0.25,-0.7) and (0.35,-1.5) .. (0.75,-1.5)
  .. controls (1.3,-1.5) and (1.5,-0.6) .. (1.15,-0.35)
  -- (0.35,-0.15) -- (0.35,0.15) -- (0.85,0.35)
  .. controls (1.1,0.9) and (1.0,1.5) .. (0.35,1.5)
  .. controls (0.15,1.5) and (0.1,0.6) .. (0.1,0.35) -- (-0.1,0.35)
  .. controls (-0.1,0.6) and (-0.15,1.5) .. (-0.35,1.5) -- cycle;
% Canal épendymaire
\fill[white, draw=popdark] (0,0) circle (0.09);
% Racine dorsale gauche + ganglion
\draw[poporange, very thick] (-0.7,1.6) .. controls (-1.2,2.6) and (-1.6,3.0) .. (-2.4,3.2);
\fill[poporange] (-1.65,2.95) circle (0.22);
% Racine ventrale gauche
\draw[popgreen!70!black, very thick] (-0.9,-1.4) .. controls (-1.4,-2.3) and (-1.9,-2.7) .. (-2.35,-2.9);
% Nerf rachidien gauche
\draw[popdark, very thick] (-2.4,3.2) .. controls (-3.1,3.35) and (-3.4,2.2) .. (-3.5,0.2);
\draw[popdark, very thick] (-2.35,-2.9) .. controls (-3.0,-3.2) and (-3.6,-1.6) .. (-3.5,0.2);
% Racine dorsale droite + ganglion
\draw[poporange, very thick] (0.7,1.6) .. controls (1.2,2.6) and (1.6,3.0) .. (2.4,3.2);
\fill[poporange] (1.65,2.95) circle (0.22);
% Racine ventrale droite
\draw[popgreen!70!black, very thick] (0.9,-1.4) .. controls (1.4,-2.3) and (1.9,-2.7) .. (2.35,-2.9);
% Nerf rachidien droit
\draw[popdark, very thick] (2.4,3.2) .. controls (3.1,3.35) and (3.4,2.2) .. (3.5,0.2);
\draw[popdark, very thick] (2.35,-2.9) .. controls (3.0,-3.2) and (3.6,-1.6) .. (3.5,0.2);
% Légendes fléchées
\draw[->, popmuted] (0.09,0.1) -- (1.3,0.9) node[right, align=left]{1. Canal épendymaire};
\draw[->, popmuted] (-0.6,0.6) -- (-2.2,1.4) node[left, align=right]{2. Substance grise (H)};
\draw[->, popmuted] (2.0,-0.9) -- (3.1,-1.1) node[right, align=left]{3. Substance blanche};
\node[poporange, align=left] at (2.9,3.6) {4. Racine dorsale};
\draw[->, poporange] (2.85,3.45) -- (2.45,3.25);
\node[poporange, align=left] at (0.4,3.7) {5. Ganglion spinal};
\draw[->, poporange] (0.9,3.55) -- (1.55,3.05);
\node[popgreen!60!black, align=left] at (3.0,-3.3) {6. Racine ventrale};
\draw[->, popgreen!60!black] (2.95,-3.15) -- (2.4,-2.9);
\node[popdark, align=left] at (-4.6,0.9) {7. Nerf rachidien};
\draw[->, popdark] (-3.4,0.75) -- (-3.5,0.25);
\end{tikzpicture}
\end{center}
\legende{Coupe transversale de la moelle épinière et origine d'un nerf rachidien. 1 : canal épendymaire ; 2 : substance grise centrale en forme de H (cornes dorsales fines, cornes ventrales larges) ; 3 : substance blanche périphérique ; 4 : racine dorsale (postérieure) ; 5 : ganglion spinal ; 6 : racine ventrale (antérieure), sans ganglion ; 7 : nerf rachidien formé par la convergence des deux racines.}
\end{popfigure}

\begin{attention}[Erreur fréquente : inverser substance grise et substance blanche]
Dans la \textbf{moelle épinière}, la substance grise est \textbf{centrale} (le H) et la substance blanche est \textbf{périphérique}. C'est l'\emph{inverse} dans le cerveau, où la substance grise est périphérique (elle forme le \emph{cortex}) et la substance blanche centrale. Au Baccalauréat, beaucoup d'élèves perdent des points en recopiant la disposition cérébrale sur un schéma de moelle. Retiens : \emph{dans la moelle, le gris est au cœur}.
\end{attention}

\courssub{Coupe transversale d'un nerf}

\textbf{Observation.} On prélève un nerf (par exemple le nerf sciatique d'un mammifère) et on en réalise une coupe transversale que l'on colore et observe au microscope optique. Le nerf apparaît comme un \emph{câble électrique à plusieurs conducteurs} : il est constitué de \textbf{faisceaux de fibres nerveuses}, chaque fibre étant un axone entouré de sa gaine. Trois enveloppes de tissu conjonctif organisent l'ensemble :
\begin{itemize}
\item l'\cle{endonèvre} entoure chaque fibre nerveuse individuellement ;
\item le \cle{périnèvre} regroupe les fibres en faisceaux ;
\item l'\cle{épinèvre} enveloppe le nerf tout entier et soude les faisceaux entre eux.
\end{itemize}

\begin{popfigure}
\begin{center}
\begin{tikzpicture}[scale=1.0, every node/.style={font=\small}]
% Épinèvre
\fill[popgold!25, draw=popdark, very thick] (0,0) circle (2.5);
% Faisceaux
\foreach \cx/\cy in {-1.05,0.75 / 0.9,0.9 / 0.1,-0.55 / -1.35,-0.85 / 1.35,-0.75}{
  \fill[popblue!20, draw=popblue, thick] (\cx,\cy) circle (0.72);
}
% Fibres dans les faisceaux
\foreach \cx/\cy in {-1.05,0.75 / 0.9,0.9 / 0.1,-0.55 / -1.35,-0.85 / 1.35,-0.75}{
  \foreach \dx/\dy in {-0.3,0.25 / 0.28,0.28 / 0,0 / -0.3,-0.28 / 0.3,-0.25 / 0.02,0.45 / 0.02,-0.48}{
    \fill[white, draw=poporange, thick] ({\cx+\dx},{\cy+\dy}) circle (0.13);
    \fill[poporange] ({\cx+\dx},{\cy+\dy}) circle (0.045);
  }
}
% Légendes
\draw[->, popmuted] (2.45,0.5) -- (3.4,1.3) node[right, align=left]{1. Épinèvre};
\draw[->, popmuted] (1.62,0.9) -- (3.4,0.4) node[right, align=left]{2. Périnèvre (faisceau)};
\draw[->, popmuted] (0.4,-0.55) -- (3.4,-0.6) node[right, align=left]{3. Fibre nerveuse (axone)};
\draw[->, popmuted] (-1.35,-1.57) -- (-2.6,-2.2) node[left, align=right]{4. Endonèvre};
\end{tikzpicture}
\end{center}
\legende{Coupe transversale d'un nerf. 1 : épinèvre, gaine conjonctive externe ; 2 : périnèvre entourant chaque faisceau de fibres ; 3 : fibres nerveuses (axones vus en coupe, avec leur gaine de myéline) ; 4 : endonèvre entre les fibres. Un nerf est donc un ensemble de faisceaux de fibres nerveuses réunis par du tissu conjonctif.}
\end{popfigure}

\begin{definition}[Nerf]
Un \cle{nerf} est un ensemble de \textbf{fibres nerveuses} (axones) réunies en \textbf{faisceaux} par du tissu conjonctif (endonèvre, périnèvre, épinèvre), reliant les centres nerveux aux organes (récepteurs et effecteurs).
\end{definition}

\courssub{Le neurone : unité structurale et fonctionnelle du système nerveux}

\textbf{Observation.} L'examen au microscope de lamelles de moelle épinière colorées (méthode de Golgi) révèle des cellules étoilées pourvues de longs prolongements : ce sont les \cle{neurones}.

\begin{definition}[Neurone]
Un \cle{neurone} est une cellule nerveuse spécialisée dans la \textbf{réception}, la \textbf{conduction} et la \textbf{transmission} du message nerveux. Il comprend :
\begin{itemize}
\item un \cle{corps cellulaire} (ou péricaryon) contenant le noyau et la majeure partie du cytoplasme ;
\item des \cle{dendrites}, prolongements courts, ramifiés et nombreux, qui \emph{reçoivent} l'information et la conduisent \emph{vers} le corps cellulaire ;
\item un \cle{axone} (ou cylindraxe), prolongement unique et long (jusqu'à \SI{1}{m} pour les neurones des membres), qui conduit l'influx nerveux \emph{du corps cellulaire vers} d'autres cellules.
\end{itemize}
L'axone est souvent entouré d'une \cle{gaine de myéline}, substance blanche et graisseuse déposée par des cellules annexes (cellules de Schwann dans le système nerveux périphérique, oligodendrocytes dans le système nerveux central). Cette gaine est interrompue à intervalles réguliers par des étranglements, les \cle{nœuds de Ranvier}, qui accélèrent la conduction du message nerveux (jusqu'à \SI{120}{m.s^{-1}} dans les fibres myélinisées).
\end{definition}

\begin{propriete}[Substance grise et substance blanche : que contiennent-elles ?]
La couleur des régions des centres nerveux s'explique par leur contenu cellulaire :
\begin{itemize}
\item la \cle{substance grise} contient les \textbf{corps cellulaires} des neurones (et les dendrites) ;
\item la \cle{substance blanche} contient les \textbf{fibres myélinisées} (axones gainés de myéline, substance blanchâtre).
\end{itemize}
Ainsi, dans la moelle épinière, les corps cellulaires sont centraux (H gris) et les axones qui montent ou descendent le long de la moelle forment la substance blanche périphérique.
\end{propriete}

\textbf{Les différentes formes de neurones.} Tous les neurones ne se ressemblent pas : leur forme est adaptée à leur fonction dans la chaîne réflexe. On en distingue classiquement trois types.

\begin{itemize}
\item Le \cle{neurone sensitif} (ou neurone afférent) est \textbf{unipolaire} (plus précisément pseudo-unipolaire) : un unique prolongement quitte son corps cellulaire puis se divise en deux branches en T, l'une venant d'un récepteur sensoriel (peau, muscle), l'autre pénétrant dans la moelle par la racine dorsale. Son corps cellulaire est logé \textbf{hors des centres nerveux}, dans le \emph{ganglion spinal} — ce qui explique le renflement de la racine dorsale.
\item Le \cle{neurone moteur} (ou motoneurone) est \textbf{multipolaire} : son corps cellulaire, situé dans la \emph{corne ventrale} de la substance grise, porte de nombreuses dendrites et un long axone qui sort par la racine ventrale et se termine sur un muscle (jonction neuromusculaire).
\item Le \cle{neurone d'association} (ou interneurone) est un petit neurone multipolaire dont le corps cellulaire et tous les prolongements restent \textbf{dans la substance grise} des centres nerveux ; il fait la liaison entre le neurone sensitif et le neurone moteur.
\end{itemize}

\begin{popfigure}
\begin{center}
\begin{tikzpicture}[scale=0.92, every node/.style={font=\small}]
% ============ Neurone sensitif (unipolaire) ============
\begin{scope}[xshift=-5.2cm]
\node[font=\small\bfseries, popblue] at (0,3.1) {Neurone sensitif (unipolaire)};
% corps dans ganglion
\fill[popblue!30, draw=popdark, thick] (0,1.2) circle (0.45);
\fill[popdark] (0,1.2) circle (0.12);
% prolongement en T
\draw[popblue, very thick] (0,1.65) -- (0,2.3);
\draw[popblue, very thick] (0,2.3) -- (-1.5,2.3);
\draw[popblue, very thick] (0,2.3) -- (1.5,2.3);
\draw[popblue, very thick] (1.5,2.3) -- (1.5,0.4);
% récepteur
\draw[poporange, thick] (-1.5,2.3) -- (-1.5,2.7);
\node[poporange, align=center] at (-1.5,2.95) {Récepteur\\(peau)};
% ganglion
\draw[popmuted, dashed, thick] (0,1.2) circle (0.75);
\node[popmuted, align=center] at (0,0.15) {Corps cellulaire\\dans le ganglion spinal};
\node[popblue, align=center] at (1.5,0.1) {vers la moelle\\(racine dorsale)};
\end{scope}
% ============ Neurone moteur (multipolaire) ============
\begin{scope}[xshift=0cm]
\node[font=\small\bfseries, popgreen!60!black] at (0,3.1) {Neurone moteur (multipolaire)};
% dendrites
\foreach \a in {120,150,60,30,90}{
  \draw[popgreen!70!black, thick] (0,0.9) -- ({0.9*cos(\a)},{0.9+0.9*sin(\a)});
}
% corps
\fill[popgreen!30, draw=popdark, thick] (0,0.9) circle (0.5);
\fill[popdark] (0,0.9) circle (0.13);
% axone avec myéline
\draw[popgreen!70!black, very thick] (0,0.4) -- (0,-2.3);
\foreach \y in {-0.1,-0.7,-1.3,-1.9}{
  \draw[popgold, line width=3.2pt] (-0.09,\y) -- (0.09,\y);
}
\node[popgold!70!black, align=left] at (1.15,-1.0) {gaine de\\myéline};
\node[popmuted, align=left] at (1.15,-1.9) {nœud de\\Ranvier};
\draw[->, popmuted] (1.0,-1.75) -- (0.12,-1.6);
\node[popgreen!60!black, align=center] at (0,-2.75) {vers le muscle\\(racine ventrale)};
\node[popmuted, align=center] at (-1.7,0.9) {Corps dans la\\corne ventrale};
\end{scope}
% ============ Interneurone ============
\begin{scope}[xshift=5.0cm]
\node[font=\small\bfseries, poppurple] at (0,3.1) {Neurone d'association};
\foreach \a in {40,100,160,220,280,340}{
  \draw[poppurple, thick] (0,1.0) -- ({0.85*cos(\a)},{1.0+0.85*sin(\a)});
}
\fill[poppurple!30, draw=popdark, thick] (0,1.0) circle (0.45);
\fill[popdark] (0,1.0) circle (0.12);
\node[poppurple, align=center] at (0,-0.35) {Petit neurone multipolaire,\\entièrement situé dans\\la substance grise};
\end{scope}
\end{tikzpicture}
\end{center}
\legende{Les trois types de neurones. À gauche : neurone sensitif unipolaire, dont le corps cellulaire est dans le ganglion spinal ; au centre : neurone moteur multipolaire, dont le corps est dans la corne ventrale de la moelle, avec un axone myélinisé interrompu par les nœuds de Ranvier ; à droite : neurone d'association (interneurone), petit neurone multipolaire confiné dans la substance grise.}
\end{popfigure}

\courssub{L'arc réflexe médullaire : organisation fonctionnelle}

Les trois types de neurones ne flottent pas au hasard : ils s'enchaînent pour former un circuit fonctionnel, l'\cle{arc réflexe}. C'est le schéma minimal de tout réflexe médullaire.

\begin{experience}[Schéma fonctionnel de l'arc réflexe médullaire]
\textbf{Construction du schéma.} À partir de la coupe de moelle et des types de neurones, on représente le cheminement de l'influx nerveux lors d'un réflexe simple (ex. : réflexe de retrait de la main piquée par une épine).
\begin{enumerate}
\item \textbf{Récepteur} (terminaison sensitive libre dans la peau) $\rightarrow$ 
\item \textbf{Neurone sensitif} : corps dans le ganglion spinal, axone entrant par la \textbf{racine dorsale}, se terminant dans la \textbf{corne dorsale} (synapse).
\item \textbf{Neurone d'association (interneurone)} : corps et prolongements dans la \textbf{substance grise} (corne dorsale / zone intermédiaire), relie le neurone sensitif au neurone moteur (synapse).
\item \textbf{Neurone moteur} : corps dans la \textbf{corne ventrale}, axone sortant par la \textbf{racine ventrale}, puis le \textbf{nerf rachidien} $\rightarrow$ 
\item \textbf{Effecteur} (muscle squelettique : contraction).
\end{enumerate}
\end{experience}

\begin{popfigure}
\begin{center}
\begin{tikzpicture}[scale=0.85, every node/.style={font=\small}, >=Stealth]
% Moelle (simplifiée)
\fill[popcream, draw=popdark, thick] (-4,-2.5) rectangle (4,2.5);
\fill[popblue!40, draw=popdark] (-0.5,1.5) -- (-1.2,0) -- (-0.5,-0.2) -- (0.5,-0.2) -- (1.2,0) -- (0.5,1.5) -- cycle; % H simplifié
\fill[white, draw=popdark] (0,0) circle (0.15); % Canal
% Neurone sensitif
\fill[popblue!30, draw=popblue, thick] (-3.5,1.8) circle (0.3); % Ganglion
\fill[popdark] (-3.5,1.8) circle (0.08);
\draw[popblue, very thick, ->] (-3.5,1.5) -- (-2.5,1.5) node[midway, above] {Influx};
\draw[popblue, very thick] (-3.5,2.1) -- (-3.5,2.5);
\node[poporange, align=center] at (-3.5,2.8) {Récepteur};
% Interneurone
\fill[poppurple!30, draw=poppurple, thick] (-1.0,0.5) circle (0.25);
\fill[popdark] (-1.0,0.5) circle (0.06);
\draw[poppurple, thick] (-1.0,0.5) -- (-1.5,1.0);
\draw[poppurple, thick] (-1.0,0.5) -- (-0.5,0.0);
% Neurone moteur
\fill[popgreen!30, draw=popgreen!60!black, thick] (1.0,-1.0) circle (0.3);
\fill[popdark] (1.0,-1.0) circle (0.08);
\draw[popgreen!60!black, thick] (0.5,-0.5) -- (1.0,-1.0); % dendrite
\draw[popgreen!60!black, very thick, ->] (1.3,-1.0) -- (3.5,-1.0) node[midway, above] {Influx};
\node[popgreen!60!black, align=center] at (3.8,-1.0) {Effecteur\\(Muscle)};
% Synapses
\draw[popmuted, decorate, decoration={snake, amplitude=1pt, segment length=3pt}] (-2.5,1.5) -- (-1.5,0.8) node[midway, left, font=\tiny] {Synapse};
\draw[popmuted, decorate, decoration={snake, amplitude=1pt, segment length=3pt}] (-0.5,0.0) -- (0.5,-0.5) node[midway, right, font=\tiny] {Synapse};
% Légendes structures
\node[align=center, font=\tiny] at (-3.5, -2.0) {Ganglion spinal\\(Racine dorsale)};
\node[align=center, font=\tiny] at (0, 2.0) {Corne dorsale};
\node[align=center, font=\tiny] at (0, -2.0) {Corne ventrale\\(Racine ventrale)};
% Flèches sens
\draw[->, popdark, thick] (-3.5, -2.3) -- (-3.5, -2.8) node[below, font=\tiny] {Sens de l'influx};
\end{tikzpicture}
\end{center}
\legende{Schéma fonctionnel de l'arc réflexe médullaire (monosynaptique ou polysynaptique). L'influx va du récepteur vers l'effecteur via : 1. Neurone sensitif (bleu) entrant par la racine dorsale ; 2. Interneurone (violet) dans la substance grise ; 3. Neurone moteur (vert) sortant par la racine ventrale. Les synapses sont représentées par des traits ondulés.}
\end{popfigure}

\begin{methode}[Préparation de la grenouille pour l'étude des réflexes médullaires (Savoir-faire)]
\begin{enumerate}
\item \textbf{Démédullation} : sectionner la moelle épinière juste derrière le bulbe rachidien (entre le crâne et la première vertèbre). La grenouille perd tout réflexe médullaire mais conserve les réflexes bulbo-médullaires (respiration, battements cardiaques).
\item \textbf{Décérébration} : sectionner le cerveau au-dessus du bulbe (au niveau du mésencéphale). La grenouille conserve les réflexes médullaires et bulbo-médullaires, mais perd les comportements volontaires et complexes.
\item \textbf{Intérêt} : ces préparations permettent d'isoler le fonctionnement de la moelle épinière (centre des réflexes) de l'influence inhibitrice ou facilitatrice de l'encéphale.
\end{enumerate}
\end{methode}

\begin{exoresolu}[Identifier une structure sur une coupe de moelle]
\textbf{Énoncé.} Sur une coupe transversale de moelle épinière de rat observée à la loupe, un élève désigne la zone périphérique blanchâtre comme « substance grise » et affirme que le ganglion spinal se trouve sur la racine ventrale. Relever les deux erreurs et les corriger en les justifiant.
\tcblower
\textbf{Solution.} 
\begin{enumerate}
\item \emph{Erreur 1} : la zone périphérique n'est pas la substance grise. Dans la moelle épinière, la \textbf{substance grise est centrale} et dessine un H (ou un papillon) ; elle contient les corps cellulaires des neurones. La zone périphérique blanchâtre est la \textbf{substance blanche}, constituée d'axones entourés de myéline (substance blanche et graisseuse).
\item \emph{Erreur 2} : le ganglion spinal ne se trouve pas sur la racine ventrale. Il est porté par la \textbf{racine dorsale} (postérieure), car il contient les corps cellulaires des neurones sensitifs. La racine ventrale, formée des axones des motoneurones, ne porte \textbf{aucun ganglion}.
\end{enumerate}
\textbf{Conclusion} : dans la moelle, gris = central ; ganglion = racine dorsale.
\end{exoresolu}

\begin{attention}[Pièges de rédaction à éviter]
\begin{itemize}
\item Ne pas confondre \textbf{nerf} et \textbf{neurone} : le neurone est une \emph{cellule} ; le nerf est un \emph{ensemble de fibres} (axones de nombreux neurones) réunies par du tissu conjonctif. Un nerf est au neurone ce qu'un câble électrique est à un fil de cuivre.
\item Ne pas écrire que la racine ventrale « entre » dans la moelle : elle en \emph{sort} (axones moteurs), tandis que la racine dorsale y \emph{pénètre} (fibres sensitives). La justification fonctionnelle de cette organisation sera établie dans la section suivante par les expériences de Bell et Magendie.
\item Le canal épendymaire n'est pas un « vaisseau sanguin » : il contient le liquide céphalo-rachidien.
\item Ne pas confondre \emph{neurone sensitif} (voie afférente générale) et \emph{neurone sensoriel} (spécialisé pour un sens précis comme la vue, l'ouïe) ; au programme, on utilise \textbf{neurone sensitif} pour l'arc réflexe.
\end{itemize}
\end{attention}

\begin{exercice}[Pour t'entraîner]
\begin{enumerate}
\item Réaliser le schéma annoté d'une coupe transversale de moelle épinière montrant la substance grise, la substance blanche, le canal épendymaire, les deux racines et le ganglion spinal.
\item Compléter le tableau suivant : type de neurone / forme (uni- ou multipolaire) / localisation du corps cellulaire / rôle.
\item Expliquer pourquoi la substance blanche de la moelle est blanche, alors que la substance grise est grise.
\item Dessiner le schéma fonctionnel de l'arc réflexe médullaire en indiquant le sens de l'influx nerveux, la nature des synapses et les racines rachidiennes concernées.
\end{enumerate}
\end{exercice}

\begin{corrige}[Exercice : pour t'entraîner]
\begin{enumerate}
\item Le schéma doit reproduire le H gris central (cornes dorsales fines, cornes ventrales larges), le canal épendymaire au centre, la substance blanche tout autour, la racine dorsale portant le ganglion spinal, la racine ventrale sans ganglion, et la convergence des deux racines en un nerf rachidien (se reporter à la figure de la leçon).
\item Neurone sensitif : unipolaire (pseudo-unipolaire), corps dans le ganglion spinal, transmet l'information sensorielle des récepteurs vers les centres nerveux. Neurone moteur : multipolaire, corps dans la corne ventrale de la moelle, commande la contraction musculaire. Neurone d'association : multipolaire, corps dans la substance grise, relie neurones sensitifs et moteurs.
\item La substance blanche doit sa couleur à la \textbf{myéline}, gaine graisseuse blanchâtre qui entoure les axones ; la substance grise est constituée de corps cellulaires et de dendrites dépourvus de myéline, d'où sa teinte grisâtre.
\item Voir le schéma fonctionnel de l'arc réflexe ci-dessus : Récepteur $\rightarrow$ Neurone sensitif (racine dorsale) $\rightarrow$ Interneurone (substance grise) $\rightarrow$ Neurone moteur (racine ventrale) $\rightarrow$ Effecteur. Sens de l'influx : centripète (vers le centre) puis centrifuge (vers la périphérie).
\end{enumerate}
\end{corrige}

\begin{aretenir}[L'essentiel de la section]
\begin{itemize}
\item Les centres nerveux = encéphale (cerveau, cervelet, bulbe rachidien) + moelle épinière.
\item Dans la moelle : substance grise \textbf{centrale} en H (corps cellulaires), substance blanche \textbf{périphérique} (axones myélinisés), canal épendymaire au centre.
\item Racine dorsale = sensitive, avec ganglion spinal ; racine ventrale = motrice, sans ganglion ; leur union forme le nerf rachidien.
\item Un nerf = faisceaux de fibres nerveuses + gaines conjonctives (endonèvre, périnèvre, épinèvre).
\item Le neurone = corps cellulaire + dendrites + axone (souvent myélinisé, avec nœuds de Ranvier) ; trois formes : sensitif (unipolaire), moteur (multipolaire), d'association (interneurone).
\item L'\textbf{arc réflexe} est la chaîne fonctionnelle : Récepteur $\rightarrow$ Neurone sensitif $\rightarrow$ Interneurone(s) $\rightarrow$ Neurone moteur $\rightarrow$ Effecteur. Il est réalisé au niveau de la moelle épinière pour les réflexes médullaires.
\end{itemize}
\end{aretenir}

Nous disposons désormais de la carte anatomique du système nerveux. Mais connaître les structures ne suffit pas : comment savoir quelle racine conduit la sensibilité et laquelle commande le mouvement ? C'est ce que les expériences historiques de Bell et Magendie, que nous étudierons dans la section suivante, ont permis d'établir — et c'est précisément ce que le médecin de Bafoussam teste lorsqu'il tape sous la rotule de son joueur.

\coursec{Rôle des racines rachidiennes : expériences de Bell et Magendie}

Dans la section précédente, tu as découvert la structure des centres nerveux (encéphale et moelle épinière) et celle du neurone, unité anatomique et fonctionnelle du tissu nerveux. Tu sais notamment que la moelle épinière est reliée à la périphérie du corps par 31 paires de nerfs rachidiens, chacun naissant de la réunion de deux racines : une \cle{racine dorsale} (postérieure), renflée par un \cle{ganglion spinal}, et une \cle{racine ventrale} (antérieure), dépourvue de ganglion. Une question fondamentale se pose alors, et elle a agité toute la physiologie du XIX\textsuperscript{e} siècle : le nerf rachidien conduit-il un seul type de message nerveux ? Et si deux types de messages circulent, quelle racine conduit lequel ? C'est à cette question que répondent les célèbres expériences de \cle{Bell} et \cle{Magendie}, que nous allons analyser pas à pas, comme au laboratoire.

\courssub{Problématique et hypothèses}

\courssub{Observation de départ et problème}

Observe la situation clinique suivante, fréquente dans les services de neurologie des hôpitaux camerounais : deux patients victimes d'un accident de moto présentent tous deux des troubles au niveau d'un membre. Le premier ne sent plus les piqûres ni le contact sur la jambe gauche, mais il la bouge normalement. Le second, au contraire, ne peut plus bouger la jambe gauche, alors qu'il perçoit parfaitement les piqûres. Comment expliquer que deux lésions de la même région produisent des troubles aussi différents ?

L'observation anatomique donne un indice précieux : le nerf rachidien, qui est \emph{mixte} au niveau du tronc nerveux (il contient des fibres de plusieurs types), naît de la moelle par \textbf{deux racines distinctes}. Il est donc légitime de formuler le problème suivant : \emph{chaque racine rachidienne conduit-elle un type de message particulier ?}

\begin{attention}[Bien poser le problème]
Ne confonds jamais le \cle{nerf rachidien} (tronc commun, en aval de la jonction des deux racines) avec les \cle{racines} elles-mêmes (en amont de la jonction, au contact de la moelle). Toute la finesse des expériences de Bell et Magendie consiste à sectionner les racines \emph{séparément}, avant leur réunion, pour analyser leur rôle respectif.
\end{attention}

\courssub{Hypothèses de travail}

Face à ce problème, deux hypothèses s'opposent logiquement :
\begin{itemize}
\item \textbf{Hypothèse 1} : les deux racines conduisent le même type de message ; leur séparation n'est qu'un détail anatomique sans signification fonctionnelle.
\item \textbf{Hypothèse 2} : chaque racine conduit un type de message particulier — l'une conduirait les messages \emph{sensitifs} (de la périphérie vers la moelle, messages \cle{afférents}), l'autre les messages \emph{moteurs} (de la moelle vers les muscles, messages \cle{efférents}).
\end{itemize}

Pour trancher, la démarche expérimentale classique s'impose : sectionner sélectivement une racine, observer les troubles qui en résultent, et comparer avec l'animal témoin. C'est exactement ce qu'ont réalisé Charles Bell puis François Magendie.

\courssub{Les expériences de Bell (1811) et Magendie (1822)}

\courssub{Protocole expérimental}

\begin{experience}[Sections sélectives des racines rachidiennes]
\textbf{Matériel} : chiens ou grenouilles anesthésiés, instruments de dissection, électrodes de stimulation, appareil d'enregistrement des contractions musculaires (myographe).

\textbf{Protocole} : après mise à nu des racines rachidiennes d'un même métamère (par exemple au niveau lombaire), on réalise chez plusieurs lots d'animaux :
\begin{enumerate}
\item \textbf{Lot A} : section de la seule racine dorsale d'un côté, l'autre côté servant de témoin ;
\item \textbf{Lot B} : section de la seule racine ventrale d'un côté ;
\item \textbf{Lot C} : section des deux racines du même côté ;
\item \textbf{Lot D} : après section d'une racine, stimulation électrique séparée de chacun des deux bouts (bout central, relié à la moelle ; bout périphérique, relié au nerf).
\end{enumerate}
On teste ensuite la sensibilité (réaction à la piqûre) et la motricité (capacité de mouvement) du territoire correspondant, en comparant systématiquement le côté opéré au côté sain.
\end{experience}

\courssub{Résultats expérimentaux}

Les résultats obtenus sont remarquablement nets et reproductibles :

\begin{itemize}
\item \textbf{Section de la racine dorsale} : l'animal ne réagit plus aux piqûres du territoire correspondant (\cle{anesthésie} : perte de sensibilité), mais il bouge normalement ce territoire (motricité conservée).
\item \textbf{Section de la racine ventrale} : l'animal ne peut plus mouvoir le territoire correspondant (\cle{paralysie} motrice), mais il réagit encore aux piqûres (sensibilité conservée).
\item \textbf{Section des deux racines} : le territoire correspondant est à la fois insensible et paralysé : anesthésie \emph{et} paralysie totales.
\item \textbf{Stimulations après section de la racine ventrale} : la stimulation du \emph{bout périphérique} (relié au muscle) provoque la contraction des muscles du territoire ; la stimulation du \emph{bout central} (relié à la moelle) ne provoque aucune réponse motrice.
\item \textbf{Stimulations après section de la racine dorsale} : la stimulation du \emph{bout central} provoque des manifestations de douleur (l'animal crie, se débat) ; la stimulation du \emph{bout périphérique} ne provoque aucune sensation ni contraction directe.
\end{itemize}

\begin{popfigure}
\begin{center}
\small
\begin{tabularx}{\linewidth}{|l|X|X|}
\hline
\rowcolor{popblueL} \textbf{Expérience} & \textbf{Trouble observé} & \textbf{Interprétation} \\
\hline
Section de la racine dorsale & Perte de sensibilité (anesthésie) ; motricité conservée & La racine dorsale conduit les messages \textbf{sensitifs} vers la moelle \\
\hline
Section de la racine ventrale & Paralysie motrice ; sensibilité conservée & La racine ventrale conduit les messages \textbf{moteurs} vers les muscles \\
\hline
Section des deux racines & Anesthésie et paralysie totales du territoire & Le nerf rachidien transporte les deux types de messages \\
\hline
Stimulation du bout périphérique de la racine ventrale sectionnée & Contraction des muscles & Le message moteur circule de la moelle vers le muscle (efférent) \\
\hline
Stimulation du bout central de la racine ventrale sectionnée & Aucune réponse motrice & La racine ventrale ne reçoit pas de message sensitif \\
\hline
Stimulation du bout central de la racine dorsale sectionnée & Manifestations de douleur & Le message sensitif circule de la périphérie vers la moelle (afférent) \\
\hline
\end{tabularx}
\end{center}
\legende{Tableau récapitulatif des expériences de Bell (1811) et Magendie (1822) sur les racines rachidiennes. La comparaison systématique du côté opéré au côté témoin permet d'attribuer à chaque racine une fonction précise.}
\end{popfigure}

\courssub{Interprétation et conclusion}

L'interprétation repose sur un raisonnement rigoureux. Quand on sectionne la racine dorsale, seule la sensibilité disparaît : c'est donc que cette racine est la \emph{voie obligatoire} des messages sensitifs. Comme la stimulation de son bout central (côté moelle) provoque une douleur, le message sensitif circule bien \textbf{de la périphérie vers la moelle} : il est \cle{afférent}. Symétriquement, la section de la racine ventrale abolit seule la motricité, et la stimulation de son bout périphérique (côté muscle) provoque la contraction : le message moteur circule \textbf{de la moelle vers les effecteurs} : il est \cle{efférent}.

\begin{propriete}[Loi de Bell-Magendie]
La \cle{racine dorsale} (postérieure) des nerfs rachidiens est \textbf{sensitive} : elle conduit les messages nerveux afférents, de la périphérie (récepteurs) vers la moelle épinière. La \cle{racine ventrale} (antérieure) est \textbf{motrice} : elle conduit les messages nerveux efférents, de la moelle vers les effecteurs (muscles). Le \cle{nerf rachidien}, formé par la réunion des deux racines, est donc un nerf \textbf{mixte} : il contient à la fois des fibres sensitives et des fibres motrices. Cette propriété, établie expérimentalement, est connue sous le nom de \textbf{loi de Bell-Magendie} ; son interprétation à partir d'expériences de section est exigible au Baccalauréat.
\end{propriete}

\begin{attention}[Le ganglion spinal : un repère infaillible]
Erreur très fréquente au Bac : attribuer le \cle{ganglion spinal} à la racine ventrale. Retiens définitivement que le ganglion spinal (renflement contenant les corps cellulaires des neurones sensitifs en T) se trouve \textbf{toujours sur la racine dorsale}, c'est-à-dire sur la racine \emph{sensitive}. Sur un schéma ou une coupe, repérer le ganglion permet d'identifier immédiatement la racine dorsale, donc le sens de circulation des messages.
\end{attention}

\courssub{Les expériences de dégénérescence wallérienne}

\courssub{Principe de la dégénérescence}

Les expériences de section donnent les fonctions des racines, mais comment suivre le \emph{trajet exact} des fibres nerveuses et confirmer le \emph{sens} de circulation des messages ? C'est ce que permet la \cle{dégénérescence wallérienne}, décrite par Augustus Waller (1850).

Le principe repose sur une propriété fondamentale du neurone : le \cle{corps cellulaire} (qui contient le noyau) est le centre trophique de la cellule, c'est-à-dire le siège des synthèses indispensables à la survie du prolongement. Toute portion de fibre nerveuse \emph{séparée de son corps cellulaire} par une section ne peut plus être entretenue : elle \textbf{dégénère} et meurt en quelques jours. En revanche, la portion restée \emph{en continuité avec le corps cellulaire} survit.

\courssub{Application aux racines rachidiennes}

Appliquons ce principe aux deux racines :

\begin{itemize}
\item \textbf{Section de la racine dorsale} : les corps cellulaires des neurones sensitifs se trouvent dans le \emph{ganglion spinal}, sur la racine dorsale elle-même. Après section de la racine entre le ganglion et la moelle, la portion de fibre située \emph{entre la section et la moelle} (côté central) est séparée de son corps cellulaire : elle dégénère. La portion périphérique, restée reliée au ganglion, survit. La dégénérescence se dirige donc \emph{vers la moelle} : elle confirme que le message sensitif circule de la périphérie vers le centre (afférent).
\item \textbf{Section de la racine ventrale} : les corps cellulaires des motoneurones se trouvent dans la \emph{substance grise de la moelle} (cornes ventrales). Après section, c'est la portion \emph{périphérique} de la fibre, séparée de la moelle, qui dégénère. La dégénérescence se dirige donc \emph{vers le muscle} : elle confirme que le message moteur circule du centre vers la périphérie (efférent).
\end{itemize}

\begin{popfigure}
\begin{center}
\begin{tikzpicture}[scale=0.92, every node/.style={font=\small}]
% ==== Racine dorsale ====
\node[font=\bfseries, popblue] at (3,4.6) {Racine dorsale (sensorielle)};
% moelle
\draw[fill=popblueL, draw=popblue] (0,2.6) rectangle (1.6,4.2);
\node[align=center] at (0.8,3.4) {Moelle};
% ganglion
\draw[fill=poporangeL, draw=poporange] (3.6,3.4) ellipse (0.55 and 0.5);
\node[align=center, font=\footnotesize] at (3.6,3.4) {Ganglion\\spinal};
% récepteur
\draw[fill=popgreenL, draw=popgreen] (6.6,3.4) circle (0.4);
\node[font=\footnotesize] at (6.6,3.4) {Réc.};
% fibre survivante (ganglion -> récepteur)
\draw[very thick, popgreen] (4.15,3.4) -- (6.2,3.4);
% fibre dégénérée (moelle -> ganglion) en pointillés
\draw[very thick, poppink, dashed, decorate, decoration={zigzag,segment length=4,amplitude=1.2}] (1.6,3.4) -- (3.05,3.4);
% point de section
\draw[very thick, popdark] (2.9,3.05) -- (2.9,3.75);
\node[font=\footnotesize, popdark] at (2.9,2.85) {section};
% légendes fibres
\node[font=\footnotesize, poppink, align=center] at (2.2,4.05) {dégénère\\(séparée du corps cellulaire)};
\node[font=\footnotesize, popgreen] at (5.3,3.75) {survit};
% ==== Racine ventrale ====
\node[font=\bfseries, poppurple] at (3,0.9) {Racine ventrale (motrice)};
\draw[fill=popblueL, draw=popblue] (0,-1.1) rectangle (1.6,0.5);
\node[align=center] at (0.8,-0.3) {Moelle};
% corps cellulaire moteur
\draw[fill=poppurpleL, draw=poppurple] (0.8,-0.3) circle (0.16);
% muscle
\draw[fill=popgoldL, draw=popgold] (6.2,-0.75) rectangle (7.2,0.15);
\node[font=\footnotesize] at (6.7,-0.3) {Muscle};
% fibre survivante (moelle -> section)
\draw[very thick, popgreen] (1.6,-0.3) -- (2.9,-0.3);
% fibre dégénérée (section -> muscle)
\draw[very thick, poppink, dashed, decorate, decoration={zigzag,segment length=4,amplitude=1.2}] (2.9,-0.3) -- (6.2,-0.3);
\draw[very thick, popdark] (2.9,-0.65) -- (2.9,0.05);
\node[font=\footnotesize, popdark] at (2.9,-0.85) {section};
\node[font=\footnotesize, popgreen] at (2.2,0.05) {survit};
\node[font=\footnotesize, poppink, align=center] at (4.6,0.1) {dégénère\\(séparée du corps cellulaire)};
\end{tikzpicture}
\end{center}
\legende{Dégénérescence wallérienne après section d'une racine rachidienne. En haut : racine dorsale ; le corps cellulaire étant dans le ganglion spinal, c'est la portion centrale (vers la moelle) qui dégénère. En bas : racine ventrale ; le corps cellulaire du motoneurone étant dans la moelle, c'est la portion périphérique (vers le muscle) qui dégénère. Le sens de la dégénérescence révèle le sens de circulation du message nerveux.}
\end{popfigure}

\begin{aretenir}[Règle pratique de la dégénérescence wallérienne]
Après section d'une fibre nerveuse, \textbf{la portion qui dégénère est toujours celle qui est séparée du corps cellulaire}. Pour une racine dorsale, la dégénérescence gagne le côté de la moelle ; pour une racine ventrale, elle gagne le côté du muscle. Suivre la dégénérescence, c'est donc suivre le trajet et le sens des fibres nerveuses.
\end{aretenir}

\begin{savaistu}[Bell, Magendie et une belle querelle de priorité]
L'Écossais Charles Bell publie ses résultats en 1811, dans un opuscule confidentiel, en montrant le rôle moteur de la racine ventrale. Le Français François Magendie, en 1822, démontre de façon décisive, chez le chien, le rôle sensitif de la racine dorsale — provoquant une vive controverse sur la priorité de la découverte. L'histoire des sciences a tranché avec élégance : la loi porte aujourd'hui le nom de \textbf{loi de Bell-Magendie}, honorant les deux savants qui ont travaillé indépendamment. C'est un bel exemple de la construction collective du savoir scientifique.
\end{savaistu}

\courssub{Méthode d'interprétation et application clinique}

\begin{methode}[Interpréter une expérience de section de racine]
Face à un document décrivant une section ou une stimulation de racine rachidienne, procède toujours ainsi :
\begin{enumerate}
\item \textbf{Identifie la structure lésée ou stimulée} : racine dorsale (repère : ganglion spinal), racine ventrale, ou nerf rachidien (tronc mixte).
\item \textbf{Compare le côté opéré au côté témoin} : seule la comparaison permet d'attribuer le trouble à la lésion et non à l'anesthésie ou au choc opératoire.
\item \textbf{Distingue les troubles sensitifs des troubles moteurs} : perte de sensation (anesthésie, absence de réaction à la piqûre) = atteinte de la voie afférente ; paralysie (absence de mouvement) = atteinte de la voie efférente.
\item \textbf{Conclus sur la fonction de la structure} : la fonction perdue après section est la fonction assurée par la structure sectionnée.
\end{enumerate}
\end{methode}

\begin{exemplebox}[Paralysie sans anesthésie]
Chez un patient dont les racines ventrales lombaires ont été accidentellement sectionnées (traumatisme du rachis), les jambes sont paralysées : il ne peut ni se lever ni marcher. Pourtant, il ressent parfaitement les piqûres, le chaud et le froid sur la peau des jambes. Ce tableau s'explique parfaitement par la loi de Bell-Magendie : les racines ventrales (motrices) étant coupées, les messages moteurs n'atteignent plus les muscles des membres inférieurs ; mais les racines dorsales (sensitives) étant intactes, les messages sensitifs remontent normalement vers la moelle puis vers l'encéphale. À l'inverse, une atteinte isolée des racines dorsales (comme dans certaines complications du diabète ou du zona) provoque une anesthésie avec motricité conservée.
\end{exemplebox}

\begin{exoresolu}[Localiser une lésion nerveuse à partir des troubles observés]
\textbf{Énoncé.} À la suite d'une fracture vertébrale, un patient présente au niveau du membre inférieur droit une perte totale de sensibilité cutanée et une paralysie complète des muscles. Le membre inférieur gauche est indemne. (1) Précise, en justifiant, la nature et la localisation de la lésion. (2) Que se passerait-il si l'on stimulait électriquement le bout périphérique des fibres motrices sectionnées ? Justifie.
\tcblower
\textbf{Solution détaillée.}

(1) Le patient présente à la fois une \textbf{anesthésie} (perte de sensibilité) et une \textbf{paralysie} (perte de motricité) du même territoire. D'après la loi de Bell-Magendie, la sensibilité transite par la racine dorsale et la motricité par la racine ventrale. La perte simultanée des deux fonctions signifie donc que \textbf{les deux racines} du côté droit sont lésées, ou que le \textbf{nerf rachidien} (tronc mixte formé par leur réunion) est sectionné du côté droit. Le caractère unilatéral des troubles confirme que la lésion est située \emph{à droite}, en aval ou au niveau de la jonction des racines, et non dans la moelle elle-même (qui aurait pu provoquer des troubles bilatéraux).

(2) La stimulation du bout périphérique des fibres motrices sectionnées provoque la \textbf{contraction des muscles} du membre. En effet, ce bout est resté en continuité avec les effecteurs : l'influx artificiel provoqué par l'électrode se propage directement jusqu'aux fibres musculaires, sans passer par la moelle. C'est exactement le résultat obtenu par Magendie lors de la stimulation du bout périphérique de la racine ventrale sectionnée. (On notera qu'à long terme, cette portion de fibre finirait par dégénérer par dégénérescence wallérienne, car elle est séparée du corps cellulaire du motoneurone resté dans la moelle.)
\end{exoresolu}

\begin{exercice}[À toi de jouer]
Une grenouille subit la section de la racine dorsale droite au niveau lombaire. On pique ensuite la peau de la patte arrière droite, puis celle de la patte arrière gauche. (1) Prévois les réponses observées de chaque côté, en les justifiant. (2) Après quelques jours, quelle portion de fibre aura dégénéré ? Justifie par la localisation du corps cellulaire du neurone sensitif.
\end{exercice}

\begin{corrige}[Exercice : section de la racine dorsale chez la grenouille]
(1) Du côté droit (opéré), la piqûre ne provoque \textbf{aucune réaction} : la racine dorsale étant sectionnée, le message sensitif ne peut plus atteindre la moelle ; il y a anesthésie du territoire. Du côté gauche (témoin), la piqûre provoque un \textbf{mouvement de retrait} de la patte : la voie sensitive et la voie motrice sont intactes. (2) Le corps cellulaire du neurone sensitif se trouve dans le \textbf{ganglion spinal}, sur la racine dorsale, en dehors de la moelle. La portion de fibre située \emph{entre le point de section et la moelle} (bout central) est séparée de ce corps cellulaire : c'est elle qui dégénère par dégénérescence wallérienne. La portion périphérique, restée reliée au ganglion, survit.
\end{corrige}

\begin{aretenir}[L'essentiel de la section]
\begin{itemize}
\item La \textbf{racine dorsale} (avec ganglion spinal) est \textbf{sensitive} : message afférent, de la périphérie vers la moelle.
\item La \textbf{racine ventrale} (sans ganglion) est \textbf{motrice} : message efférent, de la moelle vers les muscles.
\item Le \textbf{nerf rachidien}, réunion des deux racines, est \textbf{mixte}.
\item La \textbf{dégénérescence wallérienne} (dégénérescence de la portion de fibre séparée du corps cellulaire) permet de suivre le trajet et le sens des fibres.
\item Application clinique : anesthésie seule = lésion dorsale ; paralysie seule = lésion ventrale ; anesthésie + paralysie = lésion des deux racines ou du nerf mixte.
\end{itemize}
\end{aretenir}

Désormais, tu sais qu'un message sensitif pénètre dans la moelle par la racine dorsale et qu'un message moteur en sort par la racine ventrale. Mais que devient le message \emph{à l'intérieur} de la moelle ? Comment une simple piqûre déclenche-t-elle, en une fraction de seconde et sans intervention de la volonté, le retrait du membre ? C'est toute la question des \emph{réflexes médullaires}, que les expériences de Pflüger chez la grenouille nous permettront d'élucider dans la section suivante.

\coursec{Les réflexes médullaires : expériences de Pflüger et arc réflexe}

Dans la section précédente, les expériences de Bell et Magendie nous ont permis de comprendre le rôle des racines rachidiennes : la racine dorsale conduit les messages \cle{sensitifs} vers la moelle épinière, tandis que la racine ventrale conduit les messages \cle{moteurs} vers les effecteurs (muscles ou glandes). La moelle épinière apparaît donc comme une voie de passage entre la périphérie du corps et l'encéphale. Mais est-elle seulement un simple « câble » de transmission ? Une question fondamentale se pose alors : une grenouille privée de son encéphale peut-elle encore réagir de manière coordonnée à une excitation ? Autrement dit, la moelle épinière possède-t-elle une activité propre, capable d'élaborer une réponse motrice ? C'est à cette question que le physiologiste allemand Eduard Pflüger a répondu au XIX\textsuperscript{e} siècle par une série d'expériences rigoureuses, devenues classiques, que nous allons analyser selon la démarche d'investigation scientifique.

\courssub{Problème et hypothèse}

\textbf{Observation.} Une grenouille décapitée (dont l'encéphale a été détruit) suspendue par la mâchoire retire brusquement sa patte postérieure lorsqu'on la plonge dans une solution d'acide dilué. Cet animal, pourtant privé de cerveau, réagit donc encore de manière coordonnée à une excitation nociceptive.

\textbf{Problème.} Quelle structure nerveuse est responsable de cette réaction motrice qui persiste en l'absence de l'encéphale ?

\textbf{Hypothèse.} La moelle épinière, restée intacte, pourrait être le centre nerveux responsable de cette réponse.

\textbf{Vérification expérimentale.} Pflüger a réalisé une série de manipulations chirurgicales sur la grenouille, en détruisant ou en sectionnant successivement chaque structure susceptible d'intervenir, puis en observant la présence ou l'absence de la réponse. C'est la méthode d'ablation : on supprime un élément, on observe ce qui disparaît, et l'on en déduit le rôle de cet élément.

\courssub{Les expériences de Pflüger}

\begin{experience}[Expérience 1 : le réflexe persiste sans encéphale]
\textbf{Protocole.} On détruit l'encéphale d'une grenouille (décérébration) en introduisant une aiguille à ébrécher dans la boîte crânienne ; la moelle épinière est laissée intacte. Après disparition des troubles de l'opération, on plonge la patte postérieure de l'animal dans une solution d'acide dilué (par exemple de l'acide acétique à 0,5 \%).

\textbf{Résultat.} La grenouille retire vivement sa patte : la réponse motrice persiste.

\textbf{Interprétation.} L'encéphale n'est pas indispensable à la réalisation de cette réponse. Le centre nerveux de cette réaction se situe donc en dehors de l'encéphale, vraisemblablement dans la moelle épinière.
\end{experience}

\begin{experience}[Expérience 2 : la destruction de la moelle abolit le réflexe]
\textbf{Protocole.} Sur la même grenouille décérébrée, on détruit maintenant la moelle épinière (démédullation), puis on plonge à nouveau la patte dans l'acide dilué.

\textbf{Résultat.} Aucune réponse motrice n'est observée : la patte ne se retire plus, quelle que soit l'intensité de l'excitation.

\textbf{Interprétation.} La moelle épinière est indispensable à la réalisation de la réponse : c'est elle le \cle{centre nerveux} de ce réflexe. L'hypothèse est validée.
\end{experience}

\begin{methode}[Technique de démédullation de la grenouille]
Pour déméduller une grenouille au laboratoire :
\begin{enumerate}
\item Décérébrer d'abord l'animal en détruisant l'encéphale avec une aiguille à ébrécher introduite dans la boîte crânienne ;
\item Introduire l'aiguille à ébrécher dans le canal rachidien par le trou occipital (foramen magnum) ;
\item Détruire la moelle épinière par des mouvements de rotation et de va-et-vient de l'aiguille sur toute la longueur du canal rachidien ;
\item Avant toute expérience, vérifier l'absence totale de réflexe (la patte excitée ne doit plus bouger) : c'est le témoin qui garantit que la démédullation est complète.
\end{enumerate}
\end{methode}

\begin{experience}[Expériences 3 et 4 : le récepteur et le nerf sont indispensables]
\textbf{Expérience 3.} Sur une grenouille décérébrée à moelle intacte, on détruit la peau de la patte postérieure (zone réceptrice) par une brûlure ou par abrasion, puis on plonge la patte dans l'acide dilué : \textbf{aucune réponse}. En revanche, si l'on excite directement le nerf sciatique de cette même patte (par un courant électrique), le muscle se contracte et la patte se retire : \textbf{réponse présente}.

\textbf{Interprétation.} La peau contient les \cle{récepteurs} (nocicepteurs) qui captent l'excitation ; ils sont indispensables pour déclencher le réflexe de manière naturelle. Le nerf et le muscle, eux, sont encore fonctionnels puisque l'excitation directe du nerf provoque la réponse.

\textbf{Expérience 4.} Sur une grenouille décérébrée à moelle intacte, on sectionne le nerf sciatique du côté droit, puis on plonge les deux pattes postérieures dans l'acide dilué : la patte droite (nerf sectionné) \textbf{ne répond pas}, la patte gauche (nerf intact) \textbf{répond normalement}.

\textbf{Interprétation.} Le nerf est la \cle{voie de conduction} indispensable : il transmet le message sensitif à la moelle (via la racine dorsale) et le message moteur au muscle (via la racine ventrale). Si la voie est coupée, le circuit est interrompu et le réflexe disparaît.
\end{experience}

\begin{center}
\begin{tabularx}{\linewidth}{|X|X|X|}
\hline
\rowcolor{popblueL} \textbf{Manipulation} & \textbf{Résultat} & \textbf{Conclusion} \\
\hline
Destruction de l'encéphale, moelle intacte ; patte plongée dans l'acide dilué & La patte se retire & L'encéphale n'est pas indispensable au réflexe \\
\hline
Destruction de la moelle épinière ; même excitation & Aucune réponse & La moelle épinière est le centre nerveux du réflexe \\
\hline
Destruction de la peau de la patte ; même excitation & Aucune réponse & Le récepteur cutané est indispensable \\
\hline
Excitation directe du nerf sciatique (peau détruite) & La patte se retire & Le nerf et le muscle sont fonctionnels \\
\hline
Section du nerf sciatique droit ; excitation des deux pattes & Patte droite : pas de réponse ; patte gauche : réponse normale & Le nerf est la voie de conduction du message \\
\hline
\end{tabularx}
\end{center}

\begin{propriete}[Centre nerveux du réflexe médullaire]
Le centre nerveux d'un réflexe médullaire est la \cle{moelle épinière}, et plus précisément sa \cle{substance grise} (cornes dorsale, intermédiaire et ventrale). L'encéphale n'est pas indispensable à sa réalisation : le circuit nerveux du réflexe médullaire est strictement médullaire. Cette propriété, établie expérimentalement par Pflüger, doit pouvoir être déduite de l'analyse des résultats expérimentaux au Baccalauréat.
\end{propriete}

\courssub{Les lois de Pflüger}

En faisant varier l'intensité de l'excitation appliquée à un membre de la grenouille décérébrée, Pflüger a observé que l'amplitude et l'étendue de la réponse dépendent de la force du stimulus. Il en a dégagé quatre lois, dites \cle{lois de Pflüger} :

\begin{enumerate}
\item \textbf{Loi de l'excitation unilatérale} : si l'excitation est faible, seul le membre excité réagit (réponse du seul côté excité) ;
\item \textbf{Loi de la symétrie} : si l'excitation est moyenne, le membre excité et le membre homolatéral controlatéral (symétrique, du côté opposé) réagissent ;
\item \textbf{Loi de l'irradiation} : si l'excitation est forte, la réponse s'étend aux membres des autres segments (par exemple les pattes antérieures en plus des pattes postérieures) ;
\item \textbf{Loi de la généralisation} : si l'excitation est très forte (intense), tout le corps de l'animal réagit (contraction généralisée).
\end{enumerate}

Ces lois montrent que l'influx nerveux, arrivé dans la moelle épinière, peut se propager à des motoneurones de segments voisins d'autant plus facilement que l'excitation est intense : on parle d'\emph{irradiation} de l'influx dans la substance grise, via les interneurones et les faisceaux de la substance blanche (faisceaux propriospinaux).

\begin{popfigure}
\begin{center}
\begin{tikzpicture}[scale=0.9, every node/.style={font=\small}]
% Corps schématique de la grenouille vue dorsale, 4 niveaux d'intensité
\foreach \i/\titre/\actifs in {0/{Excitation\\faible}/{1}, 1/{Excitation\\moyenne}/{1,2}, 2/{Excitation\\forte}/{1,2,3,4}, 3/{Excitation\\très forte}/{1,2,3,4,5}}{
\begin{scope}[xshift=\i*4.2cm]
% corps
\draw[popline, thick] (0,0) ellipse (0.55 and 1.15);
% pattes : 1 = post. droite (excitée), 2 = post. gauche, 3 = ant. droite, 4 = ant. gauche, 5 = corps entier
\draw[popline, thick] (0.55,-0.7) -- (1.15,-1.25); % patte post droite
\draw[popline, thick] (-0.55,-0.7) -- (-1.15,-1.25); % patte post gauche
\draw[popline, thick] (0.55,0.7) -- (1.15,1.25); % patte ant droite
\draw[popline, thick] (-0.55,0.7) -- (-1.15,1.25); % patte ant gauche
% stimulus
\node[poporange, font=\small\bfseries] at (1.55,-1.45) {$\star$};
% coloration des membres actifs
\ifnum\i=0
  \draw[poporange, very thick] (0.55,-0.7) -- (1.15,-1.25);
\fi
\ifnum\i=1
  \draw[poporange, very thick] (0.55,-0.7) -- (1.15,-1.25);
  \draw[poporange, very thick] (-0.55,-0.7) -- (-1.15,-1.25);
\fi
\ifnum\i=2
  \draw[poporange, very thick] (0.55,-0.7) -- (1.15,-1.25);
  \draw[poporange, very thick] (-0.55,-0.7) -- (-1.15,-1.25);
  \draw[poporange, very thick] (0.55,0.7) -- (1.15,1.25);
  \draw[poporange, very thick] (-0.55,0.7) -- (-1.15,1.25);
\fi
\ifnum\i=3
  \draw[poporange, very thick] (0.55,-0.7) -- (1.15,-1.25);
  \draw[poporange, very thick] (-0.55,-0.7) -- (-1.15,-1.25);
  \draw[poporange, very thick] (0.55,0.7) -- (1.15,1.25);
  \draw[poporange, very thick] (-0.55,0.7) -- (-1.15,1.25);
  \draw[poporange, very thick, fill=poporange!20] (0,0) ellipse (0.55 and 1.15);
\fi
\node[align=center, font=\footnotesize] at (0,-2.1) {\titre};
\end{scope}
}
\node[font=\footnotesize, align=center] at (0,2.0) {Loi de l'excitation\\unilatérale};
\node[font=\footnotesize, align=center] at (4.2,2.0) {Loi de la\\symétrie};
\node[font=\footnotesize, align=center] at (8.4,2.0) {Loi de\\l'irradiation};
\node[font=\footnotesize, align=center] at (12.6,2.0) {Loi de la\\généralisation};
\end{tikzpicture}
\end{center}
\legende{Illustration des quatre lois de Pflüger chez la grenouille décérébrée. L'étoile orange ($\star$) indique le point d'application du stimulus sur la patte postérieure droite ; les segments en orange épais sont ceux qui réagissent (contraction des fléchisseurs). Plus l'excitation est intense, plus la réponse s'étend : membre excité seul, puis membre controlatéral symétrique, puis autres segments (membres antérieurs), puis tout le corps.}
\end{popfigure}

\courssub{Le réflexe médullaire et l'arc réflexe}

\begin{definition}[Réflexe et arc réflexe]
Un \cle{réflexe} est une réponse motrice \textbf{involontaire}, \textbf{innée}, \textbf{rapide} et \textbf{toujours identique} (stéréotypée) à une même excitation. On parle de \cle{réflexe médullaire} lorsque le centre nerveux de ce réflexe est la moelle épinière.

L'\cle{arc réflexe} est l'ensemble des structures anatomiques qui interviennent dans la réalisation d'un réflexe, du récepteur à l'effecteur.
\end{definition}

L'analyse des expériences de Pflüger, complétée par celles de Bell et Magendie (section précédente), permet de dégager les cinq éléments de l'arc réflexe médullaire :

\begin{enumerate}
\item Le \cle{récepteur} : organe sensoriel (ici les nocicepteurs cutanés de la peau) qui capte l'excitation (douleur, chaleur) et la transforme en message nerveux sensitif (potentiel d'action) ;
\item La \cle{voie sensitive} : le neurone sensitif (ou fibre sensitive) dont le corps cellulaire est situé dans le \cle{ganglion spinal} (ganglion de la racine dorsale) et dont l'axone pénètre dans la moelle par la \cle{racine dorsale} ; il conduit le message de la périphérie vers le centre nerveux ;
\item Le \cle{centre nerveux} : la substance grise de la moelle épinière, où un ou plusieurs \cle{interneurones} (neurones d'association ou d'intercalation) élaborent le message moteur à partir du message sensitif reçu ; le corps cellulaire de l'interneurone est dans la substance grise ;
\item La \cle{voie motrice} : le \cle{motoneurone} (neurone moteur, fibre motrice) dont le corps cellulaire est dans la corne ventrale de la substance grise et dont l'axone quitte la moelle par la \cle{racine ventrale} ; il conduit le message moteur du centre vers la périphérie ;
\item L'\cle{effecteur} : le muscle (muscle squelettique pour les réflexes somatiques) ou la glande qui réalise la réponse, ici la contraction des muscles fléchisseurs qui retirent la patte.
\end{enumerate}

\begin{popfigure}
\begin{center}
\begin{tikzpicture}[scale=1.0, font=\small]
% ===== Moelle en coupe transversale =====
% Substance blanche (fond)
\draw[fill=popblueL, draw=popline, thick] (0,0) ellipse (2.0 and 1.7);
% Substance grise (papillon)
\draw[fill=popblue!55, draw=popdark, thick]
 (0,0.95) .. controls (0.5,0.75) and (0.45,0.3) .. (0.35,0.1)
 .. controls (0.5,-0.3) and (0.75,-0.7) .. (0.55,-1.05)
 .. controls (0.35,-1.25) and (0.1,-1.0) .. (0.08,-0.6)
 .. controls (-0.08,-0.6) and (-0.35,-1.25) .. (-0.55,-1.05)
 .. controls (-0.75,-0.7) and (-0.5,-0.3) .. (-0.35,0.1)
 .. controls (-0.45,0.3) and (-0.5,0.75) .. (0,0.95) -- cycle;
\node[font=\footnotesize, popdark] at (0,-0.15) {Substance grise};
\node[font=\footnotesize, popmuted] at (1.35,1.15) {Substance blanche};
% Canal épendymaire
\draw[fill=white, draw=popdark] (0,0.15) circle (0.06);
% ===== Racine dorsale + ganglion =====
\draw[popgreen, very thick] (-6.2,3.4) .. controls (-4.5,3.2) and (-3.2,2.6) .. (-1.6,1.5);
\draw[fill=popgreen!60, draw=popgreen, thick] (-3.1,2.55) circle (0.3);
\node[font=\footnotesize, popgreen!60!black] at (-3.1,3.15) {Ganglion spinal};
\node[font=\footnotesize, popgreen!60!black, align=center] at (-1.9,2.35) {Racine dorsale\\(sensitive)};
% ===== Récepteur =====
\draw[fill=poporangeL, draw=poporange, thick] (-6.9,3.2) rectangle (-5.9,3.7);
\node[font=\footnotesize, align=center] at (-6.4,4.0) {\textbf{1. Récepteur}\\(nocicepteur cutané)};
% ===== Neurone sensitif (dans racine dorsale) =====
\node[font=\footnotesize, popgreen!60!black, align=center] at (-5.0,2.6) {\textbf{2. Neurone}\\\textbf{sensitif}};
% ===== Interneurone =====
\draw[poppurple, very thick] (-0.9,0.75) .. controls (-0.4,0.55) and (0.4,0.55) .. (0.9,0.75);
\draw[fill=poppurple!50, draw=poppurple] (-0.9,0.75) circle (0.14);
\draw[fill=poppurple!50, draw=poppurple] (0.9,0.75) circle (0.14);
\node[font=\footnotesize, poppurple!70!black] at (0,1.35) {\textbf{3. Interneurone(s)}};
% ===== Motoneurone / racine ventrale =====
\draw[popink, very thick] (0.9,0.75) .. controls (1.4,0.2) and (1.8,-0.6) .. (2.6,-1.1);
\draw[popink, very thick] (2.6,-1.1) .. controls (4.0,-1.7) and (5.2,-1.9) .. (6.2,-1.7);
\node[font=\footnotesize, popink, align=center] at (2.6,-2.0) {Racine ventrale\\(motrice)};
\node[font=\footnotesize, popink] at (0.9,0.35) {\textbf{4. Motoneurone}};
% ===== Muscle effecteur =====
\draw[fill=poppinkL, draw=poppink, thick, rounded corners=3pt] (6.2,-2.2) rectangle (7.9,-1.2);
\node[font=\footnotesize, align=center] at (7.05,-0.85) {\textbf{5. Effecteur}\\(muscle fléchisseur)};
% Flèches de sens de l'influx
\draw[-{Stealth[length=2.5mm]}, popgreen!60!black, thick] (-5.6,3.35) -- (-5.0,3.25);
\draw[-{Stealth[length=2.5mm]}, popink, thick] (4.6,-1.85) -- (5.3,-1.8);
\end{tikzpicture}
\end{center}
\legende{Schéma fonctionnel de l'arc réflexe médullaire (coupe transversale de la moelle). Le message nerveux suit le trajet : \textbf{1} récepteur cutané (nocicepteur) $\rightarrow$ \textbf{2} neurone sensitif (corps cellulaire dans le ganglion spinal, axone dans la racine dorsale) $\rightarrow$ \textbf{3} interneurone(s) dans la substance grise de la moelle $\rightarrow$ \textbf{4} motoneurone (corps cellulaire dans la corne ventrale, axone dans la racine ventrale) $\rightarrow$ \textbf{5} muscle effecteur (fléchisseur). Ce circuit ne passe jamais par l'encéphale.}
\end{popfigure}

\begin{attention}[Erreurs fréquentes à ne pas commettre]
\begin{itemize}
\item \textbf{Confondre le centre nerveux :} Écrire que « le message passe par le cerveau » ou que « le cerveau commande le réflexe ». \textbf{Faux} pour le réflexe médullaire : les expériences de Pflüger montrent que le circuit est strictement médullaire, puisque la réponse persiste après destruction de l'encéphale. Le cerveau peut être \emph{informé ensuite} (perception consciente de la douleur), mais il ne commande pas le réflexe.
\item \textbf{Inverser les racines :} Confondre racine dorsale (sensitive, avec ganglion) et racine ventrale (motrice, sans ganglion). Rappel : \textbf{Dorsale = Détecte (sensitive), Ventrale = Va vers le muscle (motrice)}.
\item \textbf{Oublier l'interneurone :} Dire que le neurone sensitif fait synapse directement sur le motoneurone. Dans le réflexe de retrait (polysynaptique), il y a au moins un interneurone excitateur (et souvent un interneurone inhibiteur pour le muscle antagoniste : inhibition réciproque).
\item \textbf{Confondre réflexe et réaction volontaire :} Le réflexe est inné, involontaire, rapide (30--50 ms) et stéréotypé. Le mouvement volontaire est acquis, conscient, plus lent (150--300 ms) et adaptable.
\end{itemize}
\end{attention}

\begin{exemplebox}[Un réflexe médullaire de la vie quotidienne]
Lorsque tu retires brusquement ta main d'une marmite chaude posée sur le feu de bois, ou du bord d'une flamme, tu réalises un réflexe médullaire de protection (réflexe de retrait ou réflexe de flexion) : les nocicepteurs et thermorécepteurs cutanés sont excités, le message est traité par la moelle épinière (niveau cervical pour la main), et les muscles fléchisseurs du bras se contractent \emph{avant même} que tu aies conscience de la brûlure. Ce réflexe protège les tissus d'une destruction plus grave.
\end{exemplebox}

\begin{savaistu}[La rapidité remarquable des réflexes et application clinique]
Le temps de réponse d'un réflexe médullaire est de l'ordre de quelques dizaines de millisecondes (environ 30 à 50 ms), car le circuit nerveux est court : il ne comporte que deux ou trois neurones (arc polysynaptique) et ne passe pas par le cerveau. Un mouvement volontaire, qui exige l'intervention du cortex cérébral, demande au contraire 150 à 300 ms. Cette rapidité explique la valeur protectrice des réflexes : dans les services d'urgence des hôpitaux de Yaoundé, Douala ou Garoua, le médecin teste systématiquement les réflexes ostéotendineux (rotulien L2-L4, achilléen S1-S2, bicipital C5-C6) pour vérifier l'intégrité de la moelle épinière et des racines rachidiennes après un accident de la route, malheureusement fréquent sur nos axes routiers (ex. route Yaoundé-Douala). L'absence de réflexe (aréflexie) signe une lésion de l'arc réflexe (nerf, racine, moelle) ; l'hyperréflexie (réflexes vifs, clonus) signe une lésion de la voie pyramidale (voie motrice volontaire) au-dessus du segment testé.
\end{savaistu}

\begin{exoresolu}[Interprétation d'expériences sur le réflexe de la grenouille]
Énoncé. On réalise sur trois grenouilles A, B et C les opérations suivantes : A est décérébrée (moelle intacte) ; B est décérébrée puis démédullée ; C est décérébrée et on lui sectionne le nerf sciatique gauche. On plonge ensuite les pattes postérieures de chaque animal dans une solution d'acide acétique dilué. Les résultats sont : A retire ses deux pattes ; B ne réagit pas ; C retire seulement la patte droite. (1) Attribue à chaque structure (encéphale, moelle, nerf) son rôle dans le réflexe. (2) Déduis-en la définition du réflexe médullaire.
\tcblower
Solution détaillée.

(1) \textbf{Comparaison de A et B} : la seule différence entre A et B est la présence ou l'absence de la moelle épinière ; or A réagit et B ne réagit pas. La moelle épinière est donc indispensable : c'est le \cle{centre nerveux} du réflexe. De plus, A réagit bien que dépourvue d'encéphale : l'encéphale n'est donc pas indispensable à ce réflexe.

\textbf{Comparaison des deux pattes de C} : la patte gauche, dont le nerf sciatique est sectionné, ne répond pas, alors que la patte droite, dont le nerf est intact, répond. Le nerf est donc la \cle{voie de conduction} indispensable du message nerveux (sensitif à l'aller vers la moelle, moteur au retour vers le muscle).

(2) Un \cle{réflexe médullaire} est une réponse motrice involontaire, innée, rapide et stéréotypée à une excitation, dont le centre nerveux est la moelle épinière. Il met en jeu un arc réflexe comprenant : récepteur, voie sensitive (neurone sensitif + racine dorsale), centre nerveux médullaire (substance grise + interneurone(s)), voie motrice (motoneurone + racine ventrale) et effecteur.
\end{exoresolu}

\begin{aretenir}[L'essentiel de la section]
\begin{itemize}
\item Une grenouille décérébrée réagit encore à une excitation : le réflexe médullaire ne nécessite pas l'encéphale ;
\item La destruction de la moelle épinière abolit le réflexe : la moelle (substance grise) est le \textbf{centre nerveux} du réflexe médullaire ;
\item Le récepteur cutané et le nerf sont indispensables : ils assurent respectivement la \textbf{réception} (transduction) et la \textbf{conduction} du message nerveux ;
\item Les lois de Pflüger décrivent l'extension de la réponse avec l'intensité du stimulus : excitation unilatérale, symétrie, irradiation, généralisation (phénomène d'irradiation dans la substance grise) ;
\item L'arc réflexe comprend cinq éléments : récepteur $\rightarrow$ voie sensitive (neurone sensitif, racine dorsale, ganglion spinal) $\rightarrow$ centre nerveux (substance grise, interneurone) $\rightarrow$ voie motrice (motoneurone, racine ventrale) $\rightarrow$ effecteur (muscle).
\end{itemize}
\end{aretenir}

\begin{exercice}[Pour t'entraîner]
Une grenouille décérébrée à moelle intacte est excitée par une goutte d'acide déposée sur la patte postérieure droite. On fait varier l'intensité de l'excitation : avec une excitation faible, seule la patte droite se retire ; avec une excitation moyenne, les deux pattes postérieures se retirent ; avec une excitation forte, les quatre membres se retirent.
\begin{enumerate}
\item Nomme les lois de Pflüger mises en évidence par chacun de ces trois résultats.
\item Prévois le résultat obtenu avec une excitation très intense.
\item Explique, à l'aide de l'arc réflexe, pourquoi la destruction de la peau de la patte droite supprime toute réponse à l'acide déposé sur cette patte.
\end{enumerate}
\end{exercice}

\begin{corrige}[Exercice : lois de Pflüger]
\begin{enumerate}
\item Excitation faible, réponse du seul côté excité : loi de l'\textbf{excitation unilatérale}. Excitation moyenne, réponse des deux membres symétriques (controlatéraux) : loi de la \textbf{symétrie}. Excitation forte, extension aux autres segments (membres antérieurs) : loi de l'\textbf{irradiation}.
\item Avec une excitation très intense, tout le corps de la grenouille réagit (contraction généralisée) : c'est la loi de la \textbf{généralisation}.
\item La peau contient les récepteurs (nocicepteurs) qui captent l'excitation et la transforment en message nerveux (transduction). Si la peau est détruite, le premier élément de l'arc réflexe (le récepteur) est absent : aucun message sensitif n'est généré, donc la moelle ne reçoit aucune information et aucune réponse motrice n'est déclenchée, même si le nerf, la moelle et le muscle sont intacts.
\end{enumerate}
\end{corrige}

\coursec{Le réflexe myotatique : réflexe rotulien et achilléen}

Dans la section précédente, nous avons vu que les réflexes médullaires étudiés par Pflüger chez la grenouille mettent en jeu un \cle{arc réflexe} complet : récepteur cutané, voie sensitive, centre nerveux (la moelle épinière), voie motrice et effecteur. Nous avons aussi constaté que, dans la moelle, le message sensitif est relayé par un ou plusieurs \cle{interneurones} avant d'atteindre le motoneurone : l'arc est dit \emph{polysynaptique}.

Existe-t-il des réflexes plus simples encore, plus rapides ? Oui. Lors d'une consultation de routine au Centre Hospitalier d'Essos à Yaoundé, le médecin frappe doucement le tendon situé sous votre genou avec un petit marteau : votre jambe se détend brusquement vers l'avant, sans que vous l'ayez voulu. Ce geste anodin explore un réflexe d'un type particulier, le \cle{réflexe myotatique}, que nous allons maintenant étudier en détail.

\courssub{Définition et observation}

\begin{definition}[Le réflexe myotatique]
Le \cle{réflexe myotatique} (ou \cle{réflexe d'étirement}) est la \textbf{contraction d'un muscle squelettique en réponse à son propre étirement}. Le récepteur et l'effecteur se trouvent donc \emph{dans le même muscle} : c'est ce qui distingue fondamentalement ce réflexe des réflexes médullaires cutanés étudiés chez la grenouille, où le récepteur (la peau) et l'effecteur (un muscle) sont des organes différents.
\end{definition}

L'observation de départ est simple : il suffit d'étirer brusquement un muscle squelettique pour qu'il se contracte aussitôt, de façon involontaire. Mais comment étirer un muscle \emph{in vivo}, sans dissection ? L'astuce consiste à frapper son \textbf{tendon} : le tendon transmet le choc au muscle, qui s'allonge d'un coup pendant une fraction de seconde. C'est exactement ce que réalise la percussion du tendon rotulien ou du tendon d'Achille.

\courssub{Mise en évidence expérimentale}

\textbf{Le réflexe rotulien.}\par

\begin{experience}[Mise en évidence du réflexe rotulien]
\textbf{Matériel :} un sujet volontaire, un tabouret ou une table haute, un marteau à réflexe (ou le bord de la main).

\textbf{Protocole :}
\begin{enumerate}
\item Le sujet s'assoit sur un tabouret suffisamment haut pour que sa jambe pende librement, sans toucher le sol. Il doit être parfaitement détendu (les muscles relâchés).
\item L'expérimentateur repère le tendon rotulien, cordon fibreux tendu sous la rotule, entre celle-ci et le tibia.
\item Il frappe le tendon d'un coup sec et bref, perpendiculairement.
\end{enumerate}

\textbf{Observation :} immédiatement après la percussion, la jambe effectue une \textbf{extension brusque} (le pied se projette vers l'avant), traduisant la contraction du muscle \cle{quadriceps} (muscle de la face antérieure de la cuisse). La réponse est involontaire : le sujet ne peut ni la déclencher volontairement ni l'empêcher s'il est détendu.

\textbf{Interprétation :} la percussion étire brusquement le quadriceps ; celui-ci répond en se contractant. Il s'agit bien d'une contraction en réponse à un étirement : un réflexe myotatique.
\end{experience}

\begin{methode}[Conditions de réussite du test]
\begin{enumerate}
\item Le sujet doit être \textbf{détendu} : s'il contracte volontairement ses muscles, la réponse est masquée ou exagérée.
\item La jambe doit \textbf{pendre librement} : le pied ne doit toucher ni le sol ni un obstacle.
\item Il faut frapper le \textbf{tendon} et non le muscle lui-même : seul le tendon transmet un étirement bref et net au muscle.
\item Le coup doit être \textbf{sec et bref} : une pression lente n'étire pas assez rapidement le muscle pour déclencher le réflexe.
\end{enumerate}
\end{methode}

\begin{popfigure}
\begin{center}
\begin{tikzpicture}[scale=0.9, every node/.style={font=\small}]
% Tabouret
\draw[fill=popmuted!40] (0,0) rectangle (3.2,0.35);
\draw[thick] (0.3,0) -- (0.3,-1.6);
\draw[thick] (2.9,0) -- (2.9,-1.6);
\node[below, popmuted] at (1.6,-1.6) {tabouret};
% Cuisse horizontale
\draw[thick, popdark, fill=poporangeL] (0.2,0.35) .. controls (1.5,1.15) and (2.6,1.15) .. (3.1,0.75) -- (3.1,0.35) -- cycle;
\node[popdark] at (1.6,0.85) {cuisse (quadriceps)};
% Rotule
\draw[thick, popdark, fill=popgoldL] (3.05,0.55) ellipse (0.22 and 0.3);
\node[right, popdark] at (3.35,0.55) {rotule};
% Jambe pendante
\draw[thick, popdark, fill=poporangeL] (2.85,0.3) -- (3.25,0.3) -- (3.15,-2.2) -- (2.95,-2.2) -- cycle;
% Pied
\draw[thick, popdark, fill=poporangeL] (2.95,-2.2) -- (3.15,-2.2) -- (3.9,-2.5) -- (3.0,-2.55) -- cycle;
\node[below, popmuted] at (3.3,-2.55) {jambe pendant librement};
% Tendon rotulien
\draw[very thick, poppink] (3.12,0.28) -- (3.1,-0.35);
\node[right=0.1cm, poppink, align=left] at (3.2,-0.35) {tendon rotulien\\ (point de percussion)};
% Marteau
\draw[thick, popblue] (4.9,-1.3) -- (3.9,-0.75);
\draw[thick, popblue, fill=popblueL] (3.75,-0.85) rectangle (4.05,-0.45);
\draw[->, very thick, popblue] (4.0,-0.65) -- (3.25,-0.4);
% Flèche extension
\draw[->, very thick, popgreen] (3.6,-2.3) .. controls (4.6,-2.2) and (5.2,-1.6) .. (5.5,-0.9);
\node[right, popgreen, align=left] at (5.5,-1.4) {extension\\ brusque};
\end{tikzpicture}
\end{center}
\legende{Mise en évidence du réflexe rotulien : le sujet, assis et détendu, laisse pendre sa jambe ; la percussion du tendon rotulien (flèche bleue) provoque l'extension brusque de la jambe (flèche verte).}
\end{popfigure}

\textbf{Le réflexe achilléen.}\par

\begin{experience}[Mise en évidence du réflexe achilléen]
\textbf{Protocole :} le sujet est agenouillé sur une chaise, pieds pendants dans le vide (ou couché, jambe fléchie). L'expérimentateur frappe d'un coup sec le \cle{tendon d'Achille}, le gros tendon situé à l'arrière de la cheville, au-dessus du talon.

\textbf{Observation :} le pied effectue une \textbf{extension} (la pointe du pied s'abaisse, comme lorsqu'on se met sur la pointe des pieds), traduisant la contraction des muscles du mollet (muscle \emph{triceps sural}).

\textbf{Interprétation :} la percussion étire les muscles du mollet, qui répondent en se contractant : c'est un autre exemple de réflexe myotatique.
\end{experience}

\courssub{Les éléments intervenant dans le réflexe myotatique}

Où se trouve le récepteur ? Ce n'est pas la peau, comme dans les réflexes de Pflüger : la percussion à travers la peau anesthésiée déclenche encore le réflexe, alors que la destruction du muscle ou la section de son nerf l'abolit. Le récepteur est donc \emph{dans le muscle lui-même}. L'anatomie microscopique l'a identifié : c'est le \cle{fuseau neuromusculaire}.

\begin{definition}[Le fuseau neuromusculaire]
Le \cle{fuseau neuromusculaire} est un \textbf{récepteur sensoriel situé à l'intérieur du muscle squelettique}, disposé en parallèle des fibres musculaires ordinaires (fibres extrafusales). Il est \textbf{sensible à l'étirement du muscle} : plus le muscle s'allonge, plus le fuseau est excité et plus il envoie de messages nerveux vers la moelle épinière.
\end{definition}

L'ensemble des éléments de l'arc myotatique est donc le suivant :
\begin{enumerate}
\item le \cle{fuseau neuromusculaire} : récepteur d'étirement, situé dans le muscle (quadriceps pour le réflexe rotulien) ;
\item le \cle{neurone sensitif} : son corps cellulaire est dans le ganglion de la racine postérieure ; sa fibre périphérique naît du fuseau, chemine dans le nerf mixte puis dans la \textbf{racine postérieure} d'un nerf spinal et pénètre dans la moelle épinière ;
\item la \cle{moelle épinière} : centre nerveux, où s'effectue le relais synaptique unique ;
\item le \cle{motoneurone} (ou neurone moteur) : son corps cellulaire est dans la substance grise de la corne antérieure de la moelle et son axone sort par la \textbf{racine antérieure} pour rejoindre le muscle ;
\item le \cle{muscle effecteur} : le même muscle qui a été étiré (le quadriceps), qui se contracte.
\end{enumerate}

\begin{popfigure}
\begin{center}
\begin{tikzpicture}[scale=1.0, every node/.style={font=\small}]
% Muscle
\draw[thick, popdark, fill=poporangeL] (-5.5,0.4) ellipse (1.7 and 0.9);
\node[popdark] at (-5.5,1.6) {muscle (quadriceps)};
% Fuseau neuromusculaire
\draw[thick, poppink, fill=poppinkL] (-5.5,0.4) ellipse (0.85 and 0.32);
\node[poppink] at (-5.5,-0.15) {fuseau neuromusculaire};
% Moelle épinière (coupe)
\draw[thick, popdark, fill=popblueL] (2.5,0.2) ellipse (1.5 and 1.2);
\draw[thick, poppurple, fill=poppurpleL] (2.5,0.2)
  .. controls (2.1,0.9) and (2.9,0.9) .. (2.5,0.35)
  .. controls (2.2,0.1) and (2.2,-0.5) .. (2.5,-0.75)
  .. controls (2.8,-0.5) and (2.8,0.1) .. (2.5,0.35) -- cycle;
\node[poppurple] at (2.5,1.75) {moelle épinière};
\node[poppurple, font=\footnotesize] at (2.5,-1.15) {substance grise (corne antérieure)};
% Voie sensitive (racine postérieure, en haut)
\draw[->, very thick, popblue] (-4.6,0.55) .. controls (-2.5,1.4) and (0.2,1.5) .. (1.35,0.95);
\node[popblue, above] at (-1.5,1.45) {neurone sensitif (racine postérieure)};
% Synapse unique
\draw[fill=popgold] (1.75,0.75) circle (0.09);
\node[popgold, right, align=left, font=\footnotesize] at (1.9,1.05) {synapse unique\\ (pas d'interneurone)};
% Motoneurone
\draw[fill=popgreen] (2.55,-0.35) circle (0.14);
\node[popgreen, right, font=\footnotesize] at (2.75,-0.55) {motoneurone};
% Voie motrice (racine antérieure, en bas)
\draw[->, very thick, popgreen] (2.55,-0.5) .. controls (1.5,-1.3) and (-2.5,-1.4) .. (-4.4,-0.35);
\node[popgreen, below] at (-1.5,-1.5) {motoneurone (racine antérieure)};
% Contraction
\draw[->, very thick, poporange] (-6.6,0.4) -- (-7.4,0.4);
\node[left, poporange] at (-7.4,0.4) {contraction};
% Étirement
\draw[->, very thick, poppink] (-5.5,-1.6) -- (-5.5,-2.3);
\node[below, poppink] at (-5.5,-2.3) {étirement (percussion du tendon)};
\end{tikzpicture}
\end{center}
\legende{Schéma de l'arc réflexe myotatique : l'étirement du muscle excite le fuseau neuromusculaire ; le message sensitif gagne la moelle par la racine postérieure, y fait \emph{une seule} synapse avec le motoneurone, qui commande la contraction du même muscle.}
\end{popfigure}

\courssub{Un arc monosynaptique : la particularité du réflexe myotatique}

\begin{propriete}[Monosynapticité de l'arc myotatique]
Le réflexe myotatique est le \textbf{seul réflexe monosynaptique connu} : dans la moelle épinière, le neurone sensitif fait \emph{directement} synapse avec le motoneurone, \textbf{sans interneurone intercalé}. C'est le réflexe \textbf{le plus rapide} de l'organisme.
\end{propriete}

Comment a-t-on établi cette particularité ? Par le raisonnement suivant, fondé sur la \textbf{rapidité exceptionnelle} de la réponse. On sait que le message nerveux se propage le long des fibres myélinisées à une vitesse élevée (de l'ordre de \SI{50}{m.s^{-1}} à \SI{100}{m.s^{-1}}), mais que chaque \textbf{relais synaptique} introduit un retard (temps de libération du neuromédiateur et d'excitation du neurone suivant). Or le temps de réaction du réflexe rotulien est extrêmement bref, de l'ordre de \num{30} à \SI{50}{ms}, valeur à peine supérieure au temps de trajet du message dans les nerfs augmenté d'\emph{un seul} délai synaptique. On en déduit que l'arc ne comporte qu'\textbf{une seule synapse centrale} : le neurone sensitif touche directement le motoneurone. (Ce savoir est admis au niveau Terminale ; il est déduit de la rapidité de la réponse et confirmé par l'histologie.)

\begin{popfigure}
\begin{center}
\begin{tikzpicture}[scale=0.95, every node/.style={font=\small}]
% ===== ARC POLYSYNAPTIQUE (gauche) =====
\node[font=\small\bfseries, popdark] at (-4.2,3.1) {Arc médullaire classique (polysynaptique)};
\draw[thick, popdark, fill=popcream] (-6.6,1.9) circle (0.35);
\node[below, popdark, font=\footnotesize] at (-6.6,1.5) {récepteur (peau)};
\draw[thick, popblue, fill=popblueL] (-4.9,1.9) circle (0.3);
\node[below, popblue, font=\footnotesize] at (-4.9,1.5) {neurone sensitif};
\draw[thick, poppurple, fill=poppurpleL] (-3.3,1.9) circle (0.3);
\node[below, poppurple, font=\footnotesize] at (-3.3,1.5) {interneurone};
\draw[thick, popgreen, fill=popgreenL] (-1.7,1.9) circle (0.3);
\node[below, popgreen, font=\footnotesize] at (-1.7,1.5) {motoneurone};
\draw[thick, poporange, fill=poporangeL] (-0.2,1.9) circle (0.35);
\node[below, poporange, font=\footnotesize] at (-0.2,1.5) {muscle};
\draw[->, thick] (-6.25,1.9) -- (-5.2,1.9);
\draw[->, thick] (-4.6,1.9) -- (-3.6,1.9);
\draw[->, thick] (-3.0,1.9) -- (-2.0,1.9);
\draw[->, thick] (-1.4,1.9) -- (-0.55,1.9);
% Synapses marquées
\draw[fill=popgold] (-3.45,1.9) circle (0.07);
\draw[fill=popgold] (-1.85,1.9) circle (0.07);
\node[popgold, above, font=\footnotesize] at (-2.65,2.05) {2 synapses (au moins)};
% ===== ARC MONOSYNAPTIQUE (droite) =====
\node[font=\small\bfseries, popdark] at (3.4,3.1) {Arc myotatique (monosynaptique)};
\draw[thick, poppink, fill=poppinkL] (1.0,1.9) ellipse (0.55 and 0.3);
\node[below, poppink, font=\footnotesize] at (1.0,1.5) {fuseau};
\draw[thick, popblue, fill=popblueL] (2.9,1.9) circle (0.3);
\node[below, popblue, font=\footnotesize] at (2.9,1.5) {neurone sensitif};
\draw[thick, popgreen, fill=popgreenL] (4.7,1.9) circle (0.3);
\node[below, popgreen, font=\footnotesize] at (4.7,1.5) {motoneurone};
\draw[thick, poporange, fill=poporangeL] (6.3,1.9) circle (0.35);
\node[below, poporange, font=\footnotesize] at (6.3,1.5) {muscle};
\draw[->, thick] (1.55,1.9) -- (2.6,1.9);
\draw[->, thick] (3.2,1.9) -- (4.4,1.9);
\draw[->, thick] (5.0,1.9) -- (5.95,1.9);
\draw[fill=popgold] (4.55,1.9) circle (0.07);
\node[popgold, above, font=\footnotesize] at (3.85,2.05) {1 seule synapse};
% Bandeau moelle
\draw[dashed, popmuted] (-3.9,0.9) rectangle (-1.0,2.9);
\node[popmuted, font=\footnotesize] at (-2.45,0.65) {moelle épinière};
\draw[dashed, popmuted] (3.0,0.9) rectangle (5.4,2.9);
\node[popmuted, font=\footnotesize] at (4.2,0.65) {moelle épinière};
\end{tikzpicture}
\end{center}
\legende{Comparaison des deux types d'arcs réflexes : à gauche, l'arc médullaire classique comporte au moins un interneurone intercalé (plusieurs synapses) ; à droite, l'arc myotatique relie directement le neurone sensitif au motoneurone (une seule synapse), d'où sa rapidité.}
\end{popfigure}

\begin{attention}[Erreur fréquente au Baccalauréat]
Ne dis \textbf{jamais} que le réflexe rotulien fait intervenir un interneurone ! L'arc myotatique est \textbf{monosynaptique} : le neurone sensitif fait directement synapse avec le motoneurone. C'est précisément ce qui le distingue de l'arc médullaire classique (polysynaptique) des expériences de Pflüger. Autre piège : le récepteur n'est \emph{pas} dans la peau ni dans le tendon, mais \emph{dans le muscle} (le fuseau neuromusculaire). Le tendon n'est que le point d'application du choc qui étire le muscle.
\end{attention}

\courssub{Importance physiologique du réflexe myotatique}

À quoi sert un tel réflexe, déclenché des milliers de fois par jour sans que nous en ayons conscience ?

\begin{aretenir}[Rôles du réflexe myotatique]
\begin{itemize}
\item \textbf{Maintien de la posture et de l'équilibre} : debout, le corps oscille imperceptiblement ; dès qu'un muscle postural (mollets, cuisses, dos, cou) s'étire sous l'effet d'un déséquilibre, ses fuseaux neuromusculaires déclenchent sa contraction, qui redresse le corps. Le réflexe myotatique est ainsi à la base du \cle{tonus postural}, cette contraction légère et permanente des muscles qui nous maintient debout.
\item \textbf{Régulation automatique de la longueur des muscles} : tout étirement excessif est immédiatement combattu par une contraction, ce qui protège les muscles et stabilise les articulations lors des mouvements.
\end{itemize}
\end{aretenir}

\begin{exemplebox}[Le réflexe myotatique en classe]
Tu t'assoupis en cours, assis sur ta chaise : ta tête penche lentement vers l'avant. Les muscles de la nuque s'étirent alors ; leurs fuseaux neuromusculaires sont excités et déclenchent, par réflexe myotatique, la contraction de ces muscles : ta tête se redresse brusquement — ce fameux « hochement de tête » qui te réveille en sursaut ! De même, un porteur d'eau ou de bois dans nos villes qui trébuche redresse instantanément son tronc grâce à ces réflexes d'étirement, avant même d'avoir pris conscience du déséquilibre.
\end{exemplebox}

\courssub{Importance médicale : l'examen des réflexes ostéotendineux}

\begin{savaistu}[Le marteau à réflexe, un instrument centenaire]
Les médecins utilisent le marteau à réflexe depuis la fin du XIX\textsuperscript{e} siècle (les neurologues Wilhelm Erb et Carl Westphal ont décrit le réflexe rotulien vers 1875). Ce test simple, rapide et indolore est resté, près de 150 ans plus tard, un examen neurologique de routine pratiqué dans tous les services de consultation, y compris dans les centres de santé camerounais, car il renseigne immédiatement sur l'intégrité des nerfs et de la moelle épinière.
\end{savaistu}

L'examen des réflexes ostéotendineux (rotulien, achilléen, mais aussi bicipital au coude) est un \textbf{test neurologique de routine}. Sa logique est directement issue de la structure de l'arc : si l'arc est interrompu en un point quelconque, le réflexe disparaît ; si le contrôle exercé par l'encéphale sur la moelle est supprimé, le réflexe devient exagéré.

\begin{itemize}
\item \textbf{Abolition du réflexe} (aréflexie) : elle signe une lésion de l'arc lui-même — atteinte du nerf (section, compression, neuropathie), des racines rachidiennes, ou de la moelle au niveau concerné. Par exemple, une sciatique par hernie discale comprimant la racine S1 peut abolir le réflexe achilléen ; une atteinte des motoneurones de la moelle au niveau lombaire abolit le réflexe rotulien du côté lésé.
\item \textbf{Exagération du réflexe} (hyperréflexie) : elle traduit au contraire une lésion des voies nerveuses \emph{descendantes} qui, normalement, freinent la moelle depuis l'encéphale (par exemple après un accident vasculaire cérébral ou une compression médullaire haute). La moelle, libérée de ce contrôle inhibiteur, répond de façon excessive (spasticité).
\end{itemize}

Ainsi, la simple percussion d'un tendon permet au médecin de \textbf{localiser} une lésion nerveuse : lésion de l'arc (réflexe aboli) ou lésion des centres supérieurs (réflexe exagéré).

\begin{exoresolu}[Interpréter un cas clinique]
\textbf{Énoncé.} À l'hôpital, on examine deux patients. Le patient A, victime d'une section du nerf fémoral (nerf du quadriceps) à la cuisse, ne présente \emph{aucun} réflexe rotulien du côté blessé, bien que son quadriceps soit intact. Le patient B, atteint d'une lésion de la moelle épinière au niveau thoracique, présente des réflexes rotulien et achilléen \emph{exagérés} aux deux jambes. Explique ces deux observations à partir de la structure de l'arc réflexe myotatique.

\tcblower
\textbf{Solution détaillée.}

\emph{Patient A.} L'arc myotatique du réflexe rotulien comprend : fuseau neuromusculaire du quadriceps, neurone sensitif (racine postérieure), moelle épinière (synapse unique), motoneurone (racine antérieure), quadriceps. Tous ces éléments doivent être intacts pour que le message parcoure l'arc. Or le nerf fémoral contient à la fois les fibres sensitives venant du quadriceps et les fibres motrices qui s'y rendent. Sa section \textbf{interrompt l'arc en deux points} (voie sensitive et voie motrice) : le message du fuseau ne peut plus atteindre la moelle, et l'ordre moteur ne peut plus atteindre le muscle. Le réflexe est donc aboli, même si le muscle est sain. C'est une \textbf{paralysie avec aréflexie}, signe d'une lésion de l'arc réflexe lui-même (lésion périphérique).

\emph{Patient B.} La lésion thoracique se situe \emph{au-dessus} des segments lombaires (L2-L4 pour le rotulien) et sacrés (S1-S2 pour l'achilléen) de la moelle où se font les synapses des réflexes. L'arc myotatique des membres inférieurs reste donc \textbf{intact} : les réflexes persistent. Mais les voies descendantes venant de l'encéphale (voies pyramidales et extrapyramidales), qui exercent normalement un contrôle inhibiteur (freinateur) sur les motoneurones médullaires, sont interrompues. La moelle lombaire et sacrée, « libérée » de ce frein, répond de façon excessive : les réflexes sont \textbf{exagérés} (hyperréflexie, signe de lésion centrale du neurone moteur supérieur).

\emph{Conclusion.} L'abolition d'un réflexe indique une lésion de l'arc (nerf, racines, moelle au niveau du réflexe) ; son exagération indique une lésion des voies de contrôle descendantes, au-dessus du niveau de l'arc. L'examen des réflexes ostéotendineux permet ainsi de localiser la lésion (topographie diagnostique).
\end{exoresolu}

\begin{exercice}[Pour t'entraîner]
Un sujet anesthésié localement au niveau de la peau du genou conserve un réflexe rotulien normal. En revanche, après injection d'un curare (substance qui bloque la jonction entre le nerf moteur et le muscle), le réflexe disparaît alors que le message sensitif parvient normalement à la moelle.
\begin{enumerate}
\item Que prouve la première observation quant à la localisation du récepteur du réflexe rotulien ?
\item À quel niveau le curare interrompt-il l'arc réflexe ?
\item Rappelle les cinq éléments de l'arc myotatique dans l'ordre.
\end{enumerate}
\end{exercice}

\begin{corrige}[Exercice]
\begin{enumerate}
\item L'anesthésie de la peau supprime la sensibilité cutanée sans abolir le réflexe : le récepteur n'est donc \textbf{pas dans la peau}. Il est situé \textbf{dans le muscle} : c'est le fuseau neuromusculaire, excité par l'étirement du quadriceps transmis par le tendon.
\item Le curare bloque la \textbf{jonction neuromusculaire} (synapse entre le motoneurone et la fibre musculaire), c'est-à-dire le dernier relais de l'arc, entre la voie motrice et l'effecteur. Le message sensitif et le centre nerveux sont intacts, mais l'ordre moteur ne peut plus déclencher la contraction.
\item Fuseau neuromusculaire (récepteur) $\rightarrow$ neurone sensitif (racine postérieure) $\rightarrow$ moelle épinière (synapse unique avec le motoneurone) $\rightarrow$ motoneurone (racine antérieure) $\rightarrow$ muscle effecteur (le même muscle qui a été étiré).
\end{enumerate}
\end{corrige}

\begin{aretenir}[L'essentiel de la section]
\begin{itemize}
\item Le \cle{réflexe myotatique} est la contraction d'un muscle en réponse à son propre étirement ; récepteur et effecteur sont dans le même muscle.
\item Le récepteur est le \cle{fuseau neuromusculaire}, sensible à l'étirement.
\item L'arc myotatique est \textbf{monosynaptique} (une seule synapse, sans interneurone) : c'est le réflexe le plus rapide de l'organisme, contrairement à l'arc médullaire classique, polysynaptique.
\item Il assure le \textbf{maintien de la posture} (tonus postural) et la régulation automatique de la longueur des muscles.
\item L'examen des réflexes rotulien et achilléen est un test neurologique de routine : réflexe \textbf{aboli} = lésion de l'arc (névrite, radiculaire, moelle au niveau) ; réflexe \textbf{exagéré} = lésion des voies de contrôle descendantes (hémorragie cérébrale, compression médullaire haute).
\end{itemize}
\end{aretenir}

\coursec{Les réflexes acquis ou conditionnels}

Dans la section précédente, nous avons étudié le réflexe myotatique (rotulien, achilléen) : un réflexe \cle{inné}, présent dès la naissance, dont le centre nerveux est la moelle épinière et dont le circuit anatomique, l'arc réflexe, est fixe et identique chez tous les individus. Frapper le tendon rotulien provoque toujours la même réponse, chez tout individu sain, sans aucun apprentissage. Pourtant, l'observation quotidienne montre que certains comportements ressemblent à des réflexes — ils sont rapides, automatiques, involontaires — alors qu'ils n'existent pas à la naissance : le conducteur camerounais qui freine instantanément à la vue d'un feu rouge n'est pas né avec ce comportement ; il l'a \emph{appris}. Comment l'organisme peut-il acquérir de nouvelles réponses automatiques ? Où se situe leur centre nerveux ? Peuvent-elles disparaître ? C'est à ces questions que répond l'étude des \cle{réflexes acquis ou conditionnels}, dont le siège n'est plus la moelle épinière mais le \cle{cerveau}.

\courssub{Distinction entre réflexes innés et réflexes acquis}

\begin{definition}[Réflexe inné et réflexe acquis]
Un \cle{réflexe inné} (ou congénital) est une réponse motrice ou sécrétoire, involontaire et stéréotypée, \textbf{présente dès la naissance}, identique chez tous les individus de l'espèce, et dont le centre nerveux est la moelle épinière (réflexes médullaires) ou le tronc cérébral. Un \cle{réflexe acquis} (ou conditionnel) est une réponse motrice ou sécrétoire, involontaire, \textbf{acquise par l'apprentissage} au cours de la vie individuelle ; il résulte de l'association répétée d'un stimulus indifférent à un stimulus naturel, et son centre nerveux est le \textbf{cortex cérébral} (cerveau).
\end{definition}

\begin{exemplebox}[Réflexes innés et réflexes acquis dans la vie courante]
\begin{itemize}
\item \textbf{Réflexes innés} : le réflexe rotulien ; le clignement des paupières quand un objet s'approche brusquement de l'œil ; le retrait de la main au contact d'une braise ; la salivation provoquée par la présence de nourriture dans la bouche ; la succion chez le nouveau-né.
\item \textbf{Réflexes acquis} : la salivation d'un élève à la simple vue d'un plat de ndolé bien préparé ; le freinage automatique du conducteur à la vue d'un feu rouge ; l'arrêt des passants au son strident du sifflet d'un agent de circulation à Douala.
\end{itemize}
Remarque essentielle : dans les deux cas, la réponse est \textbf{rapide et involontaire}. Ce qui les distingue, c'est l'\emph{origine} (naissance ou apprentissage) et le \emph{centre nerveux} (moelle ou cerveau).
\end{exemplebox}

\courssub{Mise en évidence expérimentale : l'expérience de Pavlov}

\textbf{Observation de départ.} Ivan Pavlov, physiologiste russe, étudiait les sécrétions digestives du chien. Il remarqua un fait troublant : ses chiens se mettaient à saliver \emph{avant même} de recevoir la nourriture, dès qu'ils entendaient les pas de l'expérimentateur ou le bruit de la porte de la salle d'alimentation. Or la salivation est un réflexe inné dont le stimulus naturel est la présence de nourriture dans la bouche. Comment un bruit, qui ne provoque normalement aucune salivation, pouvait-il déclencher cette réponse ? Pavlov formula l'hypothèse qu'une \emph{association apprise} s'était établie entre le bruit et la nourriture, et il conçut une expérience rigoureuse pour la vérifier.

\begin{experience}[L'expérience de Pavlov : le conditionnement de la salivation chez le chien]
\textbf{Matériel} : un chien placé dans un dispositif d'isolement, muni d'une fistule salivaire (petit tube greffé sur le canal d'une glande salivaire) permettant de recueillir et de mesurer la salive goutte à goutte ; une cloche (ou un métronome) ; de la nourriture (viande en poudre).\par
\textbf{Protocole et observations} :
\begin{enumerate}
\item \textbf{Étape 1 — Vérifications préalables.} On présente la nourriture seule : le chien salive abondamment (réflexe inné : la nourriture est le \cle{stimulus naturel} ou inconditionnel). On sonne la cloche seule : le chien dresse l'oreille mais \textbf{ne salive pas} (la cloche est un \cle{stimulus indifférent} pour la salivation).
\item \textbf{Étape 2 — Association.} On fait sonner la cloche, puis on présente la nourriture quelques secondes après. On répète cette association \textbf{cloche + nourriture} un grand nombre de fois (plusieurs dizaines d'essais).
\item \textbf{Étape 3 — Test.} On sonne la cloche \textbf{seule}, sans nourriture : le chien \textbf{salive}. La salivation est mesurable à la fistule.
\end{enumerate}
\textbf{Interprétation} : le stimulus indifférent (la cloche), associé de façon répétée au stimulus naturel (la nourriture), est devenu capable de déclencher à lui seul la réponse. Il est devenu un \cle{stimulus conditionnel}, et la salivation qu'il provoque est un \cle{réflexe conditionnel} (ou acquis).\par
\textbf{Conclusion} : l'apprentissage par association permet à l'organisme d'acquérir de nouvelles réponses réflexes.
\end{experience}

\begin{popfigure}
\begin{center}
\begin{tikzpicture}[font=\small, >=Stealth]
% Étape 1
\node[draw=popblue, fill=popblueL, rounded corners, text width=3.4cm, align=center] (e1) at (0,0) {\textbf{Avant l'association}\\[2pt] Cloche seule\\ $\Rightarrow$ pas de salivation\\ (stimulus indifférent)};
% Étape 2
\node[draw=poporange, fill=poporangeL, rounded corners, text width=3.4cm, align=center] (e2) at (5.2,0) {\textbf{Pendant l'association}\\[2pt] Cloche \emph{puis} nourriture\\ répétées ensemble\\ (de nombreux essais)};
% Étape 3
\node[draw=popgreen, fill=popgreenL, rounded corners, text width=3.4cm, align=center] (e3) at (10.4,0) {\textbf{Après l'association}\\[2pt] Cloche seule\\ $\Rightarrow$ salivation\\ (stimulus conditionnel)};
\draw[->, very thick, popdark] (e1) -- (e2);
\draw[->, very thick, popdark] (e2) -- (e3);
% Réflexe inné de référence
\node[draw=popmuted, dashed, rounded corners, text width=3.4cm, align=center] (r1) at (0,-2.6) {Nourriture seule\\ $\Rightarrow$ salivation\\ (réflexe inné)};
\draw[->, dashed, popmuted] (r1.east) to[out=0,in=-90] node[midway, below, text=popmuted, font=\scriptsize]{renforcement} (e2.south);
\end{tikzpicture}
\end{center}
\legende{Les trois étapes du conditionnement de Pavlov. Avant l'apprentissage, la cloche est indifférente pour la salivation ; l'association répétée \og cloche puis nourriture \fg{} la transforme en stimulus conditionnel capable de déclencher seul la salivation. Le réflexe inné (nourriture $\rightarrow$ salivation) sert de renforcement à chaque essai.}
\end{popfigure}

\begin{savaistu}[Ivan Pavlov (1849--1936), prix Nobel de physiologie]
Ivan Petrovitch Pavlov, physiologiste russe, reçut le \textbf{prix Nobel de physiologie ou médecine en 1904} pour ses travaux sur la physiologie de la digestion (sécrétions salivaires et gastriques). C'est au cours de ces recherches qu'il découvrit, presque par hasard, le réflexe conditionnel, qu'il appela \og réflexe conditionné \fg. Ses travaux ont fondé la science de l'apprentissage et influencé toute la psychologie du \textsc{xx}\textsuperscript{e} siècle. L'expression \og chien de Pavlov \fg{} est restée célèbre dans le langage courant pour désigner une réaction automatique acquise.
\end{savaistu}

\courssub{Mécanisme du réflexe conditionnel (en bref)}

Comment l'association de deux stimuli crée-t-elle une nouvelle réponse ? Pavlov lui-même, puis ses successeurs, ont proposé l'interprétation suivante, admise au programme.

Lors de l'association répétée, \textbf{deux centres nerveux du cerveau sont excités simultanément} : le centre sensitif recevant le message issu de la cloche (centre auditif) et le centre commandant la salivation, excité par la nourriture. Cette double excitation répétée établit entre ces deux centres une \cle{connexion temporaire}, une sorte de \og chemin fonctionnel \fg{} nouveau dans le cortex cérébral. Dès lors, l'excitation du centre auditif par la cloche seule suffit à rejoindre, par cette connexion, le centre de la salivation : le stimulus indifférent est devenu stimulus conditionnel.

\begin{aretenir}[L'essentiel du mécanisme]
\begin{itemize}
\item Le réflexe conditionnel repose sur l'établissement d'une \textbf{connexion temporaire} entre deux centres excités simultanément dans le \textbf{cerveau} (cortex cérébral).
\item Le stimulus indifférent devient \textbf{stimulus conditionnel} grâce à son association répétée avec le stimulus naturel.
\item Le \textbf{siège des réflexes conditionnels est le cortex cérébral}, contrairement aux réflexes innés médullaires dont le centre est la moelle épinière. C'est pourquoi un animal décervebré conserve ses réflexes innés mais perd ses réflexes acquis.
\end{itemize}
\end{aretenir}

\courssub{Conditions d'établissement et extinction du réflexe conditionnel}

L'expérimentation montre que le conditionnement n'est efficace que si certaines conditions sont respectées :

\begin{propriete}[Conditions d'établissement d'un réflexe conditionnel]
\begin{enumerate}
\item \textbf{Répétition de l'association} : le stimulus indifférent et le stimulus naturel doivent être présentés ensemble un grand nombre de fois ; une seule association ne suffit généralement pas.
\item \textbf{Ordre temporel} : le stimulus indifférent doit \textbf{précéder légèrement} le stimulus naturel (quelques secondes). Si la nourriture est donnée avant la cloche, le conditionnement est très difficile, voire impossible.
\item \textbf{Renforcement} : à chaque essai, le stimulus conditionnel doit être \og renforcé \fg{} par la présentation effective du stimulus naturel.
\end{enumerate}
\end{propriete}

\begin{propriete}[Extinction du réflexe conditionnel]
Un réflexe conditionnel peut \textbf{s'éteindre} : si l'on présente le stimulus conditionnel seul (la cloche sans nourriture) de nombreuses fois de suite, la réponse (salivation) devient de plus en plus faible, puis disparaît. Le réflexe conditionnel est donc un acquis \textbf{réversible} et \textbf{temporaire}, qui doit être entretenu par le renforcement. Cette propriété le distingue fondamentalement du réflexe inné, qui est permanent.
\end{propriete}

\begin{exemplebox}[Un réflexe conditionnel à l'école]
Un élève d'un lycée de Yaoundé entend chaque jour la cloche de la récréation, immédiatement suivie de la sortie en cour et de la détente. Après quelques semaines, dès que la cloche sonne, il ressent \emph{automatiquement} l'envie de sortir et une sensation de soulagement, avant même tout mouvement volontaire. La cloche (stimulus indifférent devenu conditionnel) a été associée de façon répétée à la sortie (situation naturellement agréable). Si, pendant les vacances, la cloche sonne sans jamais être suivie de récréation, cette réaction s'atténue progressivement : c'est l'\textbf{extinction}.
\end{exemplebox}

\courssub{Comparaison synthétique : réflexe inné et réflexe acquis}

\begin{popfigure}
\begin{center}
\begin{tikzpicture}[font=\small]
\node[draw=popblue, fill=popblue, text=white, rounded corners, minimum width=3.2cm, minimum height=0.7cm] at (0,0) {\textbf{Critère}};
\node[draw=poporange, fill=poporange, text=white, rounded corners, minimum width=5.2cm, minimum height=0.7cm] at (4.25,0) {\textbf{Réflexe inné}};
\node[draw=popgreen, fill=popgreen, text=white, rounded corners, minimum width=5.2cm, minimum height=0.7cm] at (9.75,0) {\textbf{Réflexe acquis (conditionnel)}};
\begin{scope}[every node/.style={draw=popline, fill=popcream, rounded corners, align=center}]
\node[minimum width=3.2cm, minimum height=0.75cm] at (0,-0.85) {Origine};
\node[minimum width=5.2cm, minimum height=0.75cm, text width=4.8cm] at (4.25,-0.85) {Congénital : présent à la naissance, sans apprentissage};
\node[minimum width=5.2cm, minimum height=0.75cm, text width=4.8cm] at (9.75,-0.85) {Acquis par l'apprentissage (association répétée)};
\node[minimum width=3.2cm, minimum height=0.75cm] at (0,-1.75) {Centre nerveux};
\node[minimum width=5.2cm, minimum height=0.75cm, text width=4.8cm] at (4.25,-1.75) {Moelle épinière (ou tronc cérébral)};
\node[minimum width=5.2cm, minimum height=0.75cm, text width=4.8cm] at (9.75,-1.75) {Cortex cérébral (cerveau)};
\node[minimum width=3.2cm, minimum height=0.75cm] at (0,-2.65) {Caractères};
\node[minimum width=5.2cm, minimum height=0.75cm, text width=4.8cm] at (4.25,-2.65) {Identique chez tous les individus ; permanent ; circuit fixe};
\node[minimum width=5.2cm, minimum height=0.75cm, text width=4.8cm] at (9.75,-2.65) {Variable selon les individus ; réversible (extinction possible) ; connexion temporaire};
\node[minimum width=3.2cm, minimum height=0.75cm] at (0,-3.55) {Exemples};
\node[minimum width=5.2cm, minimum height=0.75cm, text width=4.8cm] at (4.25,-3.55) {Réflexe rotulien ; retrait de la main au contact d'une braise ; salivation à la nourriture en bouche};
\node[minimum width=5.2cm, minimum height=0.75cm, text width=4.8cm] at (9.75,-3.55) {Salivation à la vue du ndolé ; freinage au feu rouge ; réactions apprises du conducteur};
\end{scope}
\end{tikzpicture}
\end{center}
\legende{Tableau comparatif du réflexe inné et du réflexe acquis : origine, centre nerveux, caractères et exemples. Le point capital à retenir pour le Baccalauréat est le siège nerveux : moelle épinière pour l'inné, cortex cérébral pour l'acquis.}
\end{popfigure}

\courssub{Importance des réflexes conditionnels}

Les réflexes conditionnels confèrent au comportement une \textbf{souplesse} que les réflexes innés, figés, ne peuvent offrir :

\begin{itemize}
\item \textbf{Adaptation à l'environnement} : l'organisme apprend à répondre par avance à des signaux annonciateurs (le bruit des pas annonce la nourriture ; le feu rouge annonce un danger). La réponse est anticipée, donc plus efficace pour la survie.
\item \textbf{Base de l'apprentissage et de l'éducation} : l'acquisition des automatismes (conduite d'un véhicule, gestes sportifs, réactions professionnelles) repose sur les mêmes principes d'association répétée et de renforcement. L'éducation routière, par exemple, vise à transformer en réflexes conditionnels des comportements de sécurité (ralentir à l'approche d'un passage piéton).
\item \textbf{Dressage des animaux} : les chiens de garde ou de berger, les animaux de trait sont dressés par association répétée d'un signal (sifflement, geste) à une récompense.
\end{itemize}

\begin{attention}[Erreur fréquente : confondre réflexe acquis et habitude motrice]
Ne confonds pas le \textbf{réflexe conditionnel} avec une simple \textbf{habitude motrice} (comme taper vite au clavier ou faire du vélo). Le réflexe conditionnel repose sur une \emph{association de deux stimuli} : un stimulus indifférent devient capable de déclencher une réponse qui appartenait à un réflexe inné. L'habitude motrice, elle, résulte de la \emph{répétition d'un geste volontaire} qui devient de plus en plus automatique, sans qu'un stimulus indifférent se substitue à un stimulus naturel. Au Baccalauréat, citer la cloche de Pavlov comme exemple d'habitude, ou le vélo comme exemple de réflexe conditionnel, est une erreur classique. Autre piège : écrire que le centre du réflexe conditionnel est la moelle épinière — c'est le \textbf{cerveau} (cortex cérébral).
\end{attention}

\begin{exoresolu}[Analyse d'une expérience de conditionnement]
Un expérimentateur place un chien en chambre insonorisée. Il allume une lampe rouge, puis, deux secondes après, dépose de la viande dans la gueule de l'animal ; il répète l'opération trente fois par jour pendant une semaine. Le huitième jour, il allume la lampe seule : le chien salive abondamment. Puis il cesse toute présentation de viande et n'allume plus que la lampe : au bout de quelques jours, la lampe ne provoque plus de salivation.\par
\textbf{Questions} : 1) Identifier le stimulus naturel, le stimulus indifférent devenu conditionnel, et la réponse. 2) Nommer le type de réflexe mis en évidence et préciser son centre nerveux. 3) Expliquer la disparition finale de la salivation.
\tcblower
\textbf{Solution détaillée.}\par
1) Le \textbf{stimulus naturel} (inconditionnel) est la viande dans la gueule : il déclenche la salivation par un réflexe inné. Le \textbf{stimulus indifférent} est la lumière rouge : avant l'expérience, elle ne provoque aucune salivation. Devenu \textbf{stimulus conditionnel} après association répétée, il déclenche seul la réponse, qui est la \textbf{salivation}.\par
2) Il s'agit d'un \textbf{réflexe acquis ou conditionnel} : réponse acquise par l'apprentissage, grâce à l'association répétée (trente essais par jour pendant une semaine) du stimulus indifférent au stimulus naturel, le stimulus conditionnel précédant toujours le stimulus naturel. Son centre nerveux est le \textbf{cortex cérébral} (cerveau), siège de la connexion temporaire établie entre le centre visuel excité par la lampe et le centre de la salivation excité par la viande.\par
3) La disparition finale de la salivation traduit l'\textbf{extinction} du réflexe conditionnel : le stimulus conditionnel (lampe) n'étant plus renforcé par le stimulus naturel (viande), la connexion temporaire s'affaiblit progressivement et la réponse s'éteint. Cela confirme que le réflexe conditionnel est un acquis \textbf{réversible}, contrairement au réflexe inné qui est permanent.
\end{exoresolu}

\begin{exercice}[Pour t'entraîner]
Un automobiliste de Bafoussam freine automatiquement dès qu'il aperçoit un feu rouge, même lorsqu'aucun véhicule ne croise sa route. 1) Ce comportement est-il un réflexe inné ou acquis ? Justifie par deux arguments. 2) Quel est son centre nerveux ? 3) Que deviendrait ce comportement si le feu, resté longtemps en panne, n'était plus jamais associé à un danger ni à une sanction ?
\end{exercice}

\begin{corrige}[Exercice : le freinage au feu rouge]
1) C'est un \textbf{réflexe acquis (conditionnel)}. Arguments : (a) il n'existe pas à la naissance — le nourrisson ne réagit pas au feu rouge — il a été \emph{appris} lors de l'apprentissage de la conduite ; (b) il résulte de l'\emph{association répétée} d'un stimulus indifférent (la lumière rouge) à des conséquences naturelles (danger, amende), association renforcée par l'éducation routière.\par
2) Son centre nerveux est le \textbf{cortex cérébral} (cerveau), siège des réflexes conditionnels, et non la moelle épinière.\par
3) Le réflexe s'\textbf{éteindrait progressivement} (extinction) : le stimulus conditionnel n'étant plus renforcé, la connexion temporaire s'affaiblirait et le conducteur finirait par ne plus freiner automatiquement devant ce feu.
\end{corrige}

\begin{aretenir}[Ce qu'il faut retenir de la section]
\begin{itemize}
\item Un réflexe acquis (conditionnel) est une réponse involontaire \textbf{acquise par l'apprentissage}, par association répétée d'un stimulus indifférent à un stimulus naturel ; le stimulus indifférent devient stimulus conditionnel.
\item L'expérience de \textbf{Pavlov} (cloche + nourriture $\rightarrow$ salivation à la cloche seule) en est la démonstration historique.
\item Le mécanisme repose sur une \textbf{connexion temporaire} entre deux centres excités simultanément ; le siège des réflexes conditionnels est le \textbf{cortex cérébral}.
\item Conditions : répétition de l'association, stimulus conditionnel \emph{précédant} le stimulus naturel, renforcement. Sans renforcement, le réflexe s'\textbf{éteint} (extinction).
\item Importance : adaptation du comportement à l'environnement, base de l'apprentissage, de l'éducation et du dressage.
\end{itemize}
\end{aretenir}

\coursec{Application : limitation et explication des dysfonctionnements}

Dans la section précédente, nous avons vu que les réflexes acquis montrent la plasticité de nos centres nerveux : le cerveau peut construire de nouvelles réponses à partir de l'arc réflexe inné. Mais que se passe-t-il lorsque l'un des maillons de cet arc est détruit par un accident, une infection ou une intoxication ? Au service de neurologie de l'Hôpital Général de Yaoundé, le médecin qui reçoit un accidenté de la route ne se contente pas de constater la paralysie : il cherche \emph{où} se situe la lésion, car c'est la localisation qui détermine le pronostic et le traitement. Cette démarche diagnostique, tu peux la reproduire toi-même en suivant l'arc réflexe de proche en proche. C'est l'objet de cette dernière section : utiliser nos connaissances sur l'arc réflexe pour \cle{expliquer} les dysfonctionnements et proposer des moyens de les \cle{limiter}.

\courssub{La démarche diagnostique : localiser la lésion sur l'arc réflexe}

\textbf{Le principe : cinq maillons, cinq hypothèses.} Un acte réflexe exige l'intégrité de cinq éléments disposés en chaîne : le \cle{récepteur} sensoriel, la \cle{voie sensitive} (fibre afférente, racine dorsale), le \cle{centre nerveux} (moelle épinière), la \cle{voie motrice} (motoneurone, racine ventrale, nerf moteur) et l'\cle{effecteur} (muscle). Si la réponse est absente ou anormale, au moins l'un de ces maillons est lésé. La démarche du médecin — et celle que l'on attend de toi au Baccalauréat — consiste à recueillir les signes (paralysie ? anesthésie ? réflexes abolis, conservés ou exagérés ?) puis à remonter logiquement vers le maillon défaillant.

Deux constatations cliniques orientent immédiatement le raisonnement :
\begin{itemize}
\item la \cle{paralysie} (impossibilité de mouvement volontaire) traduit une atteinte de la voie motrice, du centre ou de la voie descendante du cerveau ;
\item l'\cle{anesthésie} (perte de sensibilité) traduit une atteinte de la voie sensitive ou des voies ascendantes vers le cerveau.
\end{itemize}

\begin{methode}[Expliquer une paralysie : suivre l'arc réflexe de proche en proche]
\begin{enumerate}
\item \textbf{Identifier les troubles} : y a-t-il paralysie motrice ? anesthésie ? les deux ? Quel territoire est atteint ?
\item \textbf{Tester les réflexes} (rotulien, achilléen) : abolis, conservés ou exagérés ?
\item \textbf{Raisonner par dissociation} :
\begin{itemize}
\item paralysie \emph{avec} sensibilité conservée $\Rightarrow$ lésion \textbf{motrice} (motoneurone, racine ventrale, nerf moteur) ;
\item anesthésie \emph{sans} paralysie $\Rightarrow$ lésion \textbf{sensitive} (racine dorsale, nerf sensitif) ;
\item paralysie \emph{et} anesthésie d'un même territoire $\Rightarrow$ lésion du \textbf{nerf mixte} périphérique ou de la \textbf{moelle} elle-même ;
\item réflexes \textbf{abolis} $\Rightarrow$ l'arc réflexe local est rompu (lésion périphérique ou médullaire au niveau de l'arc) ;
\item réflexes \textbf{conservés ou exagérés} avec paralysie $\Rightarrow$ l'arc local est intact : la lésion est \textbf{au-dessus} (section médullaire haute ou atteinte cérébrale).
\end{itemize}
\item \textbf{Conclure} en désignant le maillon lésé et en justifiant chaque signe observé.
\end{enumerate}
\end{methode}

\begin{popfigure}
\begin{center}
\begin{tikzpicture}[font=\small, >=Stealth]
% Récepteur
\draw[fill=popgreenL, draw=popgreen] (0,0) ellipse (0.9 and 0.5);
\node at (0,0) {\textbf{Récepteur}};
\node[below=0.55cm of {(0,0)}] {(peau, fuseau)};
% Moelle (coupe)
\draw[fill=poppurpleL, draw=poppurple] (6.2,0) circle (1.15);
\draw[fill=white, draw=poppurple] (6.2,0) ++(0,0) ;
\draw[fill=white, draw=poppurple, rotate=0] (6.2,0.15) ellipse (0.45 and 0.6);
\node[above=1.2cm of {(6.2,0)}, align=center] {\textbf{Moelle épinière}\\ (centre)};
% Effecteur
\draw[fill=poporangeL, draw=poporange, rounded corners] (11.6,-0.5) rectangle (13.4,0.5);
\node at (12.5,0) {\textbf{Muscle}};
\node[below=0.7cm of {(12.5,0)}] {(effecteur)};
% Voie sensitive (arc supérieur)
\draw[->, very thick, popblue] (0.9,0.25) .. controls (3,1.6) .. (5.2,0.55);
\node[popblue, above] at (3,1.35) {voie sensitive (racine dorsale)};
% Voie motrice (arc inférieur)
\draw[->, very thick, poppink] (7.1,-0.55) .. controls (9.5,-1.7) .. (11.6,-0.25);
\node[poppink, below] at (9.4,-1.45) {voie motrice (racine ventrale)};
% Points de lésion
\foreach \x/\y/\n in {2.6/1.15/1, 6.2/0.15/2, 9.3/-1.05/3, 12.5/0/4}{
  \node[circle, fill=popink, text=white, inner sep=1.5pt, font=\scriptsize\bfseries] at (\x,\y) {\n};
}
% Cerveau et voie descendante
\draw[fill=popgoldL, draw=popgold] (6.2,3.1) ellipse (1.1 and 0.55);
\node at (6.2,3.1) {\textbf{Cerveau}};
\draw[->, thick, popgold, dashed] (6.2,2.55) -- (6.2,1.2);
\node[popgold, right, align=left] at (6.4,1.9) {\scriptsize contrôle\\ \scriptsize volontaire};
\node[circle, fill=popink, text=white, inner sep=1.5pt, font=\scriptsize\bfseries] at (6.2,1.9) {5};
\end{tikzpicture}
\end{center}
\legende{Schéma fonctionnel de l'arc réflexe avec les cinq points de lésion possibles. \textbf{1} : lésion de la voie sensitive $\rightarrow$ anesthésie + abolition du réflexe. \textbf{2} : lésion du centre médullaire $\rightarrow$ paralysie et anesthésie du territoire, réflexe aboli. \textbf{3} : lésion de la voie motrice $\rightarrow$ paralysie flasque, sensibilité conservée, réflexe aboli. \textbf{4} : atteinte du muscle $\rightarrow$ paralysie sans trouble sensitif. \textbf{5} : rupture du contrôle cérébral (section médullaire haute) $\rightarrow$ paralysie volontaire mais arc intact : réflexes conservés, voire exagérés.}
\end{popfigure}

\courssub{Cas 1 — La section de la moelle épinière : le traumatisme médullaire}

\textbf{Observation.} Un motocycliste victime d'un accident sur l'axe Yaoundé--Douala présente une fracture vertébrale au niveau dorsal. À l'examen : ses deux jambes sont paralysées (\cle{paraplégie}) et insensibles, pourtant le médecin constate que le \emph{réflexe rotulien persiste}, et même qu'il est \emph{exagéré}. Comment expliquer cette apparente contradiction ?

\textbf{Interprétation.} La moelle épinière est à la fois un \emph{centre} réflexe et une \emph{voie de passage} : les faisceaux ascendants portent les messages sensitifs vers le cerveau, les faisceaux descendants portent les ordres moteurs volontaires. Une section complète interrompt ces voies longues :
\begin{itemize}
\item les ordres volontaires du cerveau n'atteignent plus les motoneurones situés \emph{en dessous} de la lésion $\Rightarrow$ \textbf{paralysie} des territoires sous-lésionnaires ;
\item les messages sensitifs ne remontent plus vers le cerveau $\Rightarrow$ \textbf{anesthésie} de ces mêmes territoires ;
\item mais les \textbf{arcs réflexes médullaires} situés sous la lésion (récepteur, racines, substance grise, muscle) demeurent \emph{intacts} : les réflexes locaux \textbf{persistent}.
\end{itemize}

\begin{propriete}[Conséquence d'une section complète de la moelle épinière]
Après section complète de la moelle, les territoires situés \emph{au-dessous} de la lésion sont paralysés et anesthésiés, mais les réflexes dont l'arc est situé sous la lésion \textbf{persistent} (arc intact) ; ils échappent simplement au \cle{contrôle volontaire} et au contrôle inhibiteur du cerveau, ce qui explique qu'ils soient souvent \textbf{exagérés}.
\end{propriete}

\begin{attention}[Le piège classique du paraplégique]
Erreur fréquente : affirmer qu'un paraplégique « n'a plus aucun réflexe ». C'est faux ! Ses réflexes médullaires des membres inférieurs (rotulien, achilléen) sont le plus souvent \emph{conservés, voire exagérés}, car l'arc médullaire est intact et n'est plus freiné par le cerveau. Ce qui est aboli, c'est le \emph{mouvement volontaire} et la \emph{sensibilité consciente}, pas le circuit réflexe local.
\end{attention}

\courssub{Cas 2 — La poliomyélite : la destruction des motoneurones}

\textbf{Observation clinique.} Un adulte a gardé les séquelles d'une maladie contractée dans l'enfance : sa jambe gauche est paralysée, amaigrie (\cle{amyotrophie}), et le réflexe rotulien est aboli du côté atteint. En revanche, il \emph{ressent normalement} les piqûres et les caresses sur cette jambe.

\textbf{Interprétation.} Le virus de la poliomyélite (poliovirus) détruit sélectivement les \cle{motoneurones} des \emph{cornes ventrales} de la substance grise médullaire. Appliquons la méthode :
\begin{itemize}
\item voie motrice détruite $\Rightarrow$ paralysie \textbf{flasque} (muscle relâché) et \textbf{abolition des réflexes} (l'arc est rompu au niveau du centre) ;
\item voies sensitives (racines dorsales) intactes $\Rightarrow$ \textbf{sensibilité conservée} ;
\item le muscle, privé durablement de l'influx nerveux qui assure son trophisme, dégénère $\Rightarrow$ \textbf{amyotrophie}.
\end{itemize}

\begin{exemplebox}[Le tableau typique de la poliomyélite]
Un patient poliomyélitique présente une jambe paralysée et amaigrie car les motoneurones détruits ne stimulent plus le muscle ; mais il ressent normalement les piqûres, car les fibres sensitives et les racines dorsales sont épargnées. Cette \emph{dissociation motrice-sensitive} (paralysie + sensibilité normale + réflexes abolis) est la signature d'une lésion de la \textbf{corne ventrale}.
\end{exemplebox}

\begin{savaistu}[La poliomyélite vaincue au Cameroun]
Grâce aux campagnes de vaccination répétées du Programme Élargi de Vaccination (PEV), le Cameroun a interrompu la circulation du poliovirus sauvage et la région africaine de l'OMS a été certifiée exempte de poliomyélite sauvage en 2020. Le vaccin antipoliomyélitique (VPO, puis VPI) et le vaccin antitétanique figurent parmi les vaccins obligatoires du calendrier vaccinal camerounais : quelques gouttes ou une piqûre suffisent à prévenir une paralysie définitive, car les motoneurones détruits \emph{ne se régénèrent pas}.
\end{savaistu}

\courssub{Cas 3 — Le tétanos : quand l'inhibition médullaire est bloquée}

\textbf{Observation.} Un cultivateur de la région de l'Ouest s'est blessé au pied avec une houe rouillée. Deux semaines plus tard : mâchoires crispées (trismus), dos cambré, muscles des membres en contraction permanente et douloureuse, réflexes violemment exagérés. Pourtant, aucune paralysie, aucune anesthésie.

\textbf{Interprétation.} La bactérie \emph{Clostridium tetani}, présente dans la terre et les poussières, sécrète la \cle{toxine tétanique} (tétanospasmine). Cette toxine remonte le long des axones moteurs jusqu'à la moelle épinière, où elle \textbf{bloque les interneurones inhibiteurs} (notamment ceux qui libèrent la glycine). Or, dans le fonctionnement normal, ces interneurones freinent les motoneurones et assurent la relaxation des muscles antagonistes. Conséquences :
\begin{itemize}
\item les motoneurones, libérés de tout frein, deviennent \textbf{hyperexcitables} ;
\item le moindre stimulus déclenche des décharges massives $\Rightarrow$ \textbf{contractures permanentes} (spasmes) et \textbf{réflexes exagérés} ;
\item la sensibilité et la motricité volontaire ne sont pas détruites : il n'y a ni paralysie ni anesthésie.
\end{itemize}

\begin{attention}[Tétanos $\neq$ poliomyélite]
Ne confonds pas ces deux maladies de la moelle : la \textbf{poliomyélite} détruit les motoneurones $\rightarrow$ paralysie flasque, réflexes \emph{abolis} ; le \textbf{tétanos} bloque l'inhibition $\rightarrow$ contractures, réflexes \emph{exagérés}. Dans les deux cas, la sensibilité est conservée et la prévention repose sur la \emph{vaccination}.
\end{attention}

\courssub{Cas 4 et 5 — Atteinte d'un nerf périphérique et diagnostic par le réflexe rotulien}

\textbf{Cas 4 : la section ou compression d'un nerf.} Un nerf périphérique (par exemple le nerf sciatique, comprimé par une hernie discale, ou sectionné par une plaie profonde de la cuisse) contient des fibres \emph{motrices} et des fibres \emph{sensitives} mélangées. Sa lésion entraîne donc, dans le territoire qu'il dessert : \textbf{paralysie} (fibres motrices coupées), \textbf{anesthésie} (fibres sensitives coupées) et \textbf{abolition des réflexes} dont l'arc emprunte ce nerf. C'est le seul cas où paralysie et anesthésie coexistent \emph{avec} abolition des réflexes, dans un territoire limité à la zone d'un nerf.

\textbf{Cas 5 : le réflexe rotulien comme test diagnostique.} Le coup de marteau sur le tendon rotulien est un examen de routine car il explore l'arc myotatique le plus simple du corps (fuseau neuromusculaire du quadriceps $\rightarrow$ nerf fémoral $\rightarrow$ moelle lombaire L2--L4 $\rightarrow$ nerf fémoral $\rightarrow$ quadriceps) :
\begin{itemize}
\item réflexe \textbf{aboli} $\Rightarrow$ lésion de l'arc lui-même : nerf fémoral, racines lombaires ou moelle lombaire (ex. poliomyélite, compression discale) ;
\item réflexe \textbf{exagéré} $\Rightarrow$ arc intact mais contrôle inhibiteur cérébral supprimé : lésion des \emph{voies descendantes} (section médullaire haute, accident vasculaire cérébral, tétanos).
\end{itemize}

\begin{exoresolu}[L'accidenté de la moto]
Énoncé. Un jeune homme victime d'un accident de moto présente une fracture de la colonne vertébrale au niveau de la 10$^{\text{ème}}$ vertèbre dorsale. On observe : paralysie totale des deux membres inférieurs, absence de sensibilité à la piqûre sur les jambes, mais réflexes rotuliens présents et même exagérés. (1) Localiser la lésion. (2) Expliquer chacun des trois signes.
\tcblower
Solution. (1) \textbf{Localisation} : la coexistence d'une paralysie et d'une anesthésie des territoires situés sous le niveau de la fracture, avec des réflexes locaux \emph{conservés}, indique une \textbf{section complète de la moelle épinière} au niveau dorsal (T10). L'arc réflexe lombaire, situé sous la lésion, est intact.
(2) \textbf{Explication des signes} :
\begin{itemize}
\item \emph{Paralysie des jambes} : les faisceaux descendants (voies pyramidales) qui transportent les ordres moteurs volontaires du cerveau sont interrompus ; l'influx n'atteint plus les motoneurones lombaires.
\item \emph{Anesthésie} : les faisceaux ascendants sensitifs sont également sectionnés ; les messages des récepteurs cutanés des jambes ne parviennent plus au cerveau, siège de la sensibilité consciente.
\item \emph{Réflexes rotuliens exagérés} : l'arc myotatique lombaire (fuseau $\rightarrow$ nerf fémoral $\rightarrow$ moelle lombaire $\rightarrow$ quadriceps) est anatomiquement intact, donc fonctionnel ; il est même exagéré car il n'est plus freiné par l'action inhibitrice des voies descendantes du cerveau.
\end{itemize}
Conclusion : ce patient est paraplégique ; ses réflexes médullaires persistent mais échappent au contrôle volontaire.
\end{exoresolu}

\courssub{Limiter les dysfonctionnements : la démarche de prévention}

La connaissance des mécanismes lésionnels guide directement les mesures de \cle{limitation} des dysfonctionnements, à trois niveaux :
\begin{itemize}
\item \textbf{Prévention des causes} : la \emph{vaccination} protège contre la poliomyélite (vaccin antipoliomyélitique) et le tétanos (anatoxine tétanique, rappels réguliers) ; la \emph{sécurité routière} (port du casque à moto, ceinture de sécurité, limitation de vitesse) réduit les traumatismes médullaires, première cause de paraplégie chez les jeunes au Cameroun ; l'\emph{hygiène des plaies} (nettoyage, désinfection, consultation précoce) empêche le développement de \emph{C. tetani}.
\item \textbf{Soins précoces} : l'immobilisation correcte d'un accidenté avant transport évite d'aggraver une lésion médullaire ; la chirurgie et les soins en centre spécialisé limitent les séquelles.
\item \textbf{Rééducation fonctionnelle} : kinésithérapie, stimulation des muscles paralysés et apprentissage de compensations permettent de limiter l'amyotrophie et de restaurer une autonomie partielle ; les réflexes conservés peuvent même être utilisés dans la rééducation.
\end{itemize}

\courssub{Synthèse : de la lésion aux troubles observés}

\begin{popfigure}
\begin{center}
\small
\begin{tabularx}{\linewidth}{|>{\raggedright\arraybackslash}p{2.6cm}|>{\raggedright\arraybackslash}X|>{\raggedright\arraybackslash}X|>{\raggedright\arraybackslash}p{2.4cm}|}
\hline
\rowcolor{popblueL} \textbf{Pathologie} & \textbf{Siège de la lésion} & \textbf{Troubles moteurs / sensitifs} & \textbf{État des réflexes} \\
\hline
Section de la moelle (traumatisme) & Voies ascendantes et descendantes médullaires & Paralysie \emph{et} anesthésie sous la lésion & Conservés, souvent \textbf{exagérés} (arc intact) \\
\hline
Poliomyélite & Motoneurones des cornes ventrales & Paralysie flasque + amyotrophie ; sensibilité conservée & \textbf{Abolis} \\
\hline
Tétanos & Interneurones inhibiteurs médullaires (blocage par la toxine) & Contractures permanentes ; pas de paralysie ni d'anesthésie & \textbf{Exagérés} \\
\hline
Section d'un nerf périphérique & Fibres motrices et sensitives du nerf & Paralysie \emph{et} anesthésie du territoire du nerf & \textbf{Abolis} \\
\hline
Lésion cérébrale (AVC) & Voies descendantes (pyramidales) & Paralysie controlatérale ; sensibilité variable & \textbf{Exagérés} \\
\hline
\end{tabularx}
\end{center}
\legende{Tableau de synthèse : correspondance entre le siège de la lésion, les troubles observés et l'état des réflexes. La règle d'or : réflexes \emph{abolis} = arc rompu ; réflexes \emph{exagérés} = arc intact libéré du contrôle cérébral.}
\end{popfigure}

\begin{aretenir}[L'essentiel de la section]
\begin{itemize}
\item Diagnostiquer, c'est \textbf{suivre l'arc réflexe} : récepteur $\rightarrow$ voie sensitive $\rightarrow$ centre $\rightarrow$ voie motrice $\rightarrow$ effecteur.
\item Paralysie avec sensibilité conservée = lésion \textbf{motrice} ; anesthésie sans paralysie = lésion \textbf{sensitive} ; paralysie + anesthésie localisées = lésion du \textbf{nerf mixte} ou de la \textbf{moelle}.
\item Réflexes \textbf{abolis} : l'arc est rompu (polio, nerf sectionné). Réflexes \textbf{conservés ou exagérés} : l'arc est intact, la lésion est au-dessus (section médullaire, tétanos, lésion cérébrale).
\item La limitation des dysfonctionnements repose sur la \textbf{vaccination} (polio, tétanos), la \textbf{sécurité routière}, l'\textbf{hygiène des plaies} et la \textbf{rééducation fonctionnelle}.
\end{itemize}
\end{aretenir}

\begin{exercice}[Pour t'entraîner]
Un patient présente, à la jambe droite : une paralysie flasque, une amyotrophie du quadriceps, une abolition du réflexe rotulien, mais une sensibilité cutanée parfaitement normale. (1) Quel élément de l'arc réflexe est lésé ? Justifie. (2) Propose une cause possible et un moyen de prévention.
\end{exercice}

\begin{corrige}[Exercice]
(1) La paralysie avec \emph{sensibilité conservée} évoque une lésion de la voie motrice ; l'\emph{abolition} du réflexe rotulien montre que l'arc est rompu au niveau du centre ou du motoneurone ; l'amyotrophie confirme la privation durable d'influx moteur. La lésion siège donc au niveau des \textbf{motoneurones des cornes ventrales} de la moelle lombaire (ou de la racine ventrale). (2) Cause possible : la \textbf{poliomyélite} (destruction virale des motoneurones). Prévention : la \textbf{vaccination antipoliomyélitique} selon le calendrier du PEV, avec rappels.
\end{corrige}

\newpage

\coursec{Activité d'intégration}

\coursec{Leçon N07 : Limitation des dysfonctionnements des organes ou structures intervenant dans les mouvements réflexes}

\courssub{Objectifs de l'activité}

À l'issue de cette activité d'intégration, tu seras capable de :

\begin{itemize}
\item interpréter les résultats d'expériences de section des racines rachidiennes (type Bell et Magendie) et d'expériences de dégénérescence wallérienne ;
\item exploiter les observations d'expériences de type Pflüger sur la grenouille pour dégager les éléments de l'arc réflexe et vérifier les lois de Pflüger ;
\item réaliser le schéma légendé d'un arc réflexe médullaire et d'un arc réflexe myotatique ;
\item localiser une lésion nerveuse à partir de l'examen clinique (état des réflexes, de la motricité et de la sensibilité) ;
\item expliquer des cas de paralysie ou d'absence de réflexe et proposer des mesures de prévention et de limitation des dysfonctionnements.
\end{itemize}

\courssub{Situation}

Tu es élève en Terminale D et tu effectues une visite pédagogique au service de neurologie de l'hôpital régional de Bamenda (région du Nord-Ouest). Le médecin-chef, le Dr Nkwenti, te confie un dossier documentaire en trois parties et te demande de rédiger un \textbf{compte rendu médical argumenté}.

\textbf{Partie 1 : expérience sur les racines rachidiennes.} Sur un chien anesthésié, on réalise au niveau lombaire les opérations suivantes, puis on excite les moignons des racines sectionnées. Les résultats sont consignés dans le tableau ci-dessous.

\begin{center}
\begin{tabularx}{\linewidth}{|l|X|X|}
\hline
\rowcolor{popblueL} \textbf{Opération} & \textbf{Excitation du moignon central (côté centre nerveux)} & \textbf{Excitation du moignon périphérique (côté membre)} \\
\hline
Section de la racine dorsale seule & Douleur intense (l'animal crie, se débat) ; aucun mouvement du membre & Aucune douleur ; aucun mouvement \\
\hline
Section de la racine ventrale seule & Aucune douleur ; aucun mouvement & Aucune douleur ; contraction des muscles du membre \\
\hline
Section des deux racines & Douleur ; aucun mouvement & Aucune douleur ; aucun mouvement \\
\hline
\end{tabularx}
\end{center}

Après section de la racine dorsale, on laisse l'animal vivre quelques jours : on observe au microscope une dégénérescence des fibres du moignon \emph{périphérique} de la racine dorsale et des fibres sensitives du nerf. Après section de la racine ventrale, la dégénérescence atteint le moignon périphérique de la racine ventrale et les fibres motrices du nerf (dégénérescence wallérienne).

\textbf{Partie 2 : expériences de type Pflüger sur la grenouille.} Une grenouille est décérébrée (destruction de l'encéphale, moelle épinière intacte). On applique sur la peau de la cuisse droite un papier imbibé d'acide acétique dilué : la patte droite se replie et chasse le papier (réflexe de défense). On augmente la concentration de l'acide : les deux pattes postérieures, puis les quatre pattes, entrent en mouvement. On excite la face ventrale de la cuisse : la patte se fléchit ; on excite la face dorsale : la patte s'étend. Après destruction de la moelle épinière (démédullation), aucune excitation, même forte, ne provoque plus de mouvement. Enfin, après section du nerf sciatique droit seul, l'excitation de la cuisse droite ne provoque plus de réponse du côté droit, mais la patte gauche répond encore.

\textbf{Partie 3 : trois cas cliniques.}

\textbf{Cas A.} M. T., 28 ans, moto-taximan à Bamenda, victime d'un accident de la route (collision avec un camion à l'entrée de la ville). Fracture de la 12\textsuperscript{ème} vertèbre dorsale avec section complète de la moelle épinière. Examen un mois après : paralysie des deux membres inférieurs (paraplégie), sensibilité totalement abolie au-dessous de la lésion, mais \textbf{réflexes rotuliens exagérés} (hyperréflexie) et muscles des membres inférieurs hypertoniques.

\textbf{Cas B.} Petite A., 5 ans, non vaccinée, originaire d'un village de la Menoua. Elle présente une paralysie \emph{flasque} du membre inférieur gauche : muscles atrophiés et mous, \textbf{réflexe achilléen aboli} du côté gauche, mais sensibilité cutanée parfaitement normale. Le diagnostic du laboratoire confirme la \textbf{poliomyélite} (le poliovirus détruit sélectivement les motoneurones des cornes ventrales de la moelle épinière).

\textbf{Cas C.} M. K., 45 ans, cultivateur à Bafoussam, s'est enfoncé un clou rouillé dans la plante du pied deux semaines plus tôt. Il présente des contractures douloureuses généralisées : muscles de la mâchoire crispés (trismus), cambrure du dos (opisthotonos), déclenchées par le moindre bruit. Diagnostic : \textbf{tétanos}. La toxine tétanique, produite par la bactérie \emph{Clostridium tetani} présente dans le sol, remonte le long des nerfs moteurs jusqu'à la moelle épinière où elle bloque les interneurones inhibiteurs : les motoneurones, libérés de tout frein, excitent les muscles de façon permanente et incontrôlée.

\courssub{Consignes}

\begin{enumerate}
\item À partir du tableau de la partie 1, déduis le rôle fonctionnel de chaque racine rachidienne (dorsale et ventrale). Explique comment les observations de dégénérescence wallérienne confirment ce rôle et précisent l'origine des fibres.
\item Dans la partie 2, montre que le réflexe de défense de la grenouille est un réflexe \emph{médullaire}. Identifie, dans l'ordre, les cinq éléments de l'arc réflexe mis en jeu et justifie le rôle de chacun à partir des expériences.
\item Retrouve dans les observations de la partie 2 les lois de Pflüger (loi de l'unilatéralité, loi de la symétrie, loi de l'irradiation, loi de l'adéquation) en citant l'observation correspondante.
\item Pour chacun des trois cas cliniques (A, B, C) : (a) réalise le schéma légendé de l'arc réflexe concerné et indique par une croix le siège de la lésion ; (b) explique l'état des réflexes (exagéré, aboli), de la motricité et de la sensibilité ; (c) propose au moins une mesure de prévention ou de limitation du dysfonctionnement.
\item Rédige ta conclusion sous la forme d'un compte rendu médical d'une dizaine de lignes, adressé au Dr Nkwenti.
\end{enumerate}

\courssub{Ressources à mobiliser}

\begin{aretenir}[Ce qu'il faut avoir en tête]
\begin{itemize}
\item \textbf{Racines rachidiennes} : la racine \cle{dorsale} est sensitive (elle contient les prolongements des neurones sensitifs dont les corps cellulaires sont dans le ganglion spinal) ; la racine \cle{ventrale} est motrice (elle contient les axones des motoneurones des cornes ventrales). Loi de Bell et Magendie.
\item \textbf{Dégénérescence wallérienne} : une fibre nerveuse sectionnée dégénère du côté séparé de son corps cellulaire ; elle permet donc de localiser le neurone d'origine d'une fibre.
\item \textbf{Arc réflexe} : récepteur $\rightarrow$ nerf sensitif (racine dorsale, ganglion spinal) $\rightarrow$ centre nerveux (substance grise de la moelle, avec ou sans interneurone) $\rightarrow$ nerf moteur (racine ventrale) $\rightarrow$ effecteur (muscle).
\item \textbf{Lois de Pflüger} : unilatéralité, symétrie, irradiation (selon l'intensité de l'excitation), adéquation de la réponse au point excité.
\item \textbf{Réflexe myotatique} (rotulien, achilléen) : arc monosynaptique ; le récepteur est le fuseau neuromusculaire ; il ne nécessite que la moelle épinière, mais il est normalement modéré (freiné) par l'encéphale.
\item \textbf{Interprétation clinique} : réflexe aboli = lésion d'un élément de l'arc (paralysie flasque) ; réflexe exagéré avec paralysie = lésion des voies descendantes entre l'encéphale et la moelle, l'arc médullaire restant intact (paralysie spastique) ; sensibilité abolie = atteinte des voies sensitives.
\end{itemize}
\end{aretenir}

\courssub{Démarche et corrigé commenté}

\begin{exoresolu}[Compte rendu médical argumenté — corrigé détaillé]
\textbf{Énoncé.} Interpréter les trois parties du dossier de l'hôpital régional de Bamenda et rédiger le compte rendu demandé par le Dr Nkwenti.

\tcblower

\textbf{Question 1 — Rôle des racines rachidiennes.}

Après section de la racine dorsale, l'excitation du moignon central provoque une douleur mais aucun mouvement : la racine dorsale conduit donc l'influx \emph{vers} le centre nerveux, c'est une racine \cle{sensitive}. L'excitation de son moignon périphérique ne provoque rien, car l'influx sensitif ne peut plus atteindre la moelle. Après section de la racine ventrale, seule l'excitation du moignon périphérique provoque la contraction du muscle : la racine ventrale conduit l'influx \emph{de la moelle vers} le muscle, c'est une racine \cle{motrice}.

La dégénérescence wallérienne confirme et précise : après section dorsale, les fibres qui dégénèrent sont celles du moignon périphérique, donc séparées de leurs corps cellulaires situés dans le \textbf{ganglion spinal} ; après section ventrale, les fibres qui dégénèrent sont celles du moignon périphérique, séparées des corps cellulaires des \textbf{motoneurones des cornes ventrales}. C'est la loi de Bell et Magendie : \emph{racine dorsale sensitive, racine ventrale motrice}.

\begin{popfigure}
\begin{tikzpicture}[
    neuron/.style={circle, draw=popdark, thick, minimum size=8mm, fill=popblueL},
    ganglion/.style={ellipse, draw=popdark, thick, minimum width=12mm, minimum height=8mm, fill=poporangeL},
    muscle/.style={rectangle, draw=popdark, thick, minimum width=10mm, minimum height=6mm, fill=popgreenL},
    receptor/.style={diamond, draw=popdark, thick, minimum size=8mm, fill=poppinkL},
    arrow/.style={->, thick, popdark},
    lesion/.style={red, very thick, dash pattern=on 4pt off 2pt}
]
% Recepteur
\node[receptor, align=center] (rec) at (0,0){Récepteur\\cutané};
% Neurone sensitif
\node[neuron, align=center] (sns) at (3,1.5){Neurone\\sensitif};
\node[ganglion, align=center] (gang) at (3,-1.5){Ganglion\\spinal};
% Moelle
\node[draw=popdark, thick, rounded corners, minimum width=2cm, minimum height=1.5cm, fill=popcream, align=center] (moelle) at (6,0){Substance grise\\médullaire};
% Neurone moteur
\node[neuron] (mot) at (9,0) {Motoneurone};
% Muscle
\node[muscle, align=center] (mus) at (12,0){Effecteur\\(Muscle)};

% Connexions
\draw[arrow] (rec) -- node[above, sloped, pos=0.5] {Influx sensitif} (sns);
\draw[arrow] (sns) -- (gang);
\draw[arrow] (gang) -- node[above] {Racine dorsale} (moelle);
\draw[arrow] (moelle) -- node[above] {Racine ventrale} (mot);
\draw[arrow] (mot) -- node[above] {Influx moteur} (mus);

% Légendes racines
\node[align=center, popmuted, font=\small] at (4.5, -2.5) {Racine dorsale\\(sensitive)};
\node[align=center, popmuted, font=\small] at (7.5, -2.5) {Racine ventrale\\(motrice)};
\end{tikzpicture}
\legende{Schéma de l'arc réflexe médullaire (polysynaptique type défense). La racine dorsale (sensitive) entre par le ganglion spinal ; la racine ventrale (motrice) sort de la corne ventrale.}
\end{popfigure}

\textbf{Question 2 — Le réflexe de défense est médullaire ; éléments de l'arc.}

La grenouille décérébrée (sans encéphale) répond encore à l'excitation acide : le réflexe ne nécessite donc pas l'encéphale. Après démédullation, plus aucune réponse : le centre du réflexe est donc la \textbf{moelle épinière}. Après section du nerf sciatique droit, la patte droite ne répond plus alors que la moelle est intacte : le nerf est la \textbf{voie de conduction} indispensable. L'excitation acide de la peau montre l'existence d'un \textbf{récepteur} cutané, et la réponse (mouvement) celle d'un \textbf{effecteur} musculaire. Les cinq éléments de l'arc réflexe sont donc : récepteur (peau) $\rightarrow$ fibre sensitive (racine dorsale) $\rightarrow$ centre nerveux (substance grise médullaire) $\rightarrow$ fibre motrice (racine ventrale) $\rightarrow$ effecteur (muscle de la patte).

\begin{popfigure}
\begin{tikzpicture}[
    neuron/.style={circle, draw=popdark, thick, minimum size=7mm, fill=popblueL},
    spindle/.style={ellipse, draw=popdark, thick, minimum width=10mm, minimum height=5mm, fill=poppinkL},
    muscle/.style={rectangle, draw=popdark, thick, minimum width=10mm, minimum height=6mm, fill=popgreenL},
    arrow/.style={->, thick, popdark},
    snapse/.style={circle, draw=poppurple, fill=poppurpleL, minimum size=3mm}
]
% Fuseau
\node[spindle, align=center] (fus) at (0,0){Fuseau\\neuromusculaire};
% Neurone sensitif
\node[neuron, align=center] (sns) at (3,1){Neurone\\sensitif};
% Interneurone (optionnel pour polysynaptique, mais myotatique = monosynaptique)
% Ici on montre la synapse directe
\node[snapse] (syn) at (5.5,1) {};
% Motoneurone
\node[neuron, align=center] (mot) at (7,1){Motoneurone\\(\alpha)};
% Muscle effecteur
\node[muscle, align=center] (mus) at (10,1){Muscle\\effecteur};
% Neurone gamma (fusimoteur) - optionnel mais bon pour completude
\node[neuron, fill=poporangeL, align=center] (gam) at (7,-1.5){Motoneurone\\($\gamma$)};
\node[muscle, fill=poporangeL, align=center] (fib) at (0,-1.5){Fibres\\intrafusales};

\draw[arrow] (fus) -- node[above] {Influx sensitif (Ia)} (sns);
\draw[arrow] (sns) -- (syn);
\draw[arrow] (syn) -- (mot);
\draw[arrow] (mot) -- node[above] {Influx moteur} (mus);
% Boucle gamma
\draw[arrow, poporange] (gam) .. controls (5,-1.5) and (2,-1.5) .. (fib);
\draw[arrow, poporange] (fib) .. controls (2,-0.5) and (3,0.5) .. (sns);

\node[align=center, font=\small, popmuted] at (5, -2.5) {Arc monosynaptique (réflexe myotatique)};
\node[align=center, font=\small, poporange] at (10, -2) {Commande $\gamma$ : réglage de la sensibilité du fuseau};
\end{tikzpicture}
\legende{Schéma de l'arc réflexe myotatique (rotulien ou achilléen). Arc monosynaptique : le neurone sensitif (Ia) fait synapse directement sur le motoneurone $\alpha$. Le motoneurone $\gamma$ innerve les fibres intrafusales du fuseau pour en maintenir la sensibilité.}
\end{popfigure}

\textbf{Question 3 — Lois de Pflüger.}

\begin{itemize}
\item \emph{Loi de l'unilatéralité} : excitation faible $\rightarrow$ seule la patte droite excitée répond.
\item \emph{Loi de la symétrie} : excitation plus forte $\rightarrow$ la patte controlatérale (gauche) répond aussi.
\item \emph{Loi de l'irradiation} : excitation très forte $\rightarrow$ les quatre pattes répondent (la réponse s'étend à des territoires de plus en plus larges).
\item \emph{Loi de l'adéquation} : la réponse est adaptée au point excité (flexion pour la face ventrale, extension pour la face dorsale) et tend à écarter l'agent agressif.
\end{itemize}

\textbf{Question 4 — Cas cliniques.}

\emph{Cas A (paraplégie traumatique, section médullaire T12).}
L'arc réflexe rotulien (fuseau neuromusculaire du quadriceps $\rightarrow$ racine dorsale L2-L4 $\rightarrow$ moelle lombaire $\rightarrow$ racine ventrale L2-L4 $\rightarrow$ quadriceps) est \textbf{intact}, car la lésion siège au niveau dorsal (T12), au-dessus du centre lombaire. Mais les voies descendantes de l'encéphale vers la moelle (faisceaux corticospinaux, vestibulospinaux, reticulospinaux) sont sectionnées : l'encéphale ne commande plus les mouvements volontaires (paralysie) et ne freine plus le réflexe myotatique via les voies inhibitrices descendantes, d'où l'\textbf{hyperréflexie rotulienne} et l'hypertonie (paralysie \emph{spastique}). Les voies sensitives ascendantes (faisceaux spinothalamiques, colonnes postérieures) étant coupées, la sensibilité est abolie sous la lésion.
\textbf{Prévention / Limitation :} port du casque et du gilet de protection, respect du code de la route, limitation de vitesse des moto-taxis ; limitation des séquelles : immobilisation correcte du rachis lors du transport (plan dur, collier cervical), rééducation précoce en kinésithérapie pour lutter contre les rétractions et maintenir la trophicité.

\begin{popfigure}
\begin{tikzpicture}[
    neuron/.style={circle, draw=popdark, thick, minimum size=7mm, fill=popblueL},
    spindle/.style={ellipse, draw=popdark, thick, minimum width=10mm, minimum height=5mm, fill=poppinkL},
    muscle/.style={rectangle, draw=popdark, thick, minimum width=10mm, minimum height=6mm, fill=popgreenL},
    arrow/.style={->, thick, popdark},
    lesion/.style={red, very thick}
]
% Arc reflexe intact (niveau lombaire)
\node[spindle, align=center] (fus) at (0,0){Fuseau\\quadriceps};
\node[neuron, align=center] (sns) at (3,0.8){Neurone\\sensitif};
\node[neuron] (mot) at (6,0.8) {Motoneurone $\alpha$};
\node[muscle] (mus) at (9,0.8) {Quadriceps};

\draw[arrow] (fus) -- (sns);
\draw[arrow] (sns) -- (mot);
\draw[arrow] (mot) -- (mus);

% Voies descendantes sectionnées
\node[draw=popdark, thick, rounded corners, fill=popcream, minimum width=1.5cm, minimum height=1cm, align=center] (cerveau) at (6, -2.5){Encéphale\\(Centres supérieurs)};
\draw[lesion] (4.5, -1.8) -- (4.5, -0.5) node[midway, left, red, font=\small, align=center]{Lésion T12\\$\times$};
\draw[arrow, densely dashed, popmuted] (cerveau) -- node[left, popmuted, font=\small, align=center]{Voies descendantes\\inhibitrices (coupées)} (4.5, -1.8);

% Voies sensitives montantes coupées
\draw[lesion] (4.5, 2.5) -- (4.5, 1.2) node[midway, right, red, font=\small, align=center]{Lésion T12\\$\times$};
\draw[arrow, densely dashed, popmuted] (4.5, 2.5) -- node[right, popmuted, font=\small, align=center]{Voies sensitives\\montantes (coupées)} (4.5, 3.5);

\node[align=center, font=\small, popmuted] at (4.5, 3.8) {Sensibilité abolie};
\node[align=center, font=\small, popmuted] at (4.5, -3.3) {Paralysie volontaire + Hyperréflexie};
\end{tikzpicture}
\legende{Cas A : Section médullaire T12. L'arc réflexe lombaire (L2-L4) est intact (traits pleins). Les voies descendantes (contrôle volontaire et inhibition) et ascendantes (sensibilité) sont interrompues au niveau de la lésion (croix rouges). Résultat : paraplégie spastique, hyperréflexie, anesthésie sous-lésionnelle.}
\end{popfigure}

\emph{Cas B (poliomyélite, destruction des motoneurones cornes ventrales).}
Le virus détruit les \textbf{motoneurones $\alpha$ des cornes ventrales} au niveau lombaire (L4-S1 pour le réflexe achilléen) : la branche motrice de l'arc achilléen est rompue, d'où l'abolition du réflexe achilléen et la paralysie \emph{flasque} avec atrophie musculaire par dénutrition (le muscle n'est plus stimulé par l'influx nerveux). La voie sensitive (neurone sensitif, ganglion spinal, racine dorsale, voies ascendantes) est épargnée : la sensibilité reste normale. C'est la preuve clinique de la loi de Bell et Magendie (séparation anatomique et fonctionnelle des racines).
\textbf{Prévention / Limitation :} \textbf{vaccination antipoliomyélitique} systématique (vaccin oral bivalent ou inactivé selon calendrier EPI Cameroun, campagnes de vaccination de routine et Journées nationales de vaccination) ; surveillance épidémiologique des paralysies flasques aiguës (PFA).

\begin{popfigure}
\begin{tikzpicture}[
    neuron/.style={circle, draw=popdark, thick, minimum size=7mm, fill=popblueL},
    spindle/.style={ellipse, draw=popdark, thick, minimum width=10mm, minimum height=5mm, fill=poppinkL},
    muscle/.style={rectangle, draw=popdark, thick, minimum width=10mm, minimum height=6mm, fill=popgreenL},
    arrow/.style={->, thick, popdark},
    lesion/.style={red, very thick}
]
% Arc reflexe : partie sensitive OK
\node[spindle, align=center] (fus) at (0,0){Fuseau\\triceps sural};
\node[neuron, align=center] (sns) at (3,0){Neurone\\sensitif (intact)};
% Synapse
\node[circle, draw=poppurple, fill=poppurpleL, minimum size=3mm] (syn) at (5,0) {};
% Motoneurone DETRUIT
\node[neuron, fill=popred!20, draw=red, thick, align=center] (mot) at (7,0){Motoneurone $\alpha$\\DÉTRUIT};
\node[muscle, fill=popred!10, draw=red, align=center] (mus) at (10,0){Muscle\\atrophié (flasque)};

\draw[arrow] (fus) -- (sns);
\draw[arrow] (sns) -- (syn);
\draw[lesion] (syn) -- (mot) node[midway, above, red, font=\small] {$\times$ Pas de synapse};
\draw[lesion, densely dashed] (mot) -- (mus) node[midway, above, red, font=\small] {$\times$ Pas d'influx};

\node[align=center, font=\small, popgreen] at (1.5, -1.5) {Voie sensitive\\INTACTE};
\node[align=center, font=\small, popred] at (8.5, -1.5) {Voie motrice\\ROMPUE};
\end{tikzpicture}
\legende{Cas B : Poliomyélite. Destruction sélective des motoneurones $\alpha$ des cornes ventrales (croix rouge). L'arc réflexe est rompu au niveau du motoneurone. Réflexe aboli, paralysie flasque, atrophie, sensibilité conservée.}
\end{popfigure}

\emph{Cas C (tétanos, blocage de l'inhibition médullaire).}
La toxine tétanique (tétanospasmine) migre par transport rétrograde le long des axones moteurs jusqu'aux corps cellulaires des motoneurones et des \textbf{interneurones inhibiteurs} (glycinergiques/GABAergiques) de la moelle épinière et du tronc cérébral. Elle bloque la libération des neurotransmetteurs inhibiteurs (glycine, GABA) au niveau des synapses sur les motoneurones. Les motoneurones, libérés de tout frein inhibiteur (désinhibition), déchargent de façon permanente et incontrôlée à haute fréquence. L'arc réflexe est anatomiquement intact, mais sa régulation est abolie, d'où les contractures généralisées permanentes (trismus, opisthotonos, risus sardonicus) déclenchées par le moindre stimulus sensoriel (bruit, lumière).
\textbf{Prévention / Limitation :} vaccination antitétanique (rappels tous les 10-20 ans, vaccination de la femme enceinte pour le tétanos néonatal), sérothérapie (sérum antitétanique) et soins chirurgicaux immédiats des plaies profondes/souillées (débridement), port de chaussures de protection au champ, éducation sanitaire.

\begin{popfigure}
\begin{tikzpicture}[
    neuron/.style={circle, draw=popdark, thick, minimum size=7mm, fill=popblueL},
    spindle/.style={ellipse, draw=popdark, thick, minimum width=10mm, minimum height=5mm, fill=poppinkL},
    muscle/.style={rectangle, draw=popdark, thick, minimum width=10mm, minimum height=6mm, fill=popgreenL},
    interneuron/.style={circle, draw=poppurple, thick, minimum size=7mm, fill=poppurpleL},
    arrow/.style={->, thick, popdark},
    inhib/.style={->, thick, popred, dash pattern=on 2pt off 2pt},
    lesion/.style={red, very thick}
]
% Recepteur / Fuseau
\node[spindle, align=center] (fus) at (0,1.5){Récepteur\\cutané / Fuseau};
% Neurone sensitif
\node[neuron, align=center] (sns) at (3,1.5){Neurone\\sensitif};
% Interneurone inhibiteur BLOQUE
\node[interneuron, draw=red, thick, fill=popred!20, align=center] (inh) at (5.5, 0){Interneurone\\inhibiteur (Gly/GABA)\\\textbf{BLOQUÉ par toxine}};
% Motoneurone
\node[neuron] (mot) at (5.5, 2.5) {Motoneurone $\alpha$};
% Muscle
\node[muscle, align=center] (mus) at (8.5, 2.5){Muscle\\contracturé};

% Connexions
\draw[arrow] (fus) -- (sns);
\draw[arrow] (sns) -- (mot);
\draw[arrow] (mot) -- (mus);
% Inhibition normalement
\draw[inhib] (inh) -- node[left, popred, font=\small, align=center]{Inhibition\\BLOQUÉE} (mot);
% Excitation interneurone (normalement activé par neurone sensitif)
\draw[arrow, densely dashed, popmuted] (sns) -- node[below, popmuted, font=\small, align=center]{Excitation\\normalement} (inh);

\node[align=center, font=\small, popred] at (5.5, -1.2) {Blocage de la libération\\de Glycine / GABA};
\node[align=center, font=\small, popdark] at (7, 3.5) {Décharge continue\\non contrôlée};
\end{tikzpicture}
\legende{Cas C : Tétanos. La toxine bloque les interneurones inhibiteurs (croix rouge sur la flèche d'inhibition). Les motoneurones ne sont plus freinés : décharge continue $\rightarrow$ contractures généralisées (spasticité). Arc anatomiquement intact, mais dysfonctionnement physiologique.}
\end{popfigure}

\textbf{Question 5 — Compte rendu (modèle).}

« Monsieur le Docteur,

Suite à l'analyse du dossier documentaire, je vous confirme que les trois observations illustrent la nécessité de l'intégrité des cinq éléments de l'arc réflexe pour une motricité adaptée.

L'expérience de Bell et Magendie (Partie 1) établit la ségrégation fonctionnelle des racines rachidiennes : dorsale sensitive (corps cellulaires au ganglion spinal), ventrale motrice (corps cellulaires dans les cornes ventrales), confirmée par la dégénérescence wallérienne distale.

Les expériences de Pflüger (Partie 2) prouvent que le réflexe de défense est médullaire (persistant après décérébration, aboli après démédullation) et obéit aux quatre lois (unilatéralité, symétrie, irradiation, adéquation), définissant les propriétés des réseaux neuronaux médullaires.

Cliniquement, trois types de dysfonctionnements sont distingués :
\begin{itemize}
\item \textbf{Cas A (Section médullaire T12)} : Lésion des voies descendantes (pyramidales et extrapyramidales) et ascendantes. L'arc réflexe lombaire intact est libéré du contrôle inhibiteur encéphalique $\rightarrow$ paraplégie \emph{spastique} avec hyperréflexie rotulienne et anesthésie sous-lésionnelle.
\item \textbf{Cas B (Poliomyélite)} : Destruction sélective des motoneurones $\alpha$ (cornes ventrales). Rupture de l'arc réflexe $\rightarrow$ paralysie \emph{flasque}, aréflexie, amyotrophie, sensibilité conservée.
\item \textbf{Cas C (Tétanos)} : Neurotoxine bloquant les interneurones inhibiteurs médullaires. Désinhibition des motoneurones $\rightarrow$ contractures généralisées (trismus, opisthotonos) sans lésion anatomique de l'arc.
\end{itemize}

La limitation de ces dysfonctionnements repose sur la prévention primaire (vaccination antipoliomyélitique et antitétanique dans le PEV, sécurité routière pour les moto-taximans), la prise en charge urgente (sérothérapie antitetanique, immobilisation rachidienne) et la rééducation fonctionnelle précoce pour les séquelles motrices.

Respectueusement, » 
\end{exoresolu}

\begin{attention}[Pièges fréquents au Bac]
\begin{itemize}
\item Ne confonds pas \textbf{paralysie flasque} (réflexe aboli, hypotonie, atrophie : lésion de l'arc réflexe lui-même — motoneurone, racine ventrale, nerf périphérique) et \textbf{paralysie spastique} (réflexe exagéré, hypertonie, pas d'atrophie initiale : lésion des voies descendantes \emph{au-dessus} du centre réflexe).
\item Ne dis jamais que le réflexe rotulien « part du cerveau » ou « est commandé par le cerveau » : son centre est médullaire (L2-L4) ; l'encéphale ne fait que le \emph{moduler} (frein inhibiteur).
\item Dans les schémas, place toujours le \textbf{ganglion spinal} sur la racine dorsale (corps cellulaires des neurones sensitifs). La racine ventrale n'a pas de ganglion (corps cellulaires dans la corne ventrale).
\item Pour le tétanos, précise que la toxine bloque la \textbf{libération des neurotransmetteurs inhibiteurs} (glycine, GABA) au niveau des synapses des interneurones inhibiteurs sur les motoneurones, et non l'influx nerveux lui-même.
\item L'adéquation (loi de Pflüger) ne signifie pas « réponse proportionnelle » mais « réponse \emph{adaptée} au point stimulé pour écarter l'agent » (flexion/extension).
\end{itemize}
\end{attention}

\courssub{Bilan de l'activité}

Cette activité t'a permis de mobiliser simultanément tous les savoir-faire de la leçon : interpréter des expériences de section et de dégénérescence (Bell et Magendie, Waller), exploiter des expériences de type Pflüger pour dégager les éléments de l'arc réflexe et vérifier les lois des réflexes médullaires, schématiser un arc réflexe myotatique, et surtout \textbf{raisonner comme un clinicien} : l'état des réflexes (aboli, exagéré, normal) et de la sensibilité permet de localiser précisément une lésion nerveuse (névrite, radiculopathie, syndrome de la queue de cheval, hémisection médullaire, section complète, atteinte du motoneurone central vs périphérique). Tu as enfin compris que la limitation des dysfonctionnements du système nerveux repose d'abord sur la \textbf{prévention} : vaccination contre la poliomyélite et le tétanos (couverture vaccinale PEV Cameroun), sécurité routière pour les usagers des moto-taxis (port du casque, limitation vitesse), et rééducation fonctionnelle pour les accidentés médullaires.

\begin{savaistu}[Contexte camerounais : Vaccination et Sécurité routière]
Au Cameroun, le Programme Élargi de Vaccination (PEV) intègre le vaccin antipoliomyélitique (VPO bivalent ou VIP) et le vaccin antitétanique (associé DTC/HepB/Hib). Les Journées Nationales de Vaccination (JNV) visent à interrompre la circulation du poliovirus sauvage. Malgré la certification « Afrique exempte de polio sauvage » (2020), le risque de poliovirus dérivés de souches vaccinales (cVDPV) persiste, justifiant le maintien de la surveillance des PFA. Pour le tétanos, la vaccination de la femme enceinte (VAT 1 à 5) protège le nouveau-né (tétanos néonatal). Concernant les accidents de moto-taxi (\"benskin\"), très fréquents à Bamenda, Douala, Yaoundé, le non-port du casque et la vitesse excessive sont les premières causes de traumatismes rachidiens (T12-L1 souvent) entraînant des paraplégies spastiques définitives. La formation aux gestes de premiers secours (immobilisation sur plan dur) est cruciale.
\end{savaistu}

\newpage

\coursec{Méthodes et exercices résolus}

\coursec{Limitation des dysfonctionnements des organes ou structures intervenant dans les mouvements réflexes}

\courssub{Méthodes et exercices résolus}

\begin{aretenir}[Rappels indispensables : Centres nerveux, Neurone, Arc réflexe]
\textbf{1. Organisation des centres nerveux (Encéphale \& Moelle épinière) :}
\begin{itemize}
    \item \textbf{Encéphale :} Cerveau (2 hémisphères, cortex = substance grise, centre = substance blanche, corps calleux), Cervelet (arborisation blanche = arbre de vie, cortex grisé), Tronc cérébral (Mésencéphale, Pont, Bulbe). Méninges : Dure-mère, Arachnoïde (espace sous-arachnoïdien + LCR), Pie-mère.
    \item \textbf{Moelle épinière :} Coupe transversale : substance grise en \textbf{H} (cornes antérieures motrices, postérieures sensitives, latérales végétatives) entourée de substance blanche (faisceaux ascendants/descendants). Canal central (épendyme, LCR).
    \item \textbf{Nerf spinal (mixte) :} Réunion racine dorsale (sensitive, ganglion rachidien) + racine ventrale (motrice). Sort par le foramen de conjugaison.
\end{itemize}
\textbf{2. Neurone :} Unité structurelle et fonctionnelle. Corps cellulaire (périkaryon), dendrites (réception), axone (conduction, gaine de myéline = oligodendrocytes SNC / cellules de Schwann SNP), terminaisons synaptiques (émission). Types : Sensitif (pseudo-unipolaire, ganglion), Moteur (multipolaire, corne antérieure), Interneurone (multipolaire, substance grise, 99\%).
\textbf{3. Loi de Bell-Magendie (1822) :} \cle{Racine dorsale = Voie sensitive (afferente)} ; \cle{Racine ventrale = Voie motrice (efférente)}. Section racine dorsale $\rightarrow$ anesthésie/analgésie + aréflexie. Section racine ventrale $\rightarrow$ paralysie flasque, aréflexie, amyotrophie.
\textbf{4. Réflexe médullaire \& Lois de Pflüger (grenouille démédullée) :}
\begin{itemize}
    \item \cle{Loi 1 (Siège de l'excitation) :} Nature de la réponse = f(siège stimulation). Ex: Cuisse dorsale $\rightarrow$ Flexion (retrait) ; Plante pied $\rightarrow$ Extension (poussée).
    \item \cle{Loi 2 (État du centre) :} Intensité/facilité du réflexe = f(excitabilité centre médullaire). Fatigue réflexe (dépletion vésicules synaptiques) $\rightarrow$ disparition réponse.
\end{itemize}
\textbf{5. Arc réflexe médullaire (polysynaptique, ex: retrait) :} Récepteur $\rightarrow$ Neurone sensitif (Racine dorsale) $\rightarrow$ Interneurone(s) (Corne postérieure) $\rightarrow$ Motoneurone $\alpha$ (Corne antérieure) $\rightarrow$ Racine ventrale $\rightarrow$ Nerf spinal $\rightarrow$ Effecteur (Muscle squelettique).
\textbf{6. Réflexe myotatique (ostéotendineux, monosynaptique pour l'agoniste) :} Récepteur = \cle{Fuseau neuromusculaire} (fibres intrafusales, sensible à l'étirement). Neurone Ia (gros diamètre, myélinisé, \SI{80-120}{m/s}). Synapse \textbf{directe} (monosynaptique, Glutamate/AMPA) sur Motoneurone $\alpha$ agoniste. \textbf{Inhibition réciproque} (disynaptique) : Collatérale Ia $\rightarrow$ Interneurone inhibiteur (Glycine/GABA) $\rightarrow$ MN $\alpha$ antagoniste. \cle{Co-activation $\alpha$-$\gamma$} : MN $\gamma$ maintient tension fuseau pendant contraction.
\textbf{7. Réflexe conditionnel (Pavlov) :} Apprentissage associatif (SC + SI $\rightarrow$ RC). Règles : Contiguïté temporelle (SC avant SI), Répétition, Renforcement. Base neuronale : \cle{Convergence} voies SC/SI sur neurones détecteurs de coïncidence $\rightarrow$ \cle{LTP (Potentiation Long Terme)} (NMDA, $Ca^{2+}$, AMPA). Extinction = Nouvel apprentissage inhibiteur (LTD), pas oubli (Récupération spontanée).
\end{aretenir}

\begin{methode}[Réaliser la dissection de l'encéphale d'un mammifère et identifier les grandes régions]
\textbf{Objectif :} Observer l'organisation macroscopique de l'encéphale (cerveau, cervelet, bulbe) et les méninges.
\begin{enumerate}
    \item \textbf{Matériel :} Encéphale de mouton ou de bovin (fraîchement prélevé à l'abattoir de Ngaoundéré ou Yaoundé), plateau de dissection, instruments (scalpel, ciseaux, forceps, aiguilles), gants, lunettes, eau physiologique (\SI{0,9}{\%} NaCl).
    \item \textbf{Observation externe (Vue dorsale) :} Placer l'encéphale face dorsale vers le haut. Identifier les deux hémisphères cérébraux (convolutions/gyri et sillons/sulci), le cervelet (arborisation blanche caractéristique = arbre de vie) en arrière, le bulbe rachidien (continuité avec la moelle épinière).
    \item \textbf{Repérage des méninges :} Observer la dure-mère (épaisse, blanche, adhérente au crâne), l'arachnoïde (fine, transparente, espace sous-arachnoïdien contenant le LCR), la pie-mère (très fine, vascularisée, collée à la surface du cerveau, suit les circonvolutions).
    \item \textbf{Coupe sagittale médiane :} Inciser le long de la grande fente longitudinale pour séparer les deux hémisphères. Observer :
    \begin{itemize}
        \item Substance blanche (centre, fibres myélinisées : faisceaux d'association, commissures, projection).
        \item Substance grise (cortex, surface, corps cellulaires).
        \item Corps calleux (grande commissure blanche transversale).
        \item Ventricules : Latéraux (dans hémisphères), IIIe (diencephale), Aqueduc de Sylvius (mésencéphale), IVe (sous cervelet, losange).
    \end{itemize}
    \item \textbf{Coupe transversale (facultatif, niveau cervelet) :} Observer l'arbre de vie (substance blanche) et le cortex cérébelleux (substance grise plissée).
    \item \textbf{Nettoyage :} Éliminer les déchets biologiques selon le protocole du laboratoire (incinération/bac DASRI), désinfecter le matériel (bains désinfectants).
\end{enumerate}
\textbf{Consigne de dessin (Bac) :} Schéma annoté en vue dorsale ET coupe sagittale médiane. Échelle respectée. Légende numérotée : 1. Hémisphère cérébral, 2. Cervelet, 3. Bulbe, 4. Dure-mère, 5. Corps calleux, 6. Ventricule latéral, 7. IIIe ventricule, 8. Aqueduc de Sylvius, 9. IVe ventricule, 10. Substance blanche, 11. Substance grise (cortex).
\end{methode}

\begin{exoresolu}[Identification des structures sur une coupe sagittale médiane d'encéphale de mammifère]
\textbf{Énoncé :} On présente à un élève une coupe sagittale médiane d'un encéphale de mammifère (document photographique). L'élève doit :
\begin{enumerate}
    \item Nommer les structures repérées par les lettres A, B, C, D, E, F.
    \item Préciser la nature (substance grise ou blanche) et la composition cellulaire de la structure B.
    \item Expliquer le rôle fonctionnel de la structure C.
    \item Justifier pourquoi une lésion de la structure D entraîne une paralysie hémicorporelle controlatérale.
\end{enumerate}
\textit{Document (description textuelle) :} A : Partie antérieure bombée, surface plissée. B : Couche superficielle de A, couleur grise. C : Large bande blanche transversale reliant les deux hémisphères sous B. D : Partie descendante de C vers le tronc cérébral. E : Structure arborisée blanche postérieure, à surface plissée. F : Cavité triangulaire sous E, tapissée d'épendyme.

\tcblower
\textbf{Solution détaillée :}
\begin{enumerate}
    \item \textbf{Identification :}
    \begin{itemize}
        \item A : \textbf{Hémisphère cérébral} (Télencéphale).
        \item B : \textbf{Cortex cérébral} (Substance grise corticale).
        \item C : \textbf{Corps calleux} (Grande commissure interhémisphérique).
        \item D : \textbf{Pédoncule cérébral} (Crus cerebri, contenant les voies pyramidales/corticospinales).
        \item E : \textbf{Cervelet} (Métencéphale, arborisation blanche = arbre de vie).
        \item F : \textbf{IVe ventricule} (Cavité du losange, continuation du canal central de la moelle).
    \end{itemize}
    \item \textbf{Nature et composition de B (Cortex) :} C'est de la \textbf{substance grise}. Elle est constituée principalement de \textbf{corps cellulaires de neurones} (pyramidaux, stellaires, fusiformes), de dendrites, de terminaisons synaptiques, de cellules gliales (astrocytes, oligodendrocytes, microglie) et de capillaires sanguins. L'absence de myéline (gainée d'oligodendrocytes) explique la couleur grise.
    \item \textbf{Rôle de C (Corps calleux) :} Assurer la \textbf{communication interhémisphérique}. Il transmet les informations sensorielles, motrices et cognitives entre l'hémisphère gauche et l'hémisphère droit (intégration bilatérale, coordination bimanuelle, transfert de la mémoire, du langage).
    \item \textbf{Justification de la paralysie controlatérale par lésion de D :}
    \begin{itemize}
        \item D (Pédoncule cérébral) contient les \textbf{voies pyramidales (corticospinales)} issues du cortex moteur (aire de Brodmann 4).
        \item Ces fibres \textbf{décussent (croisent) au niveau du bulbe rachidien} (décussation des pyramides) : 80 à 85\% croisent controlatéralement (voie pyramidale croisée), 15 à 20\% restent ipsilatérales (voie pyramidale directe).
        \item Une lésion \textbf{en amont de la décussation} (au niveau du pédoncule D) interrompt les fibres \textbf{avant leur croisement}.
        \item Les motoneurones alpha de la moelle épinière (corne antérieure) ne reçoivent plus l'influx volontaire controlatéral $\rightarrow$ \textbf{paralysie flasque puis spastique controlatérale} (hémiparésie/hémiplégie controlatérale).
    \end{itemize}
\end{enumerate}
\textbf{Conclusion :} La connaissance de l'anatomie macroscopique et du trajet des voies nerveuses (notamment la décussation des pyramides) est indispensable pour localiser une lésion neurologique (diagnostic topique) et prédire les déficits moteurs (controlatéral pour voie pyramidale sus-bulbaire).
\end{exoresolu}

\begin{methode}[Démédullation et décérébration du batracien : mise en évidence du réflexe médullaire]
\textbf{Objectif :} Isoler la moelle épinière du contrôle encéphalique pour étudier les réflexes médullaires purs (lois de Pflüger).
\begin{enumerate}
    \item \textbf{Matériel :} Grenouille vivante (\textit{Xenopus laevis} ou \textit{Rana ridibunda} locale), plateau en liège, aiguilles, ciseaux, pince, solution de Ringer (NaCl \SI{0,65}{\%}, KCl \SI{0,014}{\%}, CaCl2 \SI{0,012}{\%}, pH 7,4), marteau à réflexes ou aiguille fine.
    \item \textbf{Démédullation (Section moelle/bulbe) :}
    \begin{itemize}
        \item Piquer le cerveau (destruction hémisphères + bulbe) par le foramen magnum avec une aiguille émoussée $\rightarrow$ \textbf{animal démédullé}. La moelle épinière est intacte mais isolée de l'encéphale.
        \item \textit{Vérification :} Perte des réflexes de redressement, de la respiration buccale, conservation des réflexes de retrait des membres postérieurs (pinçage orteil).
    \end{itemize}
    \item \textbf{Décérébration (Section entre mésencéphale et bulbe) :}
    \begin{itemize}
        \item Section transversale entre corps quadrijumeaux (tectum) et bulbe.
        \item \textit{Vérification :} Rigidité décérébrée (extension tonique des 4 membres), conservation réflexes médullaires, respiration conservée (centres bulbaire/pontiques intacts).
    \end{itemize}
    \item \textbf{Mise en évidence du réflexe médullaire (Lois de Pflüger) :}
    \begin{itemize}
        \item Stimuler la peau de la cuisse (piqûre aiguille / acide acétique \SI{1}{\%}) $\rightarrow$ Flexion de la jambe (retrait).
        \item Stimuler la plante du pied $\rightarrow$ Extension de la jambe (poussée).
        \item \textbf{Interprétation :} Le réflexe dépend du \textbf{siège de la stimulation} (Loi 1) et de l'\textbf{état du centre} (fatigue, lésion) (Loi 2). L'arc réflexe est intramédullaire.
    \end{itemize}
    \item \textbf{Éthique :} Anesthésie préalable (MS-222 / tricaïne méthanesulfonate \SI{0,1}{\%}) obligatoire selon déontologie. Euthanasie finale par double démédullation + destruction hémisphères.
\end{enumerate}
\textbf{Compétence Bac :} Schématiser l'arc réflexe médullaire (neurone sensitif, interneurone, motoneurone, effecteur) et expliquer les lois de Pflüger.
\end{methode}

\begin{exoresolu}[Interprétation des expériences de Pflüger sur la grenouille démédullée]
\textbf{Énoncé :} On réalise les expériences suivantes sur une grenouille démédullée maintenue humide en solution de Ringer :
\begin{enumerate}
    \item \textbf{Expérience 1 :} Piqûre de la face dorsale de la cuisse gauche $\rightarrow$ Flexion rapide de la jambe gauche (retrait).
    \item \textbf{Expérience 2 :} Piqûre de la face ventrale (plante) du pied gauche $\rightarrow$ Extension vigoureuse de la jambe gauche (poussée).
    \item \textbf{Expérience 3 :} Après l'expérience 1, on répète immédiatement la piqûre dorsale cuisse gauche toutes les 30 s. L'amplitude de la flexion diminue progressivement jusqu'à disparition.
    \item \textbf{Expérience 4 :} On sectionne les racines dorsales (sensitives) du nerf sciatique gauche. Piqûre cuisse gauche $\rightarrow$ Aucune réponse. Piqûre cuisse droite $\rightarrow$ Flexion jambe droite.
\end{enumerate}
Questions :
\begin{enumerate}
    \item Énoncer les deux lois de Pflüger illustrées par les Expériences 1 et 2.
    \item Expliquer le mécanisme physiologique de l'Expérience 3 (nommer le phénomène).
    \item Interpréter les résultats de l'Expérience 4 en précisant le rôle des racines rachidiennes.
    \item Schématiser l'arc réflexe médullaire complet pour l'Expérience 1 (5 éléments).
\end{enumerate}

\tcblower
\textbf{Solution détaillée :}
\begin{enumerate}
    \item \textbf{Lois de Pflüger :}
    \begin{itemize}
        \item \textbf{Loi 1 (Loi du siège de l'excitation) :} \textit{« La nature de la réponse réflexe dépend du siège de la stimulation. »} Exp 1 (cuisse dorsale) $\rightarrow$ Flexion (retrait). Exp 2 (pied ventral) $\rightarrow$ Extension (poussée/appui). Les circuits interneuronaux médullaires diffèrent selon le dermatome stimulé.
        \item \textbf{Loi 2 (Loi de l'état du centre) :} \textit{« L'intensité et la facilité d'obtention du réflexe dépendent de l'excitabilité du centre médullaire. »} (Illustrée par Exp 3).
    \end{itemize}
    \item \textbf{Expérience 3 : Fatigue réflexe (ou habituation synaptique).}
    \begin{itemize}
        \item Mécanisme : Dépletion des vésicules synaptiques (neurotransmetteur glutamate/aspartate) au niveau des synapses \textbf{neurone sensitif $\rightarrow$ interneurone} et \textbf{interneurone $\rightarrow$ motoneurone}.
        \item Fréquence de stimulation (toutes les 30 s) $>$ vitesse de recyclage des vésicules (reprise de la choline, resynthèse, rechargement).
        \item Diminution de l'EPSP (Potentiel Post-Synaptique Excitateur) sous le seuil de déclenchement du potentiel d'action chez le motoneurone $\rightarrow$ disparition de la réponse motrice. Récupération après repos (temps de recharge).
    \end{itemize}
    \item \textbf{Expérience 4 : Rôle des racines rachidiennes (Confirmation Bell-Magendie).}
    \begin{itemize}
        \item Section racines dorsales (ganglion rachidien) gauche $\rightarrow$ Section de l'\textbf{afference sensitive} (neurones sensitifs pseudo-unipolaires). L'influx ne peut pas entrer dans la moelle $\rightarrow$ \textbf{Abolition du réflexe} côté lésé (jambe gauche).
        \item Racines ventrales (motrices) intactes $\rightarrow$ Les motoneurones sont vivants mais ne reçoivent plus l'influx.
        \item Côté droit intact : Racines dorsales et ventrales fonctionnelles $\rightarrow$ Réflexe normal (flexion jambe droite). Preuve que l'arc réflexe est \textbf{segmentaire et unilatéral} pour ce réflexe de retrait simple (pas de commissure intersegmentaire nécessaire ici).
    \end{itemize}
    \item \textbf{Schéma de l'arc réflexe (Exp 1 - Retrait cuisse) :}
    \begin{popfigure}
    \begin{tikzpicture}[>=Stealth, node distance=1.2cm and 1.5cm, font=\small]
        % Nodes
        \node[draw, circle, fill=popblueL, minimum size=0.8cm, align=center] (Recep){Récepteur\\cutané};
        \node[draw, rectangle, rounded corners, fill=popgreenL, right=of Recep, minimum width=2.5cm, align=center] (NS){Neurone Sensitif\\ (Pseudo-unipolaire)\\(Ganglion rachidien)};
        \node[draw, rectangle, rounded corners, fill=poporangeL, below right=of NS, minimum width=2.5cm, align=center] (IN){Interneurone\\ (Corne grise postérieure)};
        \node[draw, rectangle, rounded corners, fill=poppinkL, right=of IN, minimum width=2.5cm, align=center] (MN){Motoneurone $\alpha$\\ (Corne grise antérieure)};
        \node[draw, ellipse, fill=popgoldL, right=of MN, minimum width=2cm, align=center] (Eff){Effecteur\\ Muscle squelettique\\ (Fléchisseur)};
        
        % Arrows
        \draw[->, thick, popblue] (Recep) -- node[above, sloped] {Influx nerveux} (NS);
        \draw[->, thick, popblue] (NS) -- node[above] {Racine dorsale} (IN);
        \draw[->, thick, popblue] (IN) -- node[above] {Synapse excitatrice} (MN);
        \draw[->, thick, popblue] (MN) -- node[above, align=center]{Racine ventrale \\ Nerf sciatique} (Eff);
        
        % Labels for roots
        \node[above=0.2cm of NS, text=popdark, font=\tiny\bfseries] {Racine Dorsale (Sensitive)};
        \node[below=0.2cm of MN, text=popdark, font=\tiny\bfseries] {Racine Ventrale (Motrice)};
    \end{tikzpicture}
    \legende{Schéma de l'arc réflexe médullaire polysynaptique (un interneurone) pour le réflexe de retrait (flexion) de la jambe gauche. 1. Récepteur nocicepteur cutané. 2. Neurone sensitif (corps cellulaire dans le ganglion rachidien). 3. Interneurone exciteur (corne postérieure). 4. Motoneurone alpha (corne antérieure). 5. Effecteur : muscle fléchisseur (ischio-jambiers). Flèches bleues : sens de propagation de l'influx nerveux.}
    \end{popfigure}
\end{enumerate}
\textbf{Conclusion :} Les expériences de Pflüger démontrent que la moelle épinière possède une \textbf{autonomie fonctionnelle} : elle intègre l'influx sensitif et élabore une réponse motrice adaptée (loi du siège) sans intervention de l'encéphale, mais cette capacité est modulable par l'état physiologique du centre (fatigue, facilitation/inhibition descendante).
\end{exoresolu}

\begin{methode}[Interpréter les expériences de Bell et Magendie (Racines rachidiennes)]
\textbf{Objectif :} Démontrer la séparation fonctionnelle des racines rachidiennes (Loi de Bell-Magendie).
\begin{enumerate}
    \item \textbf{Contexte historique :} François Magendie (1822) et Charles Bell (1811/1826) sur le lapin/grenouille.
    \item \textbf{Expérience type (Magendie) :}
    \begin{itemize}
        \item \textbf{Groupe Témoin :} Lapin intact. Pinçage patte antérieure $\rightarrow$ Retrait + Cri (douleur). Section nerf vague (contrôle) $\rightarrow$ Pas d'effet sur membre.
        \item \textbf{Groupe 1 : Section Racine Dorsale (Ganglion) d'un nerf spinal (ex. sciatique).}
        \item \textbf{Observation :} Pinçage zone innervée $\rightarrow$ \textbf{Aucun retrait, Aucune réaction douloureuse (analgésie)}. Mais le muscle reste contractile (stimulation électrique directe du nerf ventral ou muscle $\rightarrow$ contraction).
        \item \textbf{Conclusion :} Racine Dorsale = \textbf{Voie sensitive (afferente)}. Section $\rightarrow$ Perte de la sensibilité (anesthésie/analgésie) du dermatome correspondant.
        \item \textbf{Groupe 2 : Section Racine Ventrale du même nerf.}
        \item \textbf{Observation :} Pinçage zone innervée $\rightarrow$ \textbf{Réaction douloureuse (cri, retrait controlatéral), Sensibilité conservée}. Mais \textbf{Paralysie flasque} du myotome correspondant (atrophie ultérieure). Stimulation électrique racine ventrale proximale $\rightarrow$ Pas de contraction. Stimulation muscle direct $\rightarrow$ Contraction.
        \item \textbf{Conclusion :} Racine Ventrale = \textbf{Voie motrice (efférente)}. Section $\rightarrow$ Paralysie flasque, aréflexie, atrophie du myotome.
    \end{itemize}
    \item \textbf{Synthèse :} \textbf{Loi de Bell-Magendie :} « Les racines dorsales sont sensitives, les racines ventrales sont motrices. » Le nerf spinal (mixte) = réunion des deux racines hors du rachis.
    \item \textbf{Application clinique (Cameroun) :} Traumatisme rachidien (accident moto-taxi « bend skin » à Douala/Yaoundé) : Syndrome radiculaire (douleur radiculaire = atteinte racine dorsale ; paralysie/amyo = atteinte racine ventrale). Hernie discale L4-L5 $\rightarrow$ compression racine L5 (dorsale = douleur sciatique ; ventrale = pied bot/steppage).
\end{enumerate}
\end{methode}

\begin{exoresolu}[Analyse d'un syndrome radiculaire suite à hernie discale]
\textbf{Énoncé :} Un patient de 35 ans, cultivateur à Dschang, consulte pour une lombalgie irradiant dans la cuisse et la jambe droites depuis 3 semaines. L'examen clinique révèle :
\begin{itemize}
    \item Douleur vive à la percussion des apophyses épineuses L4-L5.
    \item Douleur projetée sur le trajet du nerf sciatique droit (Lasègue positif à 30°).
    \item Hypoesthésie (diminution sensibilité) sur la face latérale de la jambe et le dos du pied droit.
    \item Force musculaire : Extension du genou (quadriceps) = 5/5. Flexion plantaire (triceps sural) = 3/5. Dorsiflexion pied (tibial antérieur) = 3/5. Extension gros orteil (exhalteur) = 2/5.
    \item Réflexe rotulien (L3-L4) normal. Réflexe achilléen (S1) diminué/aboli à droite.
    \item IRM : Hernie discale L4-L5 postéro-latérale droite comprimant la racine L5.
\end{itemize}
Questions :
\begin{enumerate}
    \item Rappeler la loi de Bell-Magendie et préciser quelle(s) racine(s) est/sont comprimées par la hernie L4-L5 droite.
    \item Expliquer les signes cliniques (sensitifs et moteurs) par la compression de cette racine.
    \item Pourquoi le réflexe rotulien est-il normal alors que le réflexe achilléen est aboli ?
    \item Préciser le diagnostic topographique : Niveau radiculaire et nerf atteint.
\end{enumerate}

\tcblower
\textbf{Solution détaillée :}
\begin{enumerate}
    \item \textbf{Loi de Bell-Magendie :} Les \textbf{racines dorsales} (postérieures) sont \textbf{sensitives} (afferentes), les \textbf{racines ventrales} (antérieures) sont \textbf{motrices} (efférentes). Le ganglion rachidien (sur la racine dorsale) contient les corps cellulaires des neurones sensitifs pseudo-unipolaires.
    \textbf{Règle de numérotation (Rachis lombaire) :} La racine sort \textbf{sous} la vertèbre de même nom (ex: racine L5 sort par le foramen L5-S1). Une hernie discale postéro-latérale comprime la racine \textbf{traversante} qui descend vers son foramen de sortie.
    \textbf{Application :} Hernie discale L4-L5 $\rightarrow$ Comprime la racine \textbf{L5} (traversante). Hernie L5-S1 $\rightarrow$ Comprime racine \textbf{S1}.
    \textit{Donc ici : Compression de la racine \textbf{L5 droite} (racine dorsale + racine ventrale).}
    \item \textbf{Explication des signes :}
    \begin{itemize}
        \item \textbf{Sensitifs (Racine Dorsale L5) :} Hypoesthésie territoire L5 (face latérale jambe, dos du pied, espace 1er-2ème orteil). Douleur projetée (névralgie sciatique L5) : l'influx nociceptif anormal dans la racine comprimée est interprété par le cortex comme venant du territoire cutané de cette racine (loi de la projection).
        \item \textbf{Moteurs (Racine Ventrale L5) :} Parésie (force 2-3/5) des muscles innervés par L5 : Tibial antérieur (dorsiflexion), Exhalteur du gros orteil, Gluteus medius (marche de Trendelenburg probable). Quadriceps (L3-L4) intact $\rightarrow$ Force 5/5.
    \end{itemize}
    \item \textbf{Réflexes :}
    \begin{itemize}
        \item \textbf{Rotulien (L3-L4) :} Arc réflexe L3-L4 (Quadriceps). Racines L3, L4 non comprimées $\rightarrow$ \textbf{Normal}.
        \item \textbf{Achilléen (S1-S2) :} Arc réflexe S1 (Triceps sural). L'abolition du réflexe achilléen signe classiquement une atteinte \textbf{S1}. Ici, force flexion plantaire 3/5 (S1) + réflexe aboli $\rightarrow$ \textbf{Atteinte associée S1} (hernie volumineuse ou séquestre migrant caudalement comprimant S1 traversant). \textit{Au Bac, on lie le signe au niveau : Réflexe achilléen aboli = Atteinte S1.}
    \end{itemize}
    \item \textbf{Diagnostic topographique :} \textbf{Syndrome radiculaire L5 droit} (associé à une atteinte S1 droite vu le réflexe achilléen). Nerf sciatique droit (racines L4-S3) atteint à sa racine L5 (et S1). Cause : Hernie discale L4-L5 postéro-latérale droite.
\end{enumerate}
\textbf{Conclusion :} La clinique du syndrome radiculaire est l'application directe de la loi de Bell-Magendie : compression racine dorsale $\rightarrow$ troubles sensitifs (douleur, hypoesthésie) ; compression racine ventrale $\rightarrow$ troubles moteurs (parésie, aréflexie, amyotrophie). L'IRM confirme le conflit disco-radiculaire.
\end{exoresolu}

\begin{methode}[Mettre en évidence le réflexe myotatique (rotulien / achilléen) et schématiser son arc]
\textbf{Objectif :} Démontrer le réflexe d'extension (monosynaptique) et son utilité clinique.
\begin{enumerate}
    \item \textbf{Matériel :} Marteau à réflexes (Taylor ou Babinski), sujet sain (élève), table d'examen.
    \item \textbf{Réflexe Rotulien (L2-L3-L4, prédominance L4) :}
    \begin{itemize}
        \item Sujet assis, cuisses horizontales, jambes pendantes, détendu (distraction : compter à rebours, serrer les mains - manœuvre de Jendrassik).
        \item Percussion brutale du \textbf{tendon rotulien} (entre rotule et tubérosité tibiale antérieure).
        \item \textbf{Réponse normale :} Extension brève de la jambe (contraction quadriceps), rotation interne légère du pied. Notation : +++ (vif), ++ (normal), + (peu vif), 0 (aboli).
    \end{itemize}
    \item \textbf{Réflexe Achilléen (S1-S2) :}
    \begin{itemize}
        \item Sujet assis ou couché, genou fléchi, hanche en rotation externe (pied en équin).
        \item Percussion du \textbf{tendon d'Achille} (calcanéen).
        \item \textbf{Réponse normale :} Flexion plantaire brève du pied (contraction triceps sural). Notation idem.
    \end{itemize}
    \item \textbf{Schéma de l'arc réflexe myotatique (Monosynaptique pour l'agoniste) :}
    \begin{enumerate}
        \item \textbf{Récepteur :} \textbf{Fuseau neuromusculaire} (fibres intrafusales : nuclear bag/chain) dans le muscle (quadriceps/triceps). Sensible à l'\textbf{étirement} (vitesse et amplitude).
        \item \textbf{Neurone Sensitif (Ia) :} Corps cellulaire dans ganglion rachidien. Fibre myélinisée gros diamètre (vitesse \SI{80-120}{m/s}). Pénètre par racine dorsale.
        \item \textbf{Synapse Monosynaptique :} Terminaison du neurone Ia $\rightarrow$ directement sur \textbf{Motoneurone $\alpha$} (corne antérieure) du \textbf{même muscle} (agoniste). Glutamate $\rightarrow$ Récepteurs AMPA. \textbf{Pas d'interneurone}.
        \item \textbf{Motoneurone $\alpha$ :} Axone sort par racine ventrale $\rightarrow$ Nerf fémoral (rotulien) / sciatique-tibial (achilléen) $\rightarrow$ Plaque motrice fibres extrafusales (muscle agoniste) $\rightarrow$ \textbf{Contraction (raccourcissement)}.
        \item \textbf{Inhibition Réciproque (Via Interneurone Inhibiteur) :} Collatérale du neurone Ia $\rightarrow$ Interneurone inhibiteur (Glycine/GABA) $\rightarrow$ Motoneurone $\alpha$ du \textbf{muscle antagoniste} (ischio-jambiers / tibial antérieur) $\rightarrow$ \textbf{Relâchement}. Permet le mouvement non freiné.
        \item \textbf{Régulation $\gamma$ :} Motoneurones $\gamma$ $\rightarrow$ Fibres intrafusales (pôles) $\rightarrow$ Maintien tension fuseau pendant contraction (co-activation $\alpha$-$\gamma$) $\rightarrow$ Sensibilité maintenue.
    \end{enumerate}
    \item \textbf{Importance :} Régulation automatique de la \textbf{tonicité musculaire} (posture station debout), adaptation aux charges imprévues, base de la marche.
\end{enumerate}
\end{methode}

\begin{exoresolu}[Analyse de l'arc réflexe myotatique et conséquence d'une lésion]
\textbf{Énoncé :} On étudie le réflexe rotulien chez un sujet sain.
\begin{enumerate}
    \item Schématiser l'arc réflexe myotatique complet (étirement $\rightarrow$ contraction quadriceps + relâchement ischio-jambiers). Légender : Récepteur, Neurone Ia, Motoneurone $\alpha$ agoniste, Interneurone inhibiteur, Motoneurone $\alpha$ antagoniste, Motoneurone $\gamma$.
    \item Expliquer pourquoi ce réflexe est dit \textbf{monosynaptique} alors qu'il fait intervenir une synapse inhibitrice pour l'antagoniste.
    \item Un patient présente une \textbf{section complète de la moelle épinière au niveau T10} (paraplégie). Prédire l'état des réflexes rotuliens (L3-L4) et achilléens (S1) \textbf{immédiatement après le choc} (phase de choc spinal) et \textbf{à distance} (phase chronique). Justifier.
    \item Une piqûre d'acide acétique sur la cuisse provoque une flexion de la jambe (réflexe de retrait). Pourquoi ce réflexe est-il \textbf{polysynaptique} et plus lent que le réflexe myotatique ?
\end{enumerate}

\tcblower
\textbf{Solution détaillée :}
\begin{enumerate}
    \item \textbf{Schéma de l'arc réflexe myotatique (Rotulien) :}
    \begin{popfigure}
    \begin{tikzpicture}[>=Stealth, node distance=1.1cm and 1.3cm, font=\small, every node/.style={align=center}]
        % Muscle Spindle
        \node[draw, ellipse, fill=popgreenL, minimum width=2.2cm, align=center] (Fuseau){Fuseau neuromusculaire\\ (Récepteur étirement)\\(Quadriceps)};
        % Ia Afferent
        \node[draw, circle, fill=popblueL, right=of Fuseau, minimum size=1cm, align=center] (Ia){Neurone Ia\\ (Sensitif)\\(Ganglion L3-L4)};
        % Alpha MN Agonist
        \node[draw, rectangle, rounded corners, fill=poporangeL, below right=of Ia, minimum width=2.2cm, align=center] (AlphaAg){Motoneurone $\alpha$\\ (Corne antérieure L3-L4)\\(Quadriceps - Agoniste)};
        % Gamma MN
        \node[draw, rectangle, rounded corners, fill=poppurpleL, above right=of Ia, minimum width=2.2cm, align=center] (Gamma){Motoneurone $\gamma$\\ (Corne antérieure)\\(Fuseau - Coactivation)};
        % Inhibitory Interneuron
        \node[draw, diamond, fill=poppinkL, right=of AlphaAg, minimum width=2.2cm, align=center] (InhInt){Interneurone\\ Inhibiteur (Glycine/GABA)};
        % Alpha MN Antagonist
        \node[draw, rectangle, rounded corners, fill=popblueL, below right=of InhInt, minimum width=2.2cm, align=center] (AlphaAnt){Motoneurone $\alpha$\\ (Ischio-jambiers - Antagoniste)};
        % Effectors
        \node[draw, ellipse, fill=popgoldL, right=of AlphaAg, minimum width=2cm, align=center] (EffAg){Contraction\\ Quadriceps\\ (Extension jambe)};
        \node[draw, ellipse, fill=popgoldL, right=of AlphaAnt, minimum width=2cm, align=center] (EffAnt){Relâchement\\ Ischio-jambiers};
        
        % Connections
        \draw[->, thick, popblue] (Fuseau) -- node[above] {Étirement} (Ia);
        \draw[->, thick, popblue] (Ia) -- node[above left] {Racine dorsale} (AlphaAg);
        \draw[->, thick, popblue, dashed] (Ia) -- node[above] {Collatérale} (InhInt);
        \draw[->, thick, popblue] (AlphaAg) -- node[above] {Racine ventrale / Nerf fémoral} (EffAg);
        \draw[->, thick, popred] (InhInt) -- node[above] {Inhibition (IPSP)} (AlphaAnt);
        \draw[->, thick, popred] (AlphaAnt) -- node[above] {Racine ventrale / Nerf sciatique} (EffAnt);
        \draw[->, thick, poppurple] (Gamma) -- node[above] {Innervation pôles fuseau} (Fuseau);
        
        % Labels
        \node[above=0.1cm of Fuseau, font=\tiny\bfseries, text=popdark] {1. Récepteur};
        \node[above=0.1cm of Ia, font=\tiny\bfseries, text=popdark] {2. Neurone Ia};
        \node[below=0.1cm of AlphaAg, font=\tiny\bfseries, text=popdark] {3. MN $\alpha$ Agoniste};
        \node[above=0.1cm of Gamma, font=\tiny\bfseries, text=popdark] {4. MN $\gamma$};
        \node[above=0.1cm of InhInt, font=\tiny\bfseries, text=popdark] {5. Interneurone Inhibiteur};
        \node[below=0.1cm of AlphaAnt, font=\tiny\bfseries, text=popdark] {6. MN $\alpha$ Antagoniste};
    \end{tikzpicture}
    \legende{Arc réflexe myotatique rotulien (monosynaptique pour la voie excitatrice). Flèches bleues : voie excitatrice (Glutamate). Flèches rouges : voie inhibitrice réciproque (Glycine/GABA). Flèche violette : boucle de servocontrôle $\gamma$ (maintien sensibilité fuseau).}
    \end{popfigure}
    \item \textbf{Monosynaptique ?} Le terme « monosynaptique » qualifie \textbf{uniquement la voie excitatrice directe} entre le neurone sensitif Ia et le motoneurone $\alpha$ de l'agoniste (\textbf{une seule synapse chimique}). L'inhibition réciproque de l'antagoniste fait obligatoirement intervenir un \textbf{interneurone inhibiteur} (voie disynaptique : Ia $\rightarrow$ Interneurone $\rightarrow$ MN $\alpha$ antagoniste). C'est un arc \textbf{monosynaptique pour l'agoniste, disynaptique pour l'antagoniste}.
    \item \textbf{Section médullaire T10 (Paraplégie) :}
    \begin{center}
    \begin{tabularx}{\linewidth}{|l|X|X|}\hline
    \rowcolor{popblueL} \textbf{Phase} & \textbf{Réflexes Rotuliens (L3-L4) / Achilléens (S1)} & \textbf{Justification} \\ \hline
    \textbf{Immédiate (Choc spinal)} & \textbf{Abolis (0/4)} \textbf{Aréflexie totale} & Brutale privation des influences \textbf{descendantes facilitrices} (voies réticulospinales, vestibulospinales, corticospinales) qui maintiennent l'excitabilité basale des motoneurones $\alpha$ et $\gamma$. Les motoneurones sont en \textbf{dépolarisation post-synaptique inhibitoire} prolongée / inexcitables. Perte du tonus $\rightarrow$ \textbf{paralysie flasque}. \\ \hline
    \textbf{Chronique (Semaines/Mois)} & \textbf{Hyper-réfléchie (+++/4) + Clonus + Spasticité} & \textbf{Désensibilisation par déafferentation (Supersensibilité de Cannon)} : Prolifération de récepteurs post-synaptiques (dénervation supersensitivity) sur les motoneurones isolés. Perte du contrôle inhibiteur descendant (voies corticospinales, rubrospinales). \textbf{Hyperréflexie}, \textbf{Clonus} (oscillations réflexes rapides), \textbf{Spasticité} (hypertonie à composante vitesse-dépendante), \textbf{Signe de Babinski} (extension gros orteil). \\ \hline
    \end{tabularx}
    \end{center}
    \item \textbf{Réflexe de retrait (Polysynaptique) :}
    \begin{itemize}
        \item Récepteur : Nocicepteur (chimique, thermique, mécanique intense) $\rightarrow$ Fibres A$\delta$ (vite) et C (lent).
        \item Voie : Neurone sensitif $\rightarrow$ \textbf{Plusieurs interneurones} (excitateurs pour fléchisseurs, inhibiteurs pour extenseurs, interneurones commissuraux pour membre controlatéral $\rightarrow$ réflexe croisé d'extension) $\rightarrow$ Motoneurones $\alpha$.
        \item \textbf{Nombre de synapses $\ge$ 2 (souvent 3-4)} $\rightarrow$ Temps de latence synaptic cumulé ($\approx$ \SI{0,5-1}{ms} par synapse) $\rightarrow$ \textbf{Latence totale plus longue} (\SI{30-50}{ms} vs \SI{15-25}{ms} pour myotatique).
        \item \textbf{Durée plus longue} (réverbération dans les circuits interneuronaux) $\rightarrow$ Protection/Retrait soutenu.
    \end{itemize}
\end{enumerate}
\textbf{Conclusion :} Le réflexe myotatique est le prototype du réflexe monosynaptique rapide, assurant la régulation instantanée de la longueur musculaire. La section médullaire révèle la dépendance des centres médullaires vis-à-vis du contrôle supraspinal (tonus, inhibition) : le choc spinal (aréflexie) laisse place à l'hyperréflexie spastique par plasticité synaptique maladaptative.
\end{exoresolu}

\begin{methode}[Expliquer les mécanismes des réflexes acquis (conditionnels) - Pavlov]
\textbf{Objectif :} Comprendre l'apprentissage associatif simple (conditionnement classique) et sa base neuronale.
\begin{enumerate}
    \item \textbf{Définition :} Réflexe \textbf{acquis} (non inné), \textbf{conditionnel} (dépend d'une association temporelle), \textbf{transitoire} (s'éteint sans renforcement), \textbf{conscient} (chez l'homme).
    \item \textbf{Paradigme de Pavlov (Chien) :}
    \begin{itemize}
        \item \textbf{Phase 1 (Inné) :} Stimulus Inconditionnel (SI : Viande) $\rightarrow$ Réponse Inconditionnelle (RI : Salivation). Arc inné, génétiquement programmé.
        \item \textbf{Phase 2 (Conditionnement) :} Stimulus Conditionnel (SC : Sonnerie/Metronome) \textbf{PRÉCÈDE} SI (Viande) de \SI{0,5}{s} à \SI{2}{s} (contiguïté temporelle). Répétition (20-50 essais).
        \item \textbf{Phase 3 (Acquis) :} SC (Sonnerie) \textbf{SEUL} $\rightarrow$ Réponse Conditionnelle (RC : Salivation). RC $\neq$ RI (quantité, composition, latence différents).
    \end{itemize}
    \item \textbf{Règles fondamentales :}
    \begin{itemize}
        \item \textbf{Contiguïté temporelle :} SC avant SI (conditionnement en avant / \textit{forward}). Optimal : délai court (\SI{0,5}{s}). Conditionnement simultané ou arrière (SC après SI) inefficace.
        \item \textbf{Répétition / Renforcement :} Association SC-SI doit être répétée. Présentation SC seul $\rightarrow$ \textbf{Extinction} (inhibition interne).
        \item \textbf{Généralisation :} Stimuli semblables au SC provoquent RC (gradient de généralisation).
        \item \textbf{Discrimination :} Apprentissage différentiel (SC+ renforcé, SC- non renforcé).
    \end{itemize}
    \item \textbf{Bases neurales (Modèle actuel) :}
    \begin{itemize}
        \item \textbf{Convergence :} Voie SC (ex. auditif : Oreille $\rightarrow$ Tronc $\rightarrow$ Thalamus $\rightarrow$ Cortex Auditif $\rightarrow$ \textbf{Hippocampe / Cortex préfrontal / Amygdale}) et Voie SI (Gustatif $\rightarrow$ Nucleus tractus solitarius $\rightarrow$ Thalamus $\rightarrow$ Cortex Insulaire) convergent vers des \textbf{neurones détecteurs de coïncidence} (ex. cortex associatif, hippocampe, cervelet pour conditionnement paupière).
        \item \textbf{Plasticité synaptique (LTP - Potentiation Long Terme) :} « \textit{Neurons that fire together, wire together} » (Hebbe). Activation simultanée entrée faible (SC) + entrée forte (SI) $\rightarrow$ Dépolarisation forte du post-synaptique $\rightarrow$ Déblocage récepteurs NMDA $\rightarrow$ Influx $Ca^{2+}$ $\rightarrow$ Cascade kinases $\rightarrow$ Insertion récepteurs AMPA $\rightarrow$ **Potentiation synapse SC $\rightarrow$ Neurone sortie**. Le SC seul suffit désormais à déclencher la sortie.
        \item \textbf{Stockage :} Cortex associatif (temporel, préfrontal), Hippocampe (déclaratif/contexte), Cervelet (moteur simple), Amygdale (émotionnel/peur).
    \end{itemize}
    \item \textbf{Importance :} Adaptation fine à l'environnement changeant, anticipation (préparation physiologique : salivation, sécrétion gastrique, insulinémie pré-prandiale), base de l'éducation, des phobies, des addictions (conditionnement opérant = Skinner, différent).
\end{enumerate}
\end{methode}

\begin{exoresolu}[Modélisation du conditionnement classique et extinction]
\textbf{Énoncé :} On réalise une expérience de conditionnement pavlovien sur un lapin.
\begin{itemize}
    \item \textbf{SI :} Souffle d'air sur la cornée (réflexe inné de clignement paupière - RI).
    \item \textbf{SC :} Son pure (1000 Hz, \SI{80}{dB}) présenté \SI{250}{ms} avant le SI.
    \item \textbf{Protocole :} 100 essais d'acquisition (SC+SI), puis 50 essais d'extinction (SC seul), puis repos 24h, puis 10 essais de récupération spontanée (SC seul).
\end{itemize}
Résultats (Pourcentage de Réponses Conditionnelles - RC) :
\begin{center}
\begin{tabular}{|c|c|c|c|c|c|}\hline
\textbf{Bloc d'essais} & 1-10 & 21-30 & 51-60 & 91-100 & \textbf{Extinction 1-10} \\ \hline
\textbf{\% RC} & 10\% & 55\% & 85\% & 95\% & \textbf{40\%} \\ \hline
\end{tabular}
\end{center}
\begin{center}
\begin{tabular}{|c|c|c|}\hline
\textbf{Phase} & \textbf{Récup. spontanée (Essai 1)} & \textbf{Récup. spontanée (Essai 10)} \\ \hline
\textbf{\% RC} & \textbf{60\%} & \textbf{10\%} \\ \hline
\end{tabular}
\end{center}
Questions :
\begin{enumerate}
    \item Identifier SI, RI, SC, RC. Préciser le type de conditionnement (avance, trace, simultané, arrière).
    \item Interpréter la courbe d'acquisition (blocs 1-100). Quel phénomène neurophysiologique sous-tend l'augmentation du \% RC ?
    \item Expliquer la chute du \% RC pendant l'extinction (blocs 1-10 extinction). Est-ce un oubli ?
    \item Interpréter la récupération spontanée (60\% puis 10\%). Quelle conclusion tirer sur la nature de l'extinction ?
    \item Si on présentait le SI \textbf{avant} le SC (conditionnement arrière), quel serait le résultat ? Justifier.
\end{enumerate}

\tcblower
\textbf{Solution détaillée :}
\begin{enumerate}
    \item \textbf{Identification :}
    \begin{itemize}
        \item SI (Stimulus Inconditionnel) : \textbf{Souffle d'air cornéen}.
        \item RI (Réponse Inconditionnelle) : \textbf{Clignement paupière (réflexe de protection)}.
        \item SC (Stimulus Conditionnel) : \textbf{Son 1000 Hz}.
        \item RC (Réponse Conditionnelle) : \textbf{Clignement paupière anticipé (déclenché par le son)}.
    \end{itemize}
    \textbf{Type : Conditionnement \textbf{en avance (delay conditioning)}.} Le SC commence \SI{250}{ms} avant le SI et \textbf{chevauche} le SI (le son continue pendant le souffle d'air). C'est le paradigme le plus efficace. (Si SC s'arrêtait avant SI = Conditionnement \textit{trace} ; si SC après SI = \textit{backward} inefficace).
    \item \textbf{Acquisition :} Croissance sigmoïde du \% RC (10\% $\rightarrow$ 95\%).
    \begin{itemize}
        \item Phase latente (1-10) : Détection contingence.
        \item Phase exponentielle (20-60) : Renforcement synaptique rapide.
        \item Plateau asymptotique (90-100) : Performance maximale (plafond biologique, bruit).
    \end{itemize}
    \textbf{Mécanisme : Potentiation à Long Terme (LTP) hippocampique/cérébelleuse.} Co-activation entrée SC (faible, glutamate $\rightarrow$ AMPA) et entrée SI (forte, despolarisation forte $\rightarrow$ déblocage Mg$^{2+}$ récepteur NMDA $\rightarrow$ influx $Ca^{2+}$) $\rightarrow$ Cascade $CaMKII$ $\rightarrow$ Phosphorylation/Insertion récepteurs AMPA supplémentaires sur synapse SC $\rightarrow$ **Augmentation efficacité synaptique SC $\rightarrow$ Neurone sortie (Noyau interposé cervelet / Motoneurone facial)**. Le SC seul suffit à dépolariser au seuil.
    \item \textbf{Extinction :} Présentation SC seul (sans SI). Le SC active la synapse potentiée $\rightarrow$ Déclenche RC. Mais \textbf{absence de renforcement (SI)} $\rightarrow$ Pas de despolarisation forte conjointe $\rightarrow$ \textbf{Dépression à Long Terme (LTD) / Inhibition synaptique active} (interneurones inhibiteurs recrutés par le contexte « non-renforcé », endocannabinoïdes, phosphatases PP1/PP2A déphosphorylant AMPA). \textbf{Non, ce n'est pas un oubli (effacement de la trace)}, c'est un \textbf{nouvel apprentissage inhibiteur} (« SC signifie maintenant : PAS de SI ») qui masque la trace excitatory originale.
    \item \textbf{Récupération spontanée :} Après 24h repos, 1er essai SC seul $\rightarrow$ 60\% RC (vs 40\% fin extinction). L'inhibition d'extinction (dépendante du contexte temporel, fragile) a \textbf{décru spontanément} (oubli de l'inhibition), révélant la \textbf{trace excitatory originale (LTP) intacte}. Essais suivants (10) $\rightarrow$ Ré-extinction rapide (10\%) : l'inhibition se réinstalle vite (réactivation du circuit inhibitif). \textbf{Preuve :} L'extinction n'efface pas la mémoire, elle la masque par une mémoire inhibitrice concurrente.
    \item \textbf{Conditionnement arrière (SI puis SC) :} \textbf{Aucun conditionnement (RC $\approx$ 0\%).} Justification : Le SC arrive \textbf{après} la réponse inconditionnelle et la despolarisation forte du SI. Pas de coïncidence temporelle « SC avant SI » nécessaire pour la règle de Hebb (LTP). Le SC ne prédit pas le SI. Au contraire, le SC peut devenir un \textbf{inhibiteur conditionnel} (signal de « fin de SI » / sécurité).
\end{enumerate}
\textbf{Conclusion :} Le réflexe conditionnel est un modèle d'apprentissage associatif reposant sur la \textbf{plasticité synaptique (LTP/LTD)} dans des circuits de convergence (cervelet pour clignement, hippocampe/cortex pour complexe). L'extinction est un apprentissage inhibiteur distinct, réversible, et non un effacement de la trace mnésique initiale.
\end{exoresolu}

\begin{methode}[Expliquer les cas de paralysie ou d'absence de réflexe (Sémiologie neurologique)]
\textbf{Objectif :} Différencier les syndromes neuronaux (central vs périphérique) et localiser la lésion sur l'arc réflexe.
\begin{enumerate}
    \item \textbf{Analyse du réflexe ostéotendineux (ROT) :} Recherche systématique (Rotulien L3-L4, Achilléen S1, Bicipital C5-C6, Stylo-radien C6-C7, Tricipital C7-C8).
    \item \textbf{Notation :} 0 (aboli), + (peu vif), ++ (normal), +++ (vif), ++++ (clonus).
    \item \textbf{Syndrome Moteur Central (Pyramidal / Sus-segmentaire) :} Lésion voie pyramidale (Cortex, Capsule interne, Tronc, Moelle \textbf{au-dessus} du métamère).
    \begin{itemize}
        \item \textbf{Hyperréflexie (+++/++++)} : Perte de l'inhibition descendante (corticospinale, rubrospinal) sur les interneurones inhibiteurs (Renshaw) et motoneurones $\alpha$ $\rightarrow$ Hyperexcitabilité du motoneurone.
        \item \textbf{Spasticité} : Hypertonie à composante vitesse-dépendante (étirement rapide $\rightarrow$ résistance brutale « lame de couteau »).
        \item \textbf{Signe de Babinski (Extension gros orteil + éventail orteils)} : Signe de déconnexion cortico-spinale (normal < 2 ans : myélinisation incomplète).
        \item \textbf{Paralysie spastique} : Force diminuée mais tonus augmenté. Pas d'amyotrophie précoce (dénervation centrale $\neq$ dénervation périphérique).
        \item \textbf{Réflexes cutanés plantaires :} Étendus (réflexe de Babinski positif).
    \end{itemize}
    \item \textbf{Syndrome Moteur Périphérique (Inférieur / Segmentaire) :} Lésion Unité Motrice Finale (Motoneurone $\alpha$, Racine ventrale, Nerf, Jonction neuromusculaire, Muscle).
    \begin{itemize}
        \item \textbf{Aréflexie (0)} : Interruption de l'arc réflexe (afference Ia, interneurone, efference $\alpha$).
        \item \textbf{Paralysie flasque} : Tonus aboli (hypotonie/atonie).
        \item \textbf{Amyotrophie rapide} : Dénervation $\rightarrow$ Atrophie par désinnervation (fibrillations, fasciculations visibles).
        \item \textbf{Fasciculations} : Contractions spontanées unités motrices dénervées (hyperexcitabilité membrane motoneurone lésé).
        \item \textbf{Réflexes cutanés :} Abolis dans le territoire.
    \end{itemize}
    \item \textbf{Localisation topographique (Exemples Cameroun) :}
    \begin{itemize}
        \item \textbf{Hémiplégie spastique controlatérale + Babinski +} : AVC (Accident Vasculaire Cérébral) hémisphère controlatéral (HTA non traitée, hôpital Laquintinie Douala).
        \item \textbf{Paraplégie spastique + Niveau sensitif + Sphincters} : Section/Compression moelle (Traumatisme rachidien, Tumeur intramédullaire, Pott - TB vertébrale).
        \item \textbf{Paralysie flasque monomélique + Aréflexie + Amyotrophie} : Lésion nerf périphérique (Traumatisme, Compression tunnel carpien, Syndrome du pronateur) ou Racine (Hernie discale, Zona).
        \item \textbf{Paralysie flasque généralisée aréflexique} : Syndrome de Guillain-Barré (Polyradiculonévrite démyélinisante aiguë post-infectieuse - Campylobacter, Paludisme, VIH), Botulisme, Myasthénie (jonction), Poliomyélite (virus sauvage / VDPV - surveillance OMS Cameroun).
    \end{itemize}
\end{enumerate}
\end{methode}

\begin{exoresolu}[Diagnostic différentiel de paralysies : Cas cliniques]
\textbf{Énoncé :} Trois patients admis aux urgences de l'Hôpital Central de Yaoundé.
\begin{description}
    \item[\textbf{Cas 1 :}] Homme 55 ans, HTA connue, non observant. Brutalement : hémiparésie droite dense (face + bras + jambe), aphasie de Broca. Réflexes rotulien/bicipital droits +++, Babinski droit +. Sensibilité altérée droite.
    \item[\textbf{Cas 2 :}] Jeune 22 ans, chute de moto (bend skin) il y a 1h. Paraplégie flasque des 2 membres inférieurs, aréflexie rotulienne/achilléenne bilatérale, anesthésie complète sous T6, rétention vésicale, priapisme. Radiographie : fracture-dislocation T10-T11.
    \item[\textbf{Cas 3 :}] Femme 30 ans, 10 jours post-partum. Fatigue, paresthésies pieds $\rightarrow$ montée progressive (jambes, cuisses, bras, visage) en 3 jours. Tétraparésie flasque aréflexique générale, dysphagie, ophtalmoplégie. LCR : hyperprotéinorachie (\SI{1,2}{g/L}) sans cytose (dissociation cyto-protéinique). Antécédent diarrhée 2 sem. avant.
\end{description}
Questions :
\begin{enumerate}
    \item Pour chaque cas, préciser : Type de syndrome moteur (Central / Périphérique), Niveau de la lésion (Topographie), Étio-diagnostique probable.
    \item Expliquer le mécanisme physiopathologique de l'\textbf{aréflexie} dans le Cas 2 (phase aiguë) vs l'\textbf{hyperréflexie} dans le Cas 1.
    \item Pour le Cas 3, expliquer la dissociation cyto-protéinique du LCR et le mécanisme de la remontée des symptômes (céphalisation).
    \item Pourquoi le réflexe rotulien est-il aboli dans le Cas 2 alors que la lésion est à T10 (bien au-dessus de L3-L4) ?
\end{enumerate}

\tcblower
\textbf{Solution détaillée :}
\begin{enumerate}
    \item \textbf{Bilan diagnostique :}
    \begin{center}
    \begin{tabularx}{\linewidth}{|l|X|X|X|}\hline
    \rowcolor{popblueL} \textbf{Cas} & \textbf{Syndrome / Niveau} & \textbf{Topographie Lésion} & \textbf{Étiologie Probable} \\ \hline
    \textbf{1} & \textbf{Syndrome Pyramidal (Central) Hémicorporel Droit} + Aphasie (Hémisphère Gauche dominant) & \textbf{Hémisphère Cérébral Gauche} (Cortex frontal/moteur + aire de Broca + Voie pyramidale descendante + Thalamus/Capsule interne) & \textbf{AVC Ischémique (Infarctus artère cérébrale moyenne Gauche) ou Hémorragie intracérébrale} (HTA) \\ \hline
    \textbf{2} & \textbf{Syndrome Médullaire Complet (Transsection) + Choc Spinal Aigu} & \textbf{Moelle Épinière Niveau T10-T11} (Segmentaire T10 + Sus-segmentaire pour membres inférieurs) & \textbf{Traumatisme rachidien} (Fracture-dislocation T10-T11 $\rightarrow$ Compression/Section moelle) \\ \hline
    \textbf{3} & \textbf{Syndrome Névral Périphérique Généralisé (Polyradiculonévrite) Aiguë} & \textbf{Racines rachidiennes (antérieures/sensitives) + Nerfs périphériques + Racines crâniennes} (Atteinte démyélinisante ascendante) & \textbf{Syndrome de Guillain-Barré (SGB) - Variante AIDP (Aiguë Inflammatory Demyelinating Polyradiculoneuropathie)} \\ \hline
    \end{tabularx}
    \end{center}
    \item \textbf{Mécanisme Aréflexie (Cas 2 - Choc Spinal) vs Hyperréflexie (Cas 1 - Pyralidal) :}
    \begin{itemize}
        \item \textbf{Cas 1 (Lésion sus-segmentaire, voie pyramidale coupée en amont du réflexe) :} Perte de l'\textbf{inhibition descendante} permanente exercée par le cortex (voie corticospinale) sur les interneurones inhibiteurs de Renshaw (feedback négatif motoneurone $\rightarrow$ Renshaw $\rightarrow$ inhibition motoneurone) et sur les motoneurones $\alpha$ directement. $\rightarrow$ \textbf{Désinhibition} du motoneurone $\alpha$ $\rightarrow$ \textbf{Hyperexcitabilité} $\rightarrow$ \textbf{Hyperréflexie} (seuil abaissé, réponse amplifiée, clonus).
        \item \textbf{Cas 2 (Lésion segmentaire + sus-segmentaire aiguë - Choc Spinal) :} La lésion \textbf{détruit les corps cellulaires des motoneurones} (cornes antérieures T10) ET \textbf{sectionne les voies descendantes} (pyramidales, réticulospinales, vestibulospinales) qui apportent la \textbf{facilitation tonique} (glutamate, sérotonine, noradrénaline) nécessaire au maintien de l'excitabilité basale des motoneurones lombaires (L3-S1) intacts mais \textbf{désafferentés supraspinalement}. Brutale chute de l'excitabilité $\rightarrow$ \textbf{Aréflexie totale (Choc spinal)}, paralysie flasque. \textit{Note : Si survie, hyperréflexie spastique apparaîtra en 2-6 semaines (supersensibilité déafferentation).}
    \end{itemize}
    \item \textbf{Cas 3 (Guillain-Barré) :}
    \begin{itemize}
        \item \textbf{Dissociation cyto-protéinique :} Protéines \SI{1,2}{g/L} (norme < 0,45) + Cytose < 10/mm$^3$ (norme < 5). Mécanisme : \textbf{Démylinisation inflammatoire aiguë des racines} (infiltration macrophages, lymphocytes T CD4+ anti-gangliosides GM1/GD1a par mimétisme moléculaire post-\textit{Campylobacter jejuni}). Rupture barrière hémato-liquorale au niveau des \textbf{racines} (gainée de Schwann) $\rightarrow$ Passage massif protéines sériques (IgG, albumine) dans LCR. Pas de passage cellules (inflammation radiculaire, pas méningée).
        \item \textbf{Remontée (Céphalisation) :} Atteinte **proximale (racines) $\rightarrow$ distale\textbf{ ? Non, classiquement }distale $\rightarrow$ proximale (pieds $\rightarrow$ tête)**. Mais ici : Racines lombaires/sacrées (longues) atteintes en premier (vulnérabilité zone de transition SNC/SNP, racine sans dure-mère épaisse) $\rightarrow$ Remontée vers racines cervicales, bulbaires (nerfs crâniens IX, X, XII $\rightarrow$ dysphagie, ophtalmoplégie). Atteinte \textbf{longueur-dépendante} des plus longs axones (métabolisme mitochondrial, transport axonal).
    \end{itemize}
    \item \textbf{Abolition Rotulien (L3-L4) avec lésion T10 :}
    \begin{itemize}
        \item Le centre du réflexe rotulien (motoneurones $\alpha$ L3-L4) est \textbf{intact anatomiquement} (situé en dessous de T10).
        \item Mais il est \textbf{fonctionnellement déconnecté} de ses entrées facilitrices descendantes (voies réticulospinales médullaires, vestibulospinales, corticospinales) qui traversent la lésion T10.
        \item En phase aiguë (choc spinal), la perte brutale de cette \textbf{drive facilitrice tonique} (noradrénaline, sérotonine, glutamate) plonge les motoneurones L3-L4 dans un état d'\textbf{inhibition synaptique massive / inexcitabilité}.
        \item L'arc réflexe local (Ia $\rightarrow$ MN $\alpha$) est structuralement intact, mais \textbf{fonctionnellement silencieux} par manque de « biais » excitateur (tonus). C'est une \textbf{aréflexie par déafferentation supraspinale}, pas par destruction de l'arc.
    \end{itemize}
\end{enumerate}
\textbf{Conclusion :} La sémiologie réflexe (hyper/aréflexie) ne permet pas seule de localiser la lésion ; elle doit être corrélée au type de paralysie (spastique/flasque), aux signes sensitifs (niveau, dissociation), aux signes sphinctériens et à l'imagerie. Le \textbf{choc spinal} est un piège diagnostique majeur : il mime une lésion basale (aréflexie) alors que la lésion est haute.
\end{exoresolu}

\begin{attention}[Les 5 erreurs qui coûtent des points au Bac]
\begin{enumerate}
    \item \textbf{Confondre « Racine dorsale » et « Nerf spinal » dans Bell-Magendie :} \textit{Erreur :} « Section du nerf sciatique abolit le réflexe. » \textbf{Correct :} C'est la section de la \textbf{racine dorsale} (ou ganglion) qui abolit la sensibilité et le réflexe ; la section de la racine ventrale abolit la motricité. Le nerf spinal est mixte (sensitif + moteur), sa section donne un syndrome mixte. \textbf{Précision :} La loi s'applique aux \textbf{racines}, pas au nerf mixte.
    \item \textbf{Oublier la « Décussation des Pyramides » pour l'hémiplégie :} \textit{Erreur :} « Lésion hémisphère gauche $\rightarrow$ paralysie gauche. » \textbf{Correct :} Les fibres corticospinales croisent au bulbe (80-85\%). Lésion \textbf{sus-bulbaire} (cortex, capsule interne, pédoncule) $\rightarrow$ Paralysie \textbf{controlatérale}. Lésion \textbf{sous-bulbaire} (moelle) $\rightarrow$ Paralysie \textbf{ipsilatérale}. Toujours préciser le niveau de la lésion par rapport à la décussation.
    \item \textbf{Dire que le réflexe myotatique est « monosynaptique » sans nuance :} \textit{Erreur :} « L'arc réflexe myotatique ne comporte qu'une synapse. » \textbf{Correct :} La voie \textbf{excitatrice pour l'agoniste} est monosynaptique (Ia $\rightarrow$ MN $\alpha$). Mais l'\textbf{inhibition réciproque de l'antagoniste} est \textbf{disynaptique} (Ia $\rightarrow$ Interneurone inhibiteur $\rightarrow$ MN $\alpha$ antagoniste). Le schéma complet DOIT montrer l'interneurone inhibiteur. Au Bac, dessiner l'arc complet avec inhibition réciproque et motoneurone $\gamma$.
    \item \textbf{Confondre « Extinction » et « Oubli » dans le conditionnement :} \textit{Erreur :} « L'extinction efface le réflexe conditionnel / la mémoire est perdue. » \textbf{Correct :} L'extinction est un \textbf{apprentissage inhibiteur nouveau} (SC $\rightarrow$ « Pas de SI ») qui \textbf{masque} la trace excitatory originale (LTP). Preuve : \textbf{Récupération spontanée} (réapparition RC après repos), \textbf{Réacquisition rapide}, \textbf{Effet de renouvellement} (changement de contexte). Ne jamais écrire « oubli » ou « effacement ».
    \item \textbf{Mauvaise numérotation des racines rachidiennes / Hérnies discales :} \textit{Erreur :} « Hernie L4-L5 compresse racine L4. » ou « Réflexe achilléen = L5. » \textbf{Correct :} \textbf{Règle :} La racine sort \textbf{sous} la vertèbre de même nom (sauf C1-C7). Disque L4-L5 $\rightarrow$ Comprime racine \textbf{L5} (qui sort au foramen L5-S1). Disque L5-S1 $\rightarrow$ Comprime racine \textbf{S1}. Réflexe Rotulien = \textbf{L3-L4} (Quadriceps). Réflexe Achilléen = \textbf{S1-S2} (Triceps sural). Bicipital = \textbf{C5-C6}. Stylo-radien = \textbf{C6-C7}. Tricipital = \textbf{C7-C8}. Apprendre par cœur ce tableau de correspondance \textbf{Métamère / Racine / Muscle / Réflexe / Dermatome}.
\end{enumerate}
\end{attention}

\newpage

\coursec{Exercices}

\coursec{Leçon 07 : Limitation des dysfonctionnements des organes ou structures intervenant dans les mouvements réflexes}

\courssub{Je vérifie mes connaissances}

\begin{exercice}[QCM de synthèse]
Pour chaque proposition, choisir la (ou les) bonne(s) réponse(s) en justifiant brièvement.

\textbf{1.} La substance grise de la moelle épinière :
\begin{itemize}
\item[a)] est située à la périphérie de la moelle ;
\item[b)] contient les corps cellulaires des neurones ;
\item[c)] a une forme de H ou de papillon en coupe transversale ;
\item[d)] contient uniquement des fibres myélinisées.
\end{itemize}

\textbf{2.} Le ganglion spinal :
\begin{itemize}
\item[a)] se trouve sur la racine ventrale ;
\item[b)] se trouve sur la racine dorsale ;
\item[c)] contient les corps cellulaires des neurones sensitifs ;
\item[d)] contient les corps cellulaires des motoneurones.
\end{itemize}

\textbf{3.} Le réflexe myotatique :
\begin{itemize}
\item[a)] est un réflexe polysynaptique ;
\item[b)] est un réflexe monosynaptique ;
\item[c)] met en jeu un récepteur appelé fuseau neuromusculaire ;
\item[d)] nécessite obligatoirement la participation de l'encéphale.
\end{itemize}

\textbf{4.} La section de la racine dorsale d'un nerf rachidien entraîne, du côté de la section :
\begin{itemize}
\item[a)] une paralysie motrice du territoire correspondant ;
\item[b)] une perte de la sensibilité du territoire correspondant ;
\item[c)] une paralysie motrice et sensitive ;
\item[d)] aucune conséquence observable.
\end{itemize}
\end{exercice}

\begin{exercice}[Vrai ou faux justifié]
Dire si chacune des affirmations suivantes est vraie ou fausse, puis justifier la réponse en une ou deux phrases.
\begin{enumerate}
\item Chez une grenouille démédullée, l'excitation de la patte par un cristal d'acide ne provoque aucun mouvement de retrait.
\item Le nerf rachidien est un nerf purement moteur.
\item Le réflexe rotulien fait intervenir une seule synapse entre le neurone sensitif et le motoneurone.
\item Un réflexe conditionnel est inné et ne peut jamais disparaître.
\item La substance blanche de la moelle épinière doit sa couleur à la présence de myéline autour des axones.
\end{enumerate}
\end{exercice}

\begin{exercice}[Définitions]
Définir de manière précise et concise les termes suivants :
\begin{enumerate}
\item réflexe ;
\item arc réflexe ;
\item réflexe myotatique ;
\item réflexe conditionnel ;
\item dégénérescence wallérienne.
\end{enumerate}
\end{exercice}

\begin{exercice}[Calculs directs : vitesse de conduction nerveuse]
Lors de l'examen d'un patient à l'hôpital de district de Mbalmayo, le médecin mesure le temps de réponse du réflexe achilléen : le muscle du mollet se contracte \SI{30}{ms} après le choc sur le tendon d'Achille. La longueur totale du circuit nerveux parcouru par l'influx (voie afférente + voie efférente) est estimée à \SI{1,20}{m}.
\begin{enumerate}
\item Calculer la vitesse moyenne de propagation de l'influx nerveux sur ce circuit, en \si{m.s^{-1}}.
\item Sachant que le temps de relais synaptique est d'environ \SI{1}{ms} par synapse et que le réflexe achilléen est monosynaptique, calculer le temps réellement consacré à la conduction le long des fibres, puis la vitesse de conduction corrigée.
\item Comparer cette valeur à la vitesse de conduction d'une fibre myélinisée de gros diamètre (environ \SI{100}{m.s^{-1}}) et commenter.
\end{enumerate}
\end{exercice}

\courssub{Je m'entraîne}

\begin{exercice}[Exploitation des expériences de Bell et Magendie]
Chez un animal de laboratoire, on réalise les opérations suivantes sur les racines rachidiennes d'un nerf spinal. Les résultats sont consignés dans le tableau ci-dessous.

\begin{center}
\begin{tabularx}{\linewidth}{|X|X|}
\hline
\rowcolor{popblueL} \textbf{Opération réalisée} & \textbf{Résultat observé} \\
\hline
Section de la racine dorsale & Perte de la sensibilité du territoire correspondant ; motricité conservée \\
\hline
Stimulation du bout central de la racine dorsale sectionnée & Manifestations de douleur (cris, agitation) \\
\hline
Stimulation du bout périphérique de la racine dorsale sectionnée & Aucune réaction observable \\
\hline
Section de la racine ventrale & Paralysie des muscles du territoire correspondant ; sensibilité conservée \\
\hline
Stimulation du bout central de la racine ventrale sectionnée & Aucune réaction observable \\
\hline
Stimulation du bout périphérique de la racine ventrale sectionnée & Contraction des muscles du territoire correspondant \\
\hline
\end{tabularx}
\end{center}

\begin{enumerate}
\item À partir de ces résultats, déterminer la fonction de la racine dorsale et celle de la racine ventrale. Justifier chaque réponse par les expériences appropriées.
\item Expliquer pourquoi la stimulation du bout central de la racine dorsale provoque une sensation douloureuse alors que la stimulation de son bout périphérique ne provoque rien.
\item Justifier, à partir de ces résultats, que le nerf rachidien est un nerf mixte.
\item Prévoir les conséquences de la section simultanée des deux racines d'un même nerf rachidien.
\end{enumerate}
\end{exercice}

\begin{exercice}[Interprétation des expériences de Pflüger]
On réalise chez des grenouilles les opérations suivantes, et l'on observe la réaction de l'animal lorsqu'on dépose un cristal d'acide sur la peau de la patte postérieure.

\begin{center}
\begin{tabularx}{\linewidth}{|l|X|X|}
\hline
\rowcolor{popblueL} \textbf{Expérience} & \textbf{Opération} & \textbf{Résultat} \\
\hline
E$_1$ & Grenouille intacte & Retrait de la patte excitée \\
\hline
E$_2$ & Grenouille décérébrée (encéphale détruit, moelle intacte) & Retrait de la patte excitée \\
\hline
E$_3$ & Grenouille démédullée (moelle détruite) & Aucun mouvement \\
\hline
E$_4$ & Grenouille décérébrée dont on a sectionné le nerf sciatique de la patte excitée & Aucun mouvement de cette patte \\
\hline
\end{tabularx}
\end{center}

\begin{enumerate}
\item Comparer les résultats des expériences E$_1$ et E$_2$. Quelle structure nerveuse est dispensable pour la réalisation du réflexe de retrait ?
\item Comparer les résultats des expériences E$_2$ et E$_3$. Quel est le rôle de la moelle épinière ?
\item Que montre l'expérience E$_4$ quant au rôle du nerf dans le réflexe ?
\item Dégager les cinq éléments qui interviennent obligatoirement dans cet acte réflexe médullaire.
\item Construire le schéma annoté de l'arc réflexe mis en jeu, en y faisant figurer : récepteur, fibre sensitive, ganglion spinal, centre nerveux, fibre motrice, effecteur.
\end{enumerate}
\end{exercice}

\begin{exercice}[Schéma à compléter : l'arc réflexe myotatique]
Le schéma ci-dessous représente, de façon simplifiée, une coupe transversale de la moelle épinière au niveau lombaire mettant en jeu les structures du réflexe rotulien.

\begin{popfigure}
\begin{tikzpicture}[font=\small, >=Stealth]
% Moelle épinière (coupe transversale)
\draw[thick, fill=popcream] (0,0) ellipse (2.5 and 1.5); % Substance blanche
\draw[thick, fill=poporangeL] (0,0) ellipse (1.8 and 1.0); % Substance grise (forme H approx)
% Cornes
\draw[thick, fill=poporangeL] (-1.8,0.6) .. controls (-1.8,1.2) and (-1.0,1.2) .. (-1.0,0.6) .. controls (-1.0,0.2) and (-1.8,0.2) .. cycle; % Corne dorsale G
\draw[thick, fill=poporangeL] (1.0,0.6) .. controls (1.0,1.2) and (1.8,1.2) .. (1.8,0.6) .. controls (1.8,0.2) and (1.0,0.2) .. cycle; % Corne dorsale D
\draw[thick, fill=poporangeL] (-1.5,-0.6) .. controls (-1.5,-1.2) and (-0.5,-1.2) .. (-0.5,-0.6) .. controls (-0.5,-0.2) and (-1.5,-0.2) .. cycle; % Corne ventrale G
\draw[thick, fill=poporangeL] (0.5,-0.6) .. controls (0.5,-1.2) and (1.5,-1.2) .. (1.5,-0.6) .. controls (1.5,-0.2) and (0.5,-0.2) .. cycle; % Corne ventrale D
% Canal central
\draw[thick, fill=poppinkL] (0,0) circle (0.15);

% Racines et ganglions
% Racine dorsale gauche (entrante)
\draw[thick, popblue] (-2.5,1.0) -- (-1.8,0.6);
\draw[thick, fill=popblueL] (-3.0,1.0) circle (0.25); % Ganglion
\node[above] at (-3.0,1.4) {\textbf{B}};
% Racine ventrale gauche (sortante)
\draw[thick, poppink] (-1.0,-1.0) -- (-0.5,-0.6);
\node[below] at (-1.0,-1.4) {\textbf{C}};

% Neurones et synapse (schématisés dans corne ventrale droite pour clarté)
% Neurone sensitif (fibre centrale) -> entre par racine dorsale D -> va vers corne ventrale D
\draw[thick, popblue, ->] (1.8,0.6) -- (1.2,0.2) -- (1.0,-0.4);
% Synapse
\fill[popdark] (1.0,-0.4) circle (0.06);
\node[right] at (1.0,-0.4) {\textbf{D}};
% Motoneurone (corps cellulaire dans corne ventrale D)
\draw[thick, fill=poppurpleL] (1.5,-0.8) circle (0.15);
\node[right] at (1.8,-0.8) {\textbf{E}};
% Axone motoneurone -> racine ventrale D
\draw[thick, poppink, ->] (1.65,-0.8) -- (2.5,-0.8);
\node[below] at (2.5,-1.2) {\textbf{F}};

% Etiquettes principales
\node at (0,-2.0) {\textbf{A}};
% Légendes des flèches
\node[rotate=45, popblue] at (-2.2,0.8) {Afférent};
\node[rotate=-45, poppink] at (-0.8,-1.1) {Efférent};
\node[rotate=0, popdark] at (1.5,0.2) {Synapse};
\end{tikzpicture}
\legende{Coupe transversale schématique de la moelle épinière (région lombaire) illustrant l'arc réflexe myotatique rotulien. A : Moelle épinière (substance grise et blanche) ; B : Ganglion spinal (sur racine dorsale) ; C : Racine ventrale (sortie motrice) ; D : Synapse monosynaptique ; E : Motoneurone (corps cellulaire dans corne ventrale) ; F : Racine ventrale (fibres motrices). Les flèches bleues indiquent la voie sensitive, roses la voie motrice.}
\end{popfigure}

\begin{enumerate}
\item Identifier les structures légendées par les lettres A, B, C, D, E et F sur le schéma.
\item Indiquer sur le schéma (par une phrase de réponse) le sens de circulation de l'influx nerveux.
\item Rappeler le nombre de synapses présentes dans cet arc et nommer le type de réflexe correspondant.
\item Comparer cet arc à l'arc réflexe médullaire classique (réflexe de retrait polysynaptique) : quelle différence de structure explique que le réflexe rotulien soit plus rapide ?
\item Citer une raison pour laquelle le médecin teste le réflexe rotulien lors d'une consultation neurologique.
\end{enumerate}
\end{exercice}

\begin{exercice}[Analyse de données : l'expérience de Pavlov]
Au cours de ses travaux sur la digestion, Pavlov mesure le volume de salive sécrété par un chien dans différentes situations. Les résultats sont résumés ci-dessous.

\begin{center}
\begin{tabularx}{\linewidth}{|X|c|}
\hline
\rowcolor{popblueL} \textbf{Situation expérimentale} & \textbf{Volume de salive (gouttes)} \\
\hline
Présentation de viande seule (avant conditionnement) & 30 \\
\hline
Sonnerie seule (avant conditionnement) & 0 \\
\hline
Sonnerie associée à la viande, répétée plusieurs fois & 30 \\
\hline
Sonnerie seule (après conditionnement) & 28 \\
\hline
Sonnerie seule répétée sans viande, plusieurs jours de suite & diminue de 28 à 2 \\
\hline
\end{tabularx}
\end{center}

\begin{enumerate}
\item Identifier, dans cette expérience, le stimulus inconditionnel, le stimulus conditionnel et la réponse conditionnelle.
\item Expliquer, en faisant intervenir la notion d'association entre deux excitations au niveau de l'écorce cérébrale, comment le réflexe conditionnel s'établit.
\item Que devient le réflexe conditionnel lorsque le renforcement (la viande) est supprimé de façon répétée ? Comment appelle-t-on ce phénomène ?
\item Donner un exemple de réflexe conditionnel dans la vie quotidienne d'un élève camerounais et préciser son importance adaptative.
\end{enumerate}
\end{exercice}

\courssub{Je me prépare au Bac}

\begin{exercice}[Type Bac : accident de la route et paralysie]
Un moto-taxi de 28 ans est admis en urgence à l'hôpital régional de Bafoussam à la suite d'un accident de la circulation. L'examen clinique révèle :
\begin{itemize}
\item une paralysie motrice complète des deux membres inférieurs ;
\item une perte totale de la sensibilité (tactile et douloureuse) des deux membres inférieurs ;
\item des réflexes rotuliens et achilléens exagérés (hyperréflexie) ;
\item une motricité et une sensibilité normales des membres supérieurs.
\end{itemize}
La radiographie puis l'imagerie médicale montrent une lésion de la moelle épinière au niveau de la région thoracique, la moelle située au-dessous de la lésion étant intacte.

\begin{enumerate}
\item Rappeler la structure générale de la moelle épinière en coupe transversale (substance grise, substance blanche, racines) et préciser le contenu de chacune de ces régions.
\item Expliquer la paralysie motrice et la perte de sensibilité des membres inférieurs en précisant quelles voies nerveuses (ascendantes et descendantes) sont interrompues par la lésion.
\item Les réflexes rotuliens du patient sont conservés et même exagérés.
\begin{enumerate}
\item Rappeler les éléments de l'arc réflexe myotatique mis en jeu dans le réflexe rotulien.
\item Expliquer pourquoi ce réflexe persiste malgré la paralysie.
\item Proposer une explication de l'exagération du réflexe, sachant que l'encéphale exerce normalement un contrôle inhibiteur sur les centres médullaires.
\end{enumerate}
\item Localiser précisément la lésion (région de la moelle) et justifier la réponse à partir du niveau des troubles observés.
\item Un autre patient présente une paralysie des deux jambes avec abolition totale des réflexes rotuliens. En comparant les deux cas, indiquer où se situe vraisemblablement la lésion du second patient (au niveau de l'arc réflexe lui-même ou au-dessus) et justifier.
\end{enumerate}
\end{exercice}

\begin{exercice}[Type Bac : expériences sur la grenouille et arc réflexe]
Dans le laboratoire de SVT d'un lycée de Douala, un groupe d'élèves de Terminale D réalise, sous la conduite de leur professeur, une série d'expériences sur des grenouilles afin de mettre en évidence les éléments d'un acte réflexe médullaire. Le protocole et les observations sont les suivants.

\begin{center}
\begin{tabularx}{\linewidth}{|l|X|X|}
\hline
\rowcolor{popblueL} \textbf{Lot} & \textbf{Traitement} & \textbf{Observation après dépôt d'un cristal d'acide sur la cuisse} \\
\hline
A & Grenouille intacte & La patte se replie et écarte le cristal \\
\hline
B & Décérébration seule & La patte se replie et écarte le cristal \\
\hline
C & Décérébration puis destruction de la moelle épinière & Aucune réaction \\
\hline
D & Décérébration, puis section du nerf sciatique de la patte testée & Aucune réaction de cette patte ; l'autre patte réagit \\
\hline
E & Décérébration, puis ablation de la peau de la cuisse au point d'application & Réaction absente ou très faible au point dénudé \\
\hline
\end{tabularx}
\end{center}

\begin{enumerate}
\item Formuler le problème que ce groupe d'élèves cherche à résoudre, puis proposer une hypothèse de travail.
\item À partir de la comparaison des lots A et B, dégager le rôle (ou l'absence de rôle) de l'encéphale dans ce réflexe.
\item À partir des lots B et C, déterminer le siège du centre nerveux de ce réflexe. Justifier.
\item Montrer, à l'aide des lots D et E, que le nerf et les récepteurs cutanés sont indispensables à la réaction.
\item Énumérer les cinq éléments de l'arc réflexe mis en évidence et réaliser le schéma fonctionnel annoté de cet arc, en y indiquant le sens de propagation de l'influx.
\item Les élèves constatent que le temps de réaction est plus long si l'on applique l'acide très dilué. Proposer une interprétation de cette observation.
\item Citer deux lois de Pflüger concernant les réflexes médullaires et les illustrer à partir des résultats de ce TP.
\end{enumerate}
\end{exercice}

\begin{exercice}[Type Bac : réflexes myotatiques et consultation neurologique]
Lors d'une consultation de neurologie au Centre hospitalier d'Essos à Yaoundé, un médecin teste les réflexes ostéo-tendineux de trois patients à l'aide d'un marteau à réflexes. Il percute le tendon rotulien (sous la rotule, jambe pendante) puis le tendon d'Achille, et consigne ses observations.

\begin{center}
\begin{tabularx}{\linewidth}{|l|X|X|}
\hline
\rowcolor{popblueL} \textbf{Patient} & \textbf{Réflexe rotulien} & \textbf{Autres signes cliniques} \\
\hline
N° 1 & Normal (extension brève de la jambe) & Aucun \\
\hline
N° 2 & Aboli des deux côtés & Atteinte des nerfs des membres inférieurs (neuropathie diabétique) \\
\hline
N° 3 & Exagéré des deux côtés & Paralysie des jambes avec raideur, sensibilité abolie, suite à un traumatisme du dos \\
\hline
\end{tabularx}
\end{center}

On rappelle que le réflexe rotulien met en jeu : un fuseau neuromusculaire (récepteur d'étirement) dans le muscle quadriceps, une fibre sensitive qui entre dans la moelle par la racine dorsale, une synapse unique avec un motoneurone de la corne ventrale, et une fibre motrice qui rejoint le quadriceps par la racine ventrale et le nerf crural.

\begin{enumerate}
\item Définir le réflexe myotatique et préciser ce qui le distingue des autres réflexes médullaires (caractère monosynaptique, absence d'interneurone).
\item Réaliser le schéma légendé de l'arc réflexe myotatique du réflexe rotulien, en y plaçant : le fuseau neuromusculaire, la fibre sensitive, le ganglion spinal, la synapse, le motoneurone, la fibre motrice et le muscle effecteur.
\item Expliquer pourquoi le réflexe rotulien est plus rapide qu'un réflexe de retrait (polysynaptique) : on s'appuiera sur le nombre de synapses et sur le temps de relais synaptique (environ \SI{1}{ms} par synapse).
\item Pour le patient N° 2, indiquer quel(s) élément(s) de l'arc réflexe est (sont) vraisemblablement lésé(s) et justifier l'abolition du réflexe.
\item Pour le patient N° 3, montrer que l'arc réflexe est intact et expliquer l'exagération du réflexe en faisant intervenir la rupture du contrôle descendant de l'encéphale sur la moelle.
\item Déduire de l'ensemble de l'étude l'intérêt clinique du test des réflexes ostéo-tendineux pour le diagnostic et la localisation des lésions nerveuses.
\end{enumerate}
\end{exercice}

\newpage

\coursec{Corrigés des exercices}

\begin{corrige}[Exercice 1 — QCM de synthèse]
\textbf{1.} Bonnes réponses : \textbf{b} et \textbf{c}. La substance grise de la moelle épinière est \cle{centrale} (et non périphérique, contrairement à l'écorce cérébrale) ; elle contient les \cle{corps cellulaires} des neurones (motoneurones dans les cornes ventrales, interneurones et neurones relais dans les cornes dorsales) et forme en coupe transversale un \cle{H} (ou « papillon »). La proposition d) est fausse : les fibres myélinisées constituent la substance blanche, périphérique.

\textbf{2.} Bonnes réponses : \textbf{b} et \textbf{c}. Le \cle{ganglion spinal} est un renflement situé sur la \cle{racine dorsale} ; il contient les corps cellulaires des \cle{neurones sensitifs} (neurones en T). Les corps cellulaires des motoneurones se trouvent dans la corne ventrale de la moelle, jamais dans un ganglion.

\textbf{3.} Bonnes réponses : \textbf{b} et \textbf{c}. Le réflexe myotatique est le seul réflexe \cle{monosynaptique} connu : le neurone sensitif fait directement synapse avec le motoneurone dans la corne ventrale, sans interneurone. Son récepteur est le \cle{fuseau neuromusculaire}, sensible à l'étirement du muscle. Il est purement médullaire : l'encéphale n'est pas nécessaire à sa réalisation (il exerce seulement un contrôle modulateur).

\textbf{4.} Bonne réponse : \textbf{b}. La racine dorsale est \cle{sensitive} : sa section interrompt la voie afférente, d'où une perte de sensibilité du territoire correspondant, sans paralysie motrice (la racine ventrale, motrice, est intacte).
\end{corrige}

\begin{corrige}[Exercice 2 — Vrai ou faux justifié]
\begin{enumerate}
\item \textbf{Vrai.} La moelle épinière est le \cle{centre nerveux} indispensable du réflexe de retrait ; sa destruction (démédullation) rompt l'arc réflexe en son centre : aucune réponse motrice n'est possible malgré l'excitation du récepteur.
\item \textbf{Faux.} Le nerf rachidien est un nerf \cle{mixte} : il résulte de la réunion de la racine dorsale (sensitive) et de la racine ventrale (motrice) ; il contient donc des fibres sensitives et des fibres motrices.
\item \textbf{Vrai.} Le réflexe rotulien est un réflexe \cle{myotatique monosynaptique} : la fibre sensitive issue du fuseau neuromusculaire fait directement synapse avec le motoneurone dans la corne ventrale de la moelle.
\item \textbf{Faux.} Un réflexe conditionnel est \cle{acquis} (et non inné) : il s'établit par apprentissage, par association d'un stimulus indifférent à un stimulus inconditionnel, et il peut \emph{disparaître} par extinction si le renforcement cesse.
\item \textbf{Vrai.} La substance blanche doit sa couleur à la \cle{myéline}, gaine lipidique entourant les axones des fibres nerveuses ascendantes et descendantes qui la constituent.
\end{enumerate}
\end{corrige}

\begin{corrige}[Exercice 3 — Définitions]
\begin{enumerate}
\item \textbf{Réflexe} : réponse motrice (ou sécrétoire) \cle{involontaire}, rapide et stéréotypée, déclenchée par une excitation déterminée d'un récepteur, et réalisée grâce à un circuit nerveux précâblé, l'arc réflexe.
\item \textbf{Arc réflexe} : ensemble des cinq éléments constituant la voie nerveuse d'un réflexe : un \cle{récepteur}, une voie sensitive (afférente), un \cle{centre nerveux} (moelle ou encéphale), une voie motrice (efférente) et un \cle{effecteur} (muscle ou glande).
\item \textbf{Réflexe myotatique} : réflexe médullaire \cle{monosynaptique} déclenché par l'étirement brusque d'un muscle (excitation du fuseau neuromusculaire) et se traduisant par la contraction de ce même muscle ; il participe au maintien de la posture.
\item \textbf{Réflexe conditionnel} : réflexe \cle{acquis} au cours de la vie individuelle, par association répétée d'un stimulus initialement indifférent (stimulus conditionnel) à un stimulus efficace (stimulus inconditionnel) ; il implique l'écorce cérébrale et peut s'éteindre en l'absence de renforcement.
\item \textbf{Dégénérescence wallérienne} : dégénérescence puis disparition de la partie d'une fibre nerveuse située \cle{en aval de la section} (côté opposé au corps cellulaire), car séparée de son centre trophique ; la partie restée reliée au corps cellulaire survit. Elle permet d'identifier le sens de conduction et l'appartenance des fibres d'un nerf.
\end{enumerate}
\end{corrige}

\begin{corrige}[Exercice 4 — Calculs directs : vitesse de conduction nerveuse]
\begin{enumerate}
\item La vitesse moyenne est le rapport de la distance parcourue au temps total :
\[ v = \frac{d}{t} = \frac{\SI{1,20}{m}}{\SI{0,030}{s}} = \SI{40}{m.s^{-1}}. \]
\item Le réflexe achilléen est monosynaptique : il comporte \textbf{une seule synapse}, soit un temps de relais de \SI{1}{ms}. Le temps consacré à la seule conduction le long des fibres est donc :
\[ t_{\text{cond}} = 30 - 1 = \SI{29}{ms} = \SI{0,029}{s}. \]
La vitesse de conduction corrigée vaut :
\[ v_{\text{cor}} = \frac{\SI{1,20}{m}}{\SI{0,029}{s}} \approx \SI{41,4}{m.s^{-1}}. \]
\item La valeur corrigée (\SI{41,4}{m.s^{-1}}) est du même ordre de grandeur que la vitesse des fibres myélinisées de gros diamètre (jusqu'à \SI{100}{m.s^{-1}}), tout en restant inférieure. L'écart s'explique : le circuit comporte des fibres de diamètres variés, le temps mesuré inclut aussi la latence de la contraction musculaire (transmission neuromusculaire et raccourcissement du muscle), et la longueur du circuit n'est qu'estimée. Le résultat confirme que l'influx se propage très rapidement le long des fibres myélinisées et que le relais synaptique ne représente qu'une faible part du temps de réponse.
\end{enumerate}
\end{corrige}

\begin{attention}[Points de vigilance — Exercice 4]
Convertir systématiquement les millisecondes en secondes (\SI{30}{ms} = \SI{0,030}{s}) avant tout calcul ; ne pas oublier de soustraire le temps synaptique dans la question 2 ; exprimer le résultat avec son unité (\si{m.s^{-1}}).
\end{attention}

\begin{corrige}[Exercice 5 — Exploitation des expériences de Bell et Magendie]
\begin{enumerate}
\item \textbf{Fonction de la racine dorsale : elle est sensitive.} Sa section entraîne une perte de sensibilité sans paralysie ; de plus, la stimulation de son bout central (relié aux centres nerveux) provoque une sensation douloureuse. \textbf{Fonction de la racine ventrale : elle est motrice.} Sa section entraîne une paralysie musculaire sans perte de sensibilité ; la stimulation de son bout périphérique (relié aux muscles) provoque la contraction des muscles.
\item La racine dorsale conduit l'influx \cle{de la périphérie vers le centre} (voie afférente). Stimuler le bout central envoie l'influx vers l'encéphale, qui l'interprète comme une douleur ; stimuler le bout périphérique envoie l'influx vers les récepteurs, « à contre-sens », sans effet observable. Cela illustre la \cle{loi de Bell et Magendie} sur le sens de conduction dans les racines.
\item Le nerf rachidien résulte de la réunion d'une racine sensitive (dorsale) et d'une racine motrice (ventrale) : il contient donc les deux types de fibres, c'est un nerf \cle{mixte}, à la fois sensitif et moteur.
\item La section simultanée des deux racines supprime à la fois la voie afférente et la voie efférente : on observe une \cle{paralysie motrice complète} et une \cle{anesthésie totale} (perte de toute sensibilité) du territoire correspondant, avec abolition de tous les réflexes de ce territoire (l'arc réflexe est rompu des deux côtés).
\end{enumerate}
\end{corrige}

\begin{corrige}[Exercice 6 — Interprétation des expériences de Pflüger]
\begin{enumerate}
\item E$_1$ et E$_2$ donnent le même résultat (retrait de la patte) : la destruction de l'encéphale ne supprime pas le réflexe. L'\cle{encéphale est dispensable} pour ce réflexe : celui-ci ne dépend pas du cerveau.
\item E$_2$ (moelle intacte : réflexe présent) et E$_3$ (moelle détruite : réflexe aboli) montrent que la \cle{moelle épinière est le centre nerveux indispensable} du réflexe de retrait : c'est un réflexe \cle{médullaire}.
\item E$_4$ montre que la section du nerf sciatique abolit le réflexe de la patte correspondante : le \cle{nerf est indispensable}, il assure la conduction de l'influx entre la périphérie et la moelle (voies afférente et efférente).
\item Les cinq éléments obligatoires sont : un \cle{récepteur} (récepteurs cutanés de la patte), une \cle{voie sensitive} (fibre afférente, ganglion spinal), un \cle{centre nerveux} (substance grise de la moelle épinière), une \cle{voie motrice} (fibre efférente, motoneurone) et un \cle{effecteur} (muscles de la patte).
\item Schéma annoté de l'arc réflexe :
\begin{popfigure}
\begin{tikzpicture}[font=\small, >=Stealth, node distance=0.4cm]
% Effecteur / récepteur (peau)
\draw[thick, fill=popgoldL] (-5,1.5) rectangle (-3.5,2.5);
\node at (-4.25,2.0) {Récepteur};
\node[below] at (-4.25,1.5) {(peau)};
% Fibre sensitive
\draw[thick, popblue, ->] (-3.5,2.0) -- (-1.5,2.0) node[midway, above]{fibre sensitive};
% Ganglion
\draw[thick, fill=popblueL] (-2.6,2.0) circle (0.22);
\node[below, popblue] at (-2.6,1.75) {ganglion spinal};
% Centre nerveux
\draw[thick, fill=poporangeL] (-1.5,0.5) rectangle (1.5,3.0);
\node at (0,2.75) {\textbf{Centre nerveux}};
\node at (0,2.45) {(moelle épinière)};
% Interneurone
\draw[thick, poppurple] (-1.2,1.9) -- (-0.3,1.9) -- (-0.3,1.2) -- (0.4,1.2);
\fill[popdark] (0.4,1.2) circle (0.05);
\node[below, poppurple] at (-0.3,1.55) {interneurone};
% Motoneurone
\draw[thick, fill=poppurpleL] (0.8,1.2) circle (0.2);
\node[below] at (0.8,0.95) {motoneurone};
% Fibre motrice
\draw[thick, poppink, ->] (1.5,1.2) -- (3.5,1.2) node[midway, above]{fibre motrice};
% Effecteur
\draw[thick, fill=poppinkL] (3.5,0.7) rectangle (5,1.7);
\node at (4.25,1.2) {Effecteur};
\node[below] at (4.25,0.7) {(muscle)};
\end{tikzpicture}
\legende{Schéma fonctionnel de l'arc réflexe médullaire (réflexe de retrait, polysynaptique) : récepteur cutané $\rightarrow$ fibre sensitive (ganglion spinal) $\rightarrow$ centre nerveux (interneurone + motoneurone dans la moelle) $\rightarrow$ fibre motrice $\rightarrow$ muscle effecteur. La flèche bleue indique la voie afférente, la flèche rose la voie efférente.}
\end{popfigure}
\end{enumerate}
\end{corrige}

\begin{corrige}[Exercice 7 — Schéma à compléter : l'arc réflexe myotatique]
\begin{enumerate}
\item \textbf{A} : moelle épinière en coupe transversale (substance blanche périphérique et substance grise centrale en H) ; \textbf{B} : ganglion spinal (renflement de la racine dorsale contenant les corps cellulaires des neurones sensitifs) ; \textbf{C} : racine dorsale (voie sensitive entrante) ; \textbf{D} : synapse unique entre la fibre sensitive et le motoneurone ; \textbf{E} : corps cellulaire du motoneurone (dans la corne ventrale de la substance grise) ; \textbf{F} : racine ventrale emportant les fibres motrices vers le muscle.
\item L'influx circule du \cle{récepteur} (fuseau neuromusculaire du quadriceps) vers la moelle par la fibre sensitive et la \textbf{racine dorsale} (voie afférente), franchit la synapse D dans la corne ventrale, puis ressort par la \textbf{racine ventrale} (voie efférente) vers le muscle quadriceps qui se contracte.
\item Cet arc comporte \cle{une seule synapse} : c'est un arc \cle{monosynaptique}, caractéristique du réflexe \cle{myotatique}.
\item Dans le réflexe de retrait (polysynaptique), un ou plusieurs \cle{interneurones} s'intercalent entre le neurone sensitif et le motoneurone, ce qui ajoute des synapses. Or chaque relais synaptique coûte environ \SI{1}{ms} ; le réflexe rotulien, monosynaptique, est donc plus rapide car il ne comporte qu'un seul relais.
\item Le test du réflexe rotulien renseigne sur l'\cle{intégrité de l'arc réflexe} (récepteur, nerf, moelle lombaire, muscle) et sur l'équilibre du contrôle exercé par l'encéphale : un réflexe aboli évoque une lésion de l'arc lui-même (nerf, moelle), un réflexe exagéré évoque une lésion des voies descendantes (au-dessus des centres médullaires). C'est un examen simple, rapide et non invasif d'orientation diagnostique.
\end{enumerate}
\end{corrige}

\begin{corrige}[Exercice 8 — Analyse de données : l'expérience de Pavlov]
\begin{enumerate}
\item Le \cle{stimulus inconditionnel} est la viande (il déclenche naturellement la salivation : 30 gouttes) ; le \cle{stimulus conditionnel} est la sonnerie (initialement indifférente : 0 goutte) ; la \cle{réponse conditionnelle} est la salivation déclenchée par la sonnerie seule après conditionnement (28 gouttes).
\item Avant conditionnement, la viande excite un centre de l'écorce cérébrale lié à la salivation, tandis que la sonnerie excite un centre auditif sans effet. Lorsque les deux excitations sont \cle{associées de façon répétée} (sonnerie + viande), il s'établit une \cle{connexion temporaire} entre les deux zones corticales : la sonnerie seule suffit alors à déclencher la salivation. Le réflexe conditionnel est ainsi un réflexe \cle{acquis} mettant en jeu l'écorce cérébrale.
\item Si la sonnerie est présentée seule de façon répétée, sans renforcement par la viande, la réponse diminue progressivement (de 28 à 2 gouttes) : le réflexe conditionnel s'affaiblit puis disparaît. C'est le phénomène d'\cle{extinction} du réflexe conditionnel, qui montre qu'il est réversible et nécessite un renforcement périodique.
\item Exemple : un élève qui entend la cloche de midi (ou l'annonce « c'est la récréation ») ressent la faim et sécrète de la salive, ou se précipite vers la cantine ; autre exemple : le bruit des casseroles le soir déclenche l'appétit. Importance adaptative : le réflexe conditionnel permet d'\cle{anticiper} les événements (préparer la digestion, fuir un danger signalé par un bruit) et d'adapter finement le comportement à l'environnement ; il est à la base de l'apprentissage.
\end{enumerate}
\end{corrige}

\begin{corrige}[Exercice 9 — Type Bac : accident de la route et paralysie]
\begin{enumerate}
\item En coupe transversale, la moelle épinière présente : une \cle{substance grise} centrale, en forme de H (ou de papillon), contenant les \cle{corps cellulaires} des neurones (motoneurones dans les cornes ventrales, interneurones dans les cornes dorsales) et le canal épendymaire en son centre ; une \cle{substance blanche} périphérique, formée de \cle{faisceaux de fibres myélinisées} : cordons dorsaux (voies sensitives ascendantes), cordons latéraux et ventraux (voies motrices descendantes et voies sensitives). De chaque côté émergent une racine dorsale sensitive (portant le ganglion spinal) et une racine ventrale motrice, dont la réunion forme le nerf rachidien mixte.
\item La lésion thoracique sectionne les fibres de la substance blanche : les \cle{voies sensitives ascendantes} (cordons dorsaux et spino-thalamiques) ne peuvent plus transmettre à l'encéphale les informations tactiles et douloureuses issues des membres inférieurs, d'où l'anesthésie ; les \cle{voies motrices descendantes} (faisceaux pyramidal et extra-pyramidal) ne peuvent plus transmettre aux motoneurones lombaires les ordres volontaires de l'encéphale, d'où la paralysie motrice. Les membres supérieurs, dont les centres sont situés au-dessus de la lésion (moelle cervicale), restent normaux.
\item \begin{enumerate}
\item L'arc du réflexe rotulien comprend : le \cle{fuseau neuromusculaire} du quadriceps (récepteur d'étirement), la fibre sensitive (corps cellulaire dans le ganglion spinal, entrée par la racine dorsale), une \cle{synapse unique} avec le motoneurone de la corne ventrale (moelle lombaire), la fibre motrice (racine ventrale, nerf crural) et le muscle quadriceps (effecteur).
\item Ce réflexe persiste car son arc est entièrement situé \cle{au-dessous de la lésion} (moelle lombaire intacte) et ne nécessite pas l'encéphale : c'est un réflexe médullaire monosynaptique autonome.
\item Normalement, l'encéphale exerce sur les centres médullaires un contrôle modulateur à dominante \cle{inhibitrice}. La section des voies descendantes supprime ce frein : les motoneurones, « libérés », répondent de façon excessive, d'où l'\cle{hyperréflexie} (réflexes exagérés).
\end{enumerate}
\item La lésion se situe dans la \cle{région thoracique} de la moelle : elle est au-dessus des centres lombaires des membres inférieurs (paralysie + anesthésie des jambes, réflexes conservés) et au-dessous des centres cervicaux des membres supérieurs (restés normaux). C'est une paraplégie par lésion médullaire thoracique.
\item Chez le second patient, les réflexes rotuliens sont abolis : l'\cle{arc réflexe lui-même est lésé} (atteinte de la moelle lombaire, des racines ou du nerf crural), et non les seules voies descendantes. Comparaison : réflexe exagéré = lésion \emph{au-dessus} de l'arc (voies descendantes, arc intact) ; réflexe aboli = lésion \emph{de l'arc} (centre ou fibres périphériques).
\end{enumerate}
\textbf{Barème indicatif (sur 20)} : question 1 : 4 pts ; question 2 : 4 pts ; question 3 : 6 pts (2 + 2 + 2) ; question 4 : 3 pts ; question 5 : 3 pts.
\end{corrige}

\begin{attention}[Points de vigilance — Exercice 9]
Distinguer clairement voies \emph{ascendantes} (sensitives) et \emph{descendantes} (motrices) ; justifier la persistance du réflexe par le caractère \emph{médullaire et monosynaptique} de l'arc situé sous la lésion ; citer explicitement le contrôle inhibiteur de l'encéphale pour expliquer l'hyperréflexie ; ne jamais écrire que la moelle « sous la lésion » est détruite (l'énoncé précise qu'elle est intacte).
\end{attention}

\begin{corrige}[Exercice 10 — Type Bac : expériences sur la grenouille et arc réflexe]
\begin{enumerate}
\item \textbf{Problème} : quelles structures nerveuses (encéphale, moelle, nerfs, récepteurs) sont nécessaires à la réalisation d'un acte réflexe comme le retrait de la patte chez la grenouille ? \textbf{Hypothèse} : le réflexe de retrait est un acte médullaire mettant en jeu une chaîne d'éléments : récepteur cutané, nerf, moelle épinière et muscle, sans intervention obligatoire de l'encéphale.
\item Les lots A (intacte) et B (décérébrée) réagissent de la même façon : la destruction de l'encéphale ne modifie pas la réponse. L'\cle{encéphale n'intervient pas} dans ce réflexe : il n'est ni le centre, ni une voie obligatoire de l'arc.
\item Le lot B (moelle intacte) réagit, le lot C (moelle détruite) ne réagit plus : la \cle{moelle épinière est le centre nerveux} de ce réflexe. C'est donc un réflexe \cle{médullaire} : sans centre nerveux, l'influx sensitif ne peut être relayé vers les muscles.
\item Lot D : la section du nerf sciatique abolit la réponse de la patte opérée (l'autre patte, dont le nerf est intact, réagit) : le \cle{nerf est indispensable} à la conduction de l'influx (voies afférente et efférente). Lot E : l'ablation de la peau au point d'application supprime ou affaiblit fortement la réponse : les \cle{récepteurs cutanés} sont indispensables pour capter l'excitation.
\item Les cinq éléments de l'arc : \cle{récepteur} (peau), \cle{voie sensitive} (fibre afférente, ganglion spinal), \cle{centre nerveux} (moelle épinière), \cle{voie motrice} (fibre efférente, motoneurone), \cle{effecteur} (muscles de la patte). Schéma fonctionnel :
\begin{popfigure}
\begin{tikzpicture}[font=\small, >=Stealth]
\draw[thick, fill=popgoldL] (0,0) rectangle (2,1);
\node at (1,0.5) {Récepteur (peau)};
\draw[thick, popblue, ->] (2,0.5) -- (4.2,0.5) node[midway, above]{voie sensitive};
\draw[thick, fill=popblueL] (3.1,0.5) circle (0.18);
\node[below, popblue] at (3.1,0.3) {ganglion};
\draw[thick, fill=poporangeL] (4.2,-0.6) rectangle (6.6,1.6);
\node at (5.4,1.3) {Centre nerveux};
\node at (5.4,1.0) {(moelle épinière)};
\draw[thick, poppurple] (4.5,0.5) -- (5.1,0.5) -- (5.1,-0.1) -- (5.7,-0.1);
\fill[popdark] (5.7,-0.1) circle (0.05);
\draw[thick, fill=poppurpleL] (6.1,-0.1) circle (0.16);
\draw[thick, poppink, ->] (6.6,-0.1) -- (8.6,-0.1) node[midway, above]{voie motrice};
\draw[thick, fill=poppinkL] (8.6,-0.6) rectangle (10.4,0.4);
\node at (9.5,-0.1) {Effecteur (muscle)};
\end{tikzpicture}
\legende{Schéma fonctionnel annoté de l'arc réflexe médullaire de la grenouille. Sens de propagation de l'influx : récepteur cutané $\rightarrow$ fibre sensitive (ganglion spinal) $\rightarrow$ moelle épinière (interneurone puis motoneurone) $\rightarrow$ fibre motrice $\rightarrow$ muscle.}
\end{popfigure}
\item Avec un acide très dilué, l'excitation est plus faible : il faut plus de temps pour que l'excitation des récepteurs atteigne le \cle{seuil} déclenchant l'influx (sommation temporelle des excitations au niveau du récepteur et des synapses). Le temps de réaction (latence) augmente donc lorsque l'intensité du stimulus diminue.
\item Deux lois de Pflüger : \textbf{loi de l'irradiation} : plus l'excitation est intense, plus la réponse s'étend à d'autres territoires (une excitation forte de la cuisse peut provoquer le retrait des deux pattes, voire de tout le corps) ; \textbf{loi de la symétrie (ou de l'unilatéralité)} : une excitation faible ne provoque la réponse que du côté excité (lot D : seule la patte dont le nerf est intact réagit, du côté de l'excitation). On peut aussi citer la loi de l'intensité : la force de la réponse est proportionnelle à celle de l'excitation.
\end{enumerate}
\textbf{Barème indicatif (sur 20)} : question 1 : 2 pts ; question 2 : 2 pts ; question 3 : 3 pts ; question 4 : 4 pts ; question 5 : 5 pts (dont 3 pour le schéma légendé) ; question 6 : 2 pts ; question 7 : 2 pts.
\end{corrige}

\begin{attention}[Points de vigilance — Exercice 10]
Chaque conclusion doit être appuyée sur la \emph{comparaison de deux lots ne différant que par un seul facteur} (méthode du témoin) ; dans le schéma, flécher le sens de l'influx et légender les cinq éléments ; ne pas attribuer à l'encéphale un rôle dans ce réflexe ; pour les lois de Pflüger, citer la loi \emph{et} l'illustrer par un résultat du TP.
\end{attention}

\begin{corrige}[Exercice 11 — Type Bac : réflexes myotatiques et consultation neurologique]
\begin{enumerate}
\item Le \cle{réflexe myotatique} est un réflexe médullaire déclenché par l'étirement brusque d'un muscle (excitation du fuseau neuromusculaire) et se traduisant par la contraction réflexe de ce muscle. Il se distingue des autres réflexes médullaires (comme le réflexe de retrait) par son caractère \cle{monosynaptique} : la fibre sensitive fait directement synapse avec le motoneurone dans la corne ventrale, \cle{sans interneurone}.
\item Schéma légendé de l'arc myotatique rotulien :
\begin{popfigure}
\begin{tikzpicture}[font=\small, >=Stealth]
% Muscle quadriceps
\draw[thick, fill=poppinkL] (-6,0.5) ellipse (1.2 and 0.7);
\node at (-6,0.5) {Quadriceps};
\node[below] at (-6,-0.3) {(effecteur)};
% Fuseau neuromusculaire
\draw[thick, fill=popgoldL] (-6,1.6) ellipse (0.55 and 0.25);
\node[above] at (-6,1.9) {Fuseau neuromusculaire};
% Fibre sensitive vers moelle
\draw[thick, popblue, ->] (-5.5,1.7) -- (-3,2.6) -- (-1.2,2.6) node[midway, above]{fibre sensitive};
% Ganglion spinal
\draw[thick, fill=popblueL] (-2.2,2.6) circle (0.2);
\node[above, popblue] at (-2.2,2.85) {ganglion spinal};
% Moelle (corne ventrale)
\draw[thick, fill=poporangeL] (-1.2,0.2) rectangle (1.6,3.2);
\node at (0.2,2.95) {Moelle épinière};
\node at (0.2,2.65) {(corne ventrale)};
% Synapse
\draw[thick, popblue] (-1.2,2.6) -- (-0.2,2.6) -- (-0.2,1.6);
\fill[popdark] (-0.2,1.6) circle (0.06);
\node[left] at (-0.3,1.6) {synapse};
% Motoneurone
\draw[thick, fill=poppurpleL] (0.4,1.6) circle (0.18);
\node[below] at (0.4,1.35) {motoneurone};
% Fibre motrice retour muscle
\draw[thick, poppink, ->] (0.58,1.6) -- (1.6,1.6) -- (1.6,0.5) -- (-4.8,0.5) node[pos=0.5, below]{fibre motrice (nerf crural)};
\end{tikzpicture}
\legende{Arc réflexe myotatique du réflexe rotulien : le choc sur le tendon étire le quadriceps ; le fuseau neuromusculaire est excité ; l'influx remonte par la fibre sensitive (ganglion spinal, racine dorsale), franchit une synapse unique avec le motoneurone de la corne ventrale, puis redescend par la fibre motrice (racine ventrale, nerf crural) et provoque la contraction du quadriceps (extension de la jambe).}
\end{popfigure}
\item Le réflexe rotulien ne comporte qu'\cle{une synapse} (environ \SI{1}{ms} de relais), alors qu'un réflexe de retrait polysynaptique en comporte au moins deux (interneurone intercalé), soit au moins \SI{2}{ms} de relais synaptiques, auxquels s'ajoutent souvent des circuits plus longs. Le temps de latence du réflexe myotatique est donc plus court : c'est le réflexe le plus rapide de l'organisme.
\item Chez le patient n° 2 (neuropathie diabétique), le réflexe est aboli alors qu'il n'y a pas de lésion médullaire centrale : ce sont les \cle{nerfs périphériques} qui sont atteints, c'est-à-dire les fibres sensitives et/ou motrices de l'arc (voire les récepteurs). La voie afférente ou efférente étant interrompue, l'arc réflexe est rompu : le message ne peut plus circuler, le réflexe est aboli.
\item Chez le patient n° 3, le réflexe est \emph{exagéré}, donc \cle{présent} : l'arc réflexe (récepteur, fibres, moelle lombaire, muscle) est nécessairement \cle{intact}. Le traumatisme du dos a sectionné les \cle{voies descendantes} de la substance blanche au-dessus des centres lombaires (d'où paralysie et anesthésie des jambes). L'encéphale n'exerce alors plus son contrôle \cle{inhibiteur} sur les motoneurones médullaires : libérés de ce frein, ils répondent de façon excessive, d'où l'hyperréflexie.
\item Le test des réflexes ostéo-tendineux permet de \cle{localiser une lésion nerveuse} : un réflexe \emph{aboli} indique une lésion de l'arc réflexe lui-même (nerf périphérique, racine, moelle au niveau du segment) ; un réflexe \emph{exagéré} indique une lésion des voies descendantes, au-dessus de l'arc (moelle sus-jacente ou encéphale) ; un réflexe \emph{normal} atteste l'intégrité de l'ensemble. C'est donc un examen clinique simple, rapide et précieux pour orienter le diagnostic neurologique.
\end{enumerate}
\textbf{Barème indicatif (sur 20)} : question 1 : 3 pts ; question 2 : 4 pts ; question 3 : 3 pts ; question 4 : 3 pts ; question 5 : 4 pts ; question 6 : 3 pts.
\end{corrige}

\begin{attention}[Points de vigilance — Exercice 11]
Bien nommer le récepteur (\emph{fuseau neuromusculaire}) et insister sur l'\emph{absence d'interneurone} ; dans le schéma, placer le ganglion spinal sur la racine dorsale et le motoneurone dans la corne ventrale ; pour le patient n° 3, articuler explicitement : arc intact (réflexe présent) + rupture du contrôle inhibiteur descendant (réflexe exagéré) ; conclure par l'intérêt \emph{topographique} (localisation de la lésion) du test.
\end{attention}

\newpage

\coursec{Fiche bilan}

\begin{popfigure}
\begin{center}\tcbox[colback=popink,colframe=popink,coltext=white,arc=9pt,boxsep=4pt,left=6pt,right=6pt]{\bfseries\small Réflexes et limitation des dysfonctionnements}\end{center}
\vspace{-2pt}
\begin{tcbraster}[raster columns=3,raster equal height=rows,raster column skip=6pt,raster row skip=6pt,enhanced,blankest]
\begin{tcolorbox}[enhanced,colback=white,colframe=popblue,boxrule=0.9pt,arc=3pt,title={Centres nerveux et neurone},fonttitle=\bfseries\scriptsize,coltitle=white,colbacktitle=popblue,left=4pt,right=4pt,top=2pt,bottom=2pt,boxsep=1pt,toptitle=1.5pt,bottomtitle=1.5pt,halign=left]
\begin{itemize}[nosep,leftmargin=*,itemsep=1pt,font=\scriptsize]
\item Encéphale, moelle épinière
\item Corps cellulaire, dendrites, axone
\end{itemize}
\end{tcolorbox}
\begin{tcolorbox}[enhanced,colback=white,colframe=poporange,boxrule=0.9pt,arc=3pt,title={Racines rachidiennes},fonttitle=\bfseries\scriptsize,coltitle=white,colbacktitle=poporange,left=4pt,right=4pt,top=2pt,bottom=2pt,boxsep=1pt,toptitle=1.5pt,bottomtitle=1.5pt,halign=left]
\begin{itemize}[nosep,leftmargin=*,itemsep=1pt,font=\scriptsize]
\item Bell--Magendie
\item Ventral = moteur ; dorsal = sensitif
\end{itemize}
\end{tcolorbox}
\begin{tcolorbox}[enhanced,colback=white,colframe=popgreen,boxrule=0.9pt,arc=3pt,title={Réflexes médullaires},fonttitle=\bfseries\scriptsize,coltitle=white,colbacktitle=popgreen,left=4pt,right=4pt,top=2pt,bottom=2pt,boxsep=1pt,toptitle=1.5pt,bottomtitle=1.5pt,halign=left]
\begin{itemize}[nosep,leftmargin=*,itemsep=1pt,font=\scriptsize]
\item Lois de Pflüger
\item Arc réflexe : 5 éléments
\end{itemize}
\end{tcolorbox}
\begin{tcolorbox}[enhanced,colback=white,colframe=poppurple,boxrule=0.9pt,arc=3pt,title={Réflexe myotatique},fonttitle=\bfseries\scriptsize,coltitle=white,colbacktitle=poppurple,left=4pt,right=4pt,top=2pt,bottom=2pt,boxsep=1pt,toptitle=1.5pt,bottomtitle=1.5pt,halign=left]
\begin{itemize}[nosep,leftmargin=*,itemsep=1pt,font=\scriptsize]
\item Rotulien, achilléen
\item Fuseau neuromusculaire
\end{itemize}
\end{tcolorbox}
\begin{tcolorbox}[enhanced,colback=white,colframe=poppink,boxrule=0.9pt,arc=3pt,title={Réflexes conditionnels},fonttitle=\bfseries\scriptsize,coltitle=white,colbacktitle=poppink,left=4pt,right=4pt,top=2pt,bottom=2pt,boxsep=1pt,toptitle=1.5pt,bottomtitle=1.5pt,halign=left]
\begin{itemize}[nosep,leftmargin=*,itemsep=1pt,font=\scriptsize]
\item Acquis, éducables
\item Pavlov
\end{itemize}
\end{tcolorbox}
\begin{tcolorbox}[enhanced,colback=white,colframe=popgold,boxrule=0.9pt,arc=3pt,title={Dysfonctionnements},fonttitle=\bfseries\scriptsize,coltitle=white,colbacktitle=popgold,left=4pt,right=4pt,top=2pt,bottom=2pt,boxsep=1pt,toptitle=1.5pt,bottomtitle=1.5pt,halign=left]
\begin{itemize}[nosep,leftmargin=*,itemsep=1pt,font=\scriptsize]
\item Paralysies, aréflexie
\item Diagnostic par les réflexes
\end{itemize}
\end{tcolorbox}
\end{tcbraster}

\legende{Carte mentale de la leçon : autour du thème central (réflexes et dysfonctionnements), six branches résument les savoirs exigibles : structures nerveuses, rôle des racines rachidiennes, réflexes médullaires, réflexe myotatique, réflexes conditionnels et applications cliniques.}
\end{popfigure}

\begin{bilan}
\textbf{Définitions clés}
\begin{itemize}
  \item \cle{Réflexe} : réponse motrice \textbf{innée, involontaire, rapide et stéréotypée} à une stimulation, dont le centre nerveux est la moelle épinière (réflexe médullaire) ou l'encéphale.
  \item \cle{Neurone} : cellule nerveuse comprenant un \textbf{corps cellulaire} (avec noyau), des \textbf{dendrites} (convergence de l'influx) et un \textbf{axone} (divergence), souvent entouré d'une gaine de myéline.
  \item \cle{Arc réflexe} : chaîne nerveuse assurant un réflexe, comportant 5 éléments : \textbf{récepteur $\rightarrow$ voie sensitive (racine dorsale) $\rightarrow$ centre nerveux (substance grise de la moelle) $\rightarrow$ voie motrice (racine ventrale) $\rightarrow$ effecteur} (muscle).
  \item \cle{Réflexe myotatique} : réflexe de \textbf{contraction d'un muscle en réponse à son propre étirement} ; récepteur = fuseau neuromusculaire ; arc \textbf{monosynaptique} (ex. : réflexes rotulien et achilléen).
  \item \cle{Réflexe conditionnel (acquis)} : réflexe établi par \textbf{apprentissage} (expériences de Pavlov), à centre cérébral ; il est réversible et éducable.
\end{itemize}

\textbf{Résultats expérimentaux à connaître par c\oe ur}
\begin{center}
\begin{tabularx}{\linewidth}{|l|X|X|}
\hline
\rowcolor{popblueL} \textbf{Expérience} & \textbf{Protocole} & \textbf{Conclusion} \\
\hline
Bell et Magendie & Section de la racine ventrale : membre paralysé mais sensible ; section de la racine dorsale : membre mobile mais insensible & Racine \textbf{ventrale = motrice} ; racine \textbf{dorsale = sensitive} \\
\hline
Dégénérescence wallérienne & Après section d'un nerf, le segment \textbf{séparé du corps cellulaire} dégénère & Le corps cellulaire est le \textbf{centre trophique} du neurone \\
\hline
Pflüger (grenouille démédullée / décérébrée) & Grenouille sans encéphale : réflexes conservés ; après destruction de la moelle : réflexes abolis & Le centre des réflexes étudiés est la \textbf{moelle épinière} \\
\hline
Lois de Pflüger & Stimulation acide d'une patte de grenouille spinale & Loi d'\textbf{unilatéralité}, de \textbf{symétrie}, d'\textbf{irradiation} et de \textbf{généralisation} selon l'intensité du stimulus \\
\hline
\end{tabularx}
\end{center}

\textbf{Méthodes express}
\begin{itemize}
  \item \textbf{Schéma d'un arc réflexe} : toujours dessiner la coupe de moelle (substance grise en H au centre, substance blanche en périphérie), le ganglion spinal sur la racine dorsale, et flécher le circuit dans le sens récepteur $\rightarrow$ effecteur.
  \item \textbf{Diagnostiquer une lésion} : réflexe aboli + sensibilité conservée $\Rightarrow$ atteinte motrice (racine ventrale ou muscle) ; réflexe aboli + sensibilité abolie $\Rightarrow$ atteinte sensitive ou du centre ; réflexe exagéré $\Rightarrow$ lésion des voies centrales inhibitrices.
  \item \textbf{Interpréter une expérience} : annoncer la variable modifiée, comparer témoin/expérience, conclure sur le rôle de la structure étudiée.
\end{itemize}
\end{bilan}

\begin{aretenir}[Checklist avant le Bac]
\begin{itemize}
  \item $\square$ Je sais décrire les grandes lignes de la structure de l'encéphale et de la moelle épinière (substance grise centrale, substance blanche périphérique).
  \item $\square$ Je sais dessiner et légender un neurone (corps cellulaire, dendrites, axone, myéline) et une coupe transversale de nerf.
  \item $\square$ Je sais interpréter les expériences de Bell et Magendie : racine ventrale motrice, racine dorsale sensitive.
  \item $\square$ Je sais interpréter la dégénérescence wallérienne : le corps cellulaire est le centre trophique du neurone.
  \item $\square$ Je connais les expériences de Pflüger sur la grenouille et j'en déduis le rôle de centre réflexe de la moelle épinière.
  \item $\square$ Je connais les quatre lois de Pflüger (unilatéralité, symétrie, irradiation, généralisation).
  \item $\square$ Je sais citer et ordonner les 5 éléments de l'arc réflexe et en réaliser le schéma légendé.
  \item $\square$ Je sais décrire le réflexe myotatique (rotulien, achilléen), son récepteur (fuseau neuromusculaire) et son importance (tonus, posture, diagnostic médical).
  \item $\square$ Je sais distinguer réflexe inné (médullaire) et réflexe conditionnel (acquis, Pavlov).
  \item $\square$ Je sais expliquer une paralysie ou une absence de réflexe par la localisation de la lésion sur l'arc réflexe.
  \item $\square$ Je flèche toujours mes schémas dans le sens de propagation de l'influx nerveux.
  \item $\square$ Je rédige mes interprétations en trois temps : observation, comparaison, conclusion.
\end{itemize}
\end{aretenir}

\findecours

\end{document}
